%%%%%%%%%%%%%%%%%%%%%%%%%%%%%%%%%%%%%%%%%%%%%%%%%%%%%%%%%%%%%%%%%
\chapter{Unit 11: Systems of Linear Equation}
\label{chap:unit11}
%%%%%%%%%%%%%%%%%%%%%%%%%%%%%%%%%%%%%%%%%%%%%%%%%%%%%%%%%%%%%%%%%

%%===============================================================
\section{Systems of Linear Equation}
\label{sec:systems-of-linear-equation}
%%===============================================================

\paragraph{Subspaces}



For any vector space, a \textbf{subspace} is a subset that is itself a vector

space, under the inherited operations.



\paragraph{Lemma 1}



For a nonempty subset $S$ of a vector space, under the inherited

operations, the following are equivalent statements.



1.  $S$ is a subspace of that vector space

2.  $S$ is closed under linear combinations of pairs of vectors: for any

    vectors $\vec{s}_1,\vec{s}_2\in S$ and scalars $r_1,r_2$ the vector

    $r_1\vec{s}_1+r_2\vec{s}_2$ is in $S$

3.  $S$ is closed under linear combinations of any number of vectors:

    for any vectors $\vec{s}_1,\ldots,\vec{s}_n\in S$ and scalars

    $r_1, \ldots,r_n$ the vector $r_1\vec{s}_1+\cdots+r_n\vec{s}_n$ is

    in $S$.



Briefly, the way that a subset gets to be a subspace is by being closed

under linear combinations.



>   Proof:



    >>   \"The following are equivalent\" means that each pair of

        statements are equivalent.



$$(1)\!\iff\!(2)

\qquad

(2)\!\iff\!(3)

\qquad

(3)\!\iff\!(1)$$



>   



    >   We will show this equivalence by establishing that

        $(1)\implies (3)\implies (2)\implies (1)$. This strategy is

        suggested by noticing that $(1)\implies (3)$ and

        $(3)\implies (2)$ are easy and so we need only argue the single

        implication $(2)\implies (1)$.



<!-- -->



>   



    >   For that argument, assume that $S$ is a nonempty subset of a

        vector space $V$ and that $S$ is closed under combinations of

        pairs of vectors. We will show that $S$ is a vector space by

        checking the conditions.



<!-- -->



>   



    >   The first item in the vector space definition has five

        conditions. First, for closure under addition, if

        $\vec{s}_1,\vec{s}_2\in S$ then $\vec{s}_1+\vec{s}_2\in S$, as

        $\vec{s}_1+\vec{s}_2=1\cdot\vec{s}_1+1\cdot\vec{s}_2$.

    >   Second, for any $\vec{s}_1,\vec{s}_2\in S$, because addition is

        inherited from $V$, the sum $\vec{s}_1+\vec{s}_2$ in $S$ equals

        the sum $\vec{s}_1+\vec{s}_2$ in $V$, and that equals the sum

        $\vec{s}_2+\vec{s}_1$ in $V$ (because $V$ is a vector space, its

        addition is commutative), and that in turn equals the sum

        $\vec{s}_2+\vec{s}_1$ in $S$. The argument for the third

        condition is similar to that for the second.

    >   For the fourth, consider the zero vector of $V$ and note that

        closure of $S$ under linear combinations of pairs of vectors

        gives that (where $\vec{s}$ is any member of the nonempty set

        $S$) $0\cdot\vec{s}+0\cdot\vec{s}=\vec{0}$ is in $S$; showing

        that $\vec{0}$ acts under the inherited operations as the

        additive identity of $S$ is easy.

    >   The fifth condition is satisfied because for any $\vec{s}\in S$,

        closure under linear combinations shows that the vector

        $0\cdot\vec{0}+(-1)\cdot\vec{s}$ is in $S$; showing that it is

        the additive inverse of $\vec{s}$ under the inherited operations

        is routine.



We usually show that a subset is a subspace with $(2)\implies (1)$.



\paragraph{Example 1}



>   The plane

    $P=\{\begin{pmatrix} x \\\     y \\\     z \end{pmatrix}\,\big|\, x+y+z=0\}$

    is a subspace of $\mathbb{R}^3$. As specified in the definition, the

    operations are the ones inherited from the larger space, that is,

    vectors add in $P$ as they add in $\mathbb{R}^3$



<!-- -->



>   



    >   $\begin{pmatrix} x_1 \\\     y_1 \\\     z_1 \end{pmatrix}+\begin{pmatrix} x_2 \\\     y_2 \\\     z_2 \end{pmatrix}

        =\begin{pmatrix} x_1+x_2 \\\         y_1+y_2 \\\         z_1+z_2 \end{pmatrix}$



<!-- -->



>   and scalar multiplication is also the same as it is in

    $\mathbb{R}^3$. To show that $P$ is a subspace, we need only note

    that it is a subset and then verify that it is a space. Checking

    that $P$ satisfies the conditions in the definition of a vector

    space is routine. For instance, for closure under addition, just

    note that if the summands satisfy that $x_1+y_1+z_1=0$ and

    $x_2+y_2+z_2=0$ then the sum satisfies that

    $(x_1+x_2)+(y_1+y_2)+(z_1+z_2)=(x_1+y_1+z_1)+(x_2+y_2+z_2)=0$.



\paragraph{Example 2}



>   The $x$-axis in $\mathbb{R}^2$ is a subspace where the addition and

    scalar multiplication operations are the inherited ones.



    >   $\begin{pmatrix} x_1 \\\     0 \end{pmatrix}

        +

        \begin{pmatrix} x_2 \\\         0 \end{pmatrix}

        =

        \begin{pmatrix} x_1+x_2 \\\         0 \end{pmatrix}

        \qquad

        r\cdot\begin{pmatrix} x \\\         0 \end{pmatrix}

        =\begin{pmatrix} rx \\\         0 \end{pmatrix}$



<!-- -->



>   As above, to verify that this is a subspace, we simply note that it

    is a subset and then check that it satisfies the conditions in

    definition of a vector space. For instance, the two closure

    conditions are satisfied: (1) adding two vectors with a second

    component of zero results in a vector with a second component of

    zero, and (2) multiplying a scalar times a vector with a second

    component of zero results in a vector with a second component of

    zero.



\paragraph{Example 3}



>   Another subspace of $\mathbb{R}^2$ is



<!-- -->



>   



    >   $\{\begin{pmatrix} 0 \\\     0 \end{pmatrix}\}$



<!-- -->



>   which is its \textbf{trivial subspace}.



<!-- -->



>   Any vector space has a trivial subspace $\{\vec{0}\,\}$.



At the opposite extreme, any vector space has itself for a subspace.

`{{anchor|improper}}`{=mediawiki}These two are the \textbf{improper}

subspaces. `{{anchor|proper}}`{=mediawiki}Other subspaces are

\textbf{proper}.



\paragraph{Example 4}



The condition in the definition requiring that the addition and scalar

multiplication operations must be the ones inherited from the larger

space is important. Consider the subset $\{1\}$ of the vector space

$\mathbb{R}^1$. Under the operations $1+1=1$ and $r\cdot 1=1$ that set

is a vector space, specifically, a trivial space. But it is not a

subspace of $\mathbb{R}^1$ because those aren\''t the inherited

operations, since of course $\mathbb{R}^1$ has $1+1=2$.



\paragraph{Example 5}



>   All kinds of vector spaces, not just $\mathbb{R}^n$\''s, have

    subspaces. The vector space of cubic polynomials

    $\{a+bx+cx^2+dx^3\,\big|\, a,b,c,d\in\mathbb{R}\}$ has a subspace

    comprised of all linear polynomials

    $\{m+nx\,\big|\, m,n\in\mathbb{R}\}$.



\paragraph{Example 6}



>   This is a subspace of the $2 \! \times \! 2$ matrices



$$L=\{\begin{pmatrix}

a  &0  \\b  &c

\end{pmatrix}

\,\big|\, a+b+c=0\}$$



>   (checking that it is nonempty and closed under linear combinations

    is easy).

>   To parametrize, express the condition as $a=-b-c$.



$$L

=\{\begin{pmatrix}

-b-c  &0  \\b     &c

\end{pmatrix}

\,\big|\, b,c\in\mathbb{R}\}

=\{b\begin{pmatrix}

-1    &0  \\1     &0

\end{pmatrix}

+c\begin{pmatrix}

-1    &0  \\0     &1

\end{pmatrix}

\,\big|\, b,c\in\mathbb{R}\}$$



>   As above, we\''ve described the subspace as a collection of

    unrestricted linear combinations (by coincidence, also of two

    elements).



\paragraph{Span}



The \textbf{span}(or \textbf{linear closure}) of a nonempty subset $S$ of a vector

space is the set of all linear combinations of vectors from $S$.



$$[S] =\{ c_1\vec{s}_1+\cdots+c_n\vec{s}_n

\,\big|\, c_1,\ldots, c_n\in\mathbb{R}

\text{ and } \vec{s}_1,\ldots,\vec{s}_n\in S \}$$ The span of the empty

subset of a vector space is the trivial subspace. No notation for the

span is completely standard. The square brackets used here are common,

but so are \"$\mbox{span}(S)$\" and \"$\mbox{sp}(S)$\".



\paragraph{Lemma 2}



In a vector space, the span of any subset is a subspace.



Proof:



>   Call the subset $S$. If $S$ is empty then by definition its span is

    the trivial subspace. If $S$ is not empty, then we need only check

    that the span $[S]$ is closed under linear combinations. For a pair

    of vectors from that span,

    $\vec{v}=c_1\vec{s}_1+\cdots+c_n\vec{s}_n$ and

    $\vec{w}=c_{n+1}\vec{s}_{n+1}+\cdots+c_m\vec{s}_m$, a linear

    combination



    >   $p\cdot(c_1\vec{s}_1+\cdots+c_n\vec{s}_n)+

        r\cdot(c_{n+1}\vec{s}_{n+1}+\cdots+c_m\vec{s}_m)$

    >   $=

        pc_1\vec{s}_1+\cdots+pc_n\vec{s}_n

        +rc_{n+1}\vec{s}_{n+1}+\cdots+rc_m\vec{s}_m$

>   ($p$, $r$ scalars) is a linear combination of elements of $S$ and so

    is in $[S]$ (possibly some of the $\vec{s}_i$\''s forming $\vec{v}$

    equal some of the $\vec{s}_j$\''s from $\vec{w}$, but it does not

    matter).



\paragraph{Example 7}



The span of this set is all of $\mathbb{R}^2$.



$$\{\begin{pmatrix} 1 \\\ 1 \end{pmatrix},\begin{pmatrix} 1 \\\ -1 \end{pmatrix}\}$$



To check this we must show that any member of $\mathbb{R}^2$ is a linear

combination of these two vectors. So we ask: for which vectors (with

real components $x$ and $y$) are there scalars $c_1$ and $c_2$ such that

this holds?



$$c_1\begin{pmatrix} 1 \\\ 1 \end{pmatrix}+c_2\begin{pmatrix} 1 \\\ -1 \end{pmatrix}=\begin{pmatrix} x \\\ y \end{pmatrix}$$



Gauss\'' method



$$\begin{array}{rcl}

\begin{array}{*{2}{rc}r}

c_1  &+  &c_2  &=  &x  \\c_1  &-  &c_2  &=  &y

\end{array}

&\xrightarrow[]{-\rho_1+\rho_2}

&\begin{array}{*{2}{rc}r}

c_1  &+  &c_2    &=  &x  \\&   &-2c_2  &=  &-x+y

\end{array}

\end{array}$$ with back substitution gives $c_2=(x-y)/2$ and

$c_1=(x+y)/2$. These two equations show that for any $x$ and $y$ that we

start with, there are appropriate coefficients $c_1$ and $c_2$ making

the above vector equation true. For instance, for $x=1$ and $y=2$ the

coefficients $c_2=-1/2$ and $c_1=3/2$ will do. That is, any vector in

$\mathbb{R}^2$ can be written as a linear combination of the two given

vectors.



\paragraph{Linear Independence}



We first characterize when a vector can be removed from a set without

changing the span of that set.



\paragraph{Lemma 3}



Where $S$ is a subset of a vector space $V$,



$$[S]=[S\cup\{\vec{v}\}]

\quad\text{if and only if}\quad

\vec{v}\in[S]$$ for any $\vec{v}\in V$.



>   Proof: The left to right implication is easy. If

    $[S]=[S\cup\{\vec{v}\}]$ then, since $\vec{v}\in[S\cup\{\vec{v}\}]$,

    the equality of the two sets gives that $\vec{v}\in[S]$.

>   For the right to left implication assume that $\vec{v}\in [S]$ to

    show that $[S]=[S\cup\{\vec{v}\}]$ by mutual inclusion. The

    inclusion \`{=html} is obvious. For the other inclusion

    $[S]\supseteq[S\cup\{\vec{v}\}]$, write an element of

    $[S\cup\{\vec{v}\}]$ as $d_0\vec{v}+d_1\vec{s}_1+\dots+d_m\vec{s}_m$

    and substitute $\vec{v}$\''s expansion as a linear combination of

    members of the same set

    $d_0(c_0\vec{t}_0+\dots+c_k\vec{t}_k)+d_1\vec{s}_1+\dots+d_m\vec{s}_m$.

    This is a linear combination of linear combinations and so

    distributing $d_0$ results in a linear combination of vectors from

    $S$. Hence each member of $[S\cup\{\vec{v}\}]$ is also a member of

    $[S]$.



<!-- -->



>   Example: In $\mathbb{R}^3$, where



    >   $\vec{v}_1=\begin{pmatrix} 1 \\\     0 \\\     0 \end{pmatrix}\quad

        \vec{v}_2=\begin{pmatrix} 0 \\\         1 \\\         0 \end{pmatrix}\quad

        \vec{v}_3=\begin{pmatrix} 2 \\\         1 \\\         0 \end{pmatrix}$



<!-- -->



>   the spans $[\{\vec{v}_1,\vec{v}_2\}]$ and

    $[\{\vec{v}_1,\vec{v}_2,\vec{v}_3\}]$ are equal since $\vec{v}_3$ is

    in the span $[\{\vec{v}_1,\vec{v}_2\}]$.



Lemma 2 says that if we have a spanning set then we can remove a

$\vec{v}$ to get a new set $S$ with the same span if and only if

$\vec{v}$ is a linear combination of vectors from $S$. Thus, under the

second sense described above, a spanning set is minimal if and only if

it contains no vectors that are linear combinations of the others in

that set. We have a term for this important property.



\paragraph{Definition of Linear Independence}



A subset of a vector space is \textbf{linearly independent} if none of its

elements is a linear combination of the others. Otherwise it is

\textbf{linearly dependent}. }}



Here is an important observation:



$$\vec{s}_0=c_1\vec{s}_1+c_2\vec{s}_2+\cdots +c_n\vec{s}_n$$



although this way of writing one vector as a combination of the others

visually sets $\vec{s}_0$ off from the other vectors, algebraically

there is nothing special in that equation about $\vec{s}_0$. For any

$\vec{s}_i$ with a coefficient $c_i$ that is nonzero, we can rewrite the

relationship to set off $\vec{s}_i$.



$$\vec{s}_i=(1/c_i)\vec{s}_0+(-c_1/c_i)\vec{s}_1

+\dots+(-c_n/c_i)\vec{s}_n$$



When we don\''t want to single out any vector by writing it alone on one

side of the equation, we will instead say that

$\vec{s}_0,\vec{s}_1,\dots,\vec{s}_n$ are in a \textbf{linear relationship}

and write the relationship with all of the vectors on the same side. The

next result rephrases the linear independence definition in this style.

It gives what is usually the easiest way to compute whether a finite set

is dependent or independent.



\paragraph{Lemma 4}



A subset $S$ of a vector space is linearly independent if and only if

for any distinct $\vec{s}_1,\dots,\vec{s}_n\in S$ the only linear

relationship among those vectors



$$c_1\vec{s}_1+\dots+c_n\vec{s}_n=\vec{0}

\qquad c_1,\dots,c_n\in\mathbb{R}$$



is the trivial one: $c_1=0,\dots,\,c_n=0$. Proof: This is a direct

consequence of the observation above.



>   If the set $S$ is linearly independent then no vector $\vec{s}_i$

    can be written as a linear combination of the other vectors from $S$

    so there is no linear relationship where some of the $\vec{s}\,$\''s

    have nonzero coefficients. If $S$ is not linearly independent then

    some $\vec{s}_i$ is a linear combination

    $\vec{s}_i=c_1\vec{s}_1+\dots+c_{i-1}\vec{s}_{i-1} +c_{i+1}\vec{s}_{i+1}+\dots+c_n\vec{s}_n$

    of other vectors from $S$, and subtracting $\vec{s}_i$ from both

    sides of that equation gives a linear relationship involving a

    nonzero coefficient, namely the $-1$ in front of $\vec{s}_i$.



\paragraph{Example 8}



In the vector space of two-wide row vectors, the two-element set

$\{ \begin{pmatrix} 40 &15 \end{pmatrix},\begin{pmatrix} -50 &25 \end{pmatrix}\}$

is linearly independent. To check this, set



$$c_1\cdot\begin{pmatrix} 40 &15 \end{pmatrix}+c_2\cdot\begin{pmatrix} -50 &25 \end{pmatrix}=\begin{pmatrix} 0 &0 \end{pmatrix}$$



and solving the resulting system



$$\begin{array}{*{2}{rc}r}

40c_1 &- &50c_2 &= &0 \\15c_1 &+ &25c_2 &= &0

\end{array}

\;\xrightarrow[]{-(15/40)\rho_1+\rho_2}\;

\begin{array}{*{2}{rc}r}

40c_1 &- &50c_2    &= &0 \\& &(175/4)c_2 &= &0

\end{array}$$



shows that both $c_1$ and $c_2$ are zero. So the only linear

relationship between the two given row vectors is the trivial

relationship.



In the same vector space,

$\{ \begin{pmatrix} 40 &15 \end{pmatrix},\begin{pmatrix} 20 &7.5 \end{pmatrix}\}$

is linearly dependent since we can satisfy



$$c_1\begin{pmatrix} 40 &15 \end{pmatrix}+c_2\cdot\begin{pmatrix} 20 &7.5 \end{pmatrix}=\begin{pmatrix} 0 &0 \end{pmatrix}$$



with $c_1=1$ and $c_2=-2$.



\paragraph{Example 9}



The set $\{1+x,1-x\}$ is linearly independent in $\mathcal{P}_2$, the

space of quadratic polynomials with real coefficients, because



$$0+0x+0x^2

=

c_1(1+x)+c_2(1-x)

=

(c_1+c_2)+(c_1-c_2)x+0x^2$$



gives



$$\begin{array}{rcl}

\begin{array}{*{2}{rc}r}

c_1 &+ &c_2 &= &0 \\c_1 &- &c_2 &= &0

\end{array}

&\xrightarrow[]{-\rho_1+\rho_2}

&\begin{array}{*{2}{rc}r}

c_1 &+ &c_2 &= &0 \\&  &2c_2 &= &0

\end{array}

\end{array}$$ since polynomials are equal only if their coefficients are

equal. Thus, the only linear relationship between these two members of

$\mathcal{P}_2$ is the trivial one.



\paragraph{Example 10}



In $\mathbb{R}^3$, where



$$\vec{v}_1=\begin{pmatrix} 3 \\\ 4 \\\ 5 \end{pmatrix}

\quad

\vec{v}_2=\begin{pmatrix} 2 \\\ 9 \\\ 2 \end{pmatrix}

\quad

\vec{v}_3=\begin{pmatrix} 4 \\\ 18 \\\ 4 \end{pmatrix}$$



the set $S=\{\vec{v}_1,\vec{v}_2,\vec{v}_3\}$ is linearly dependent

because this is a relationship



$$0\cdot\vec{v}_1

+2\cdot\vec{v}_2

-1\cdot\vec{v}_3

=\vec{0}$$



where not all of the scalars are zero (the fact that some of the scalars

are zero doesn\''t matter).



\paragraph{Resources}



\begin{itemize}
    \item \href{https://textbooks.math.gatech.edu/ila/subspaces.html}{Subspaces},
\end{itemize}
    Interactive Linear Algebra from Georgia Tech



\paragraph{Licensing}



Content obtained and/or adapted from:



\begin{itemize}
    \item \href{https://en.wikibooks.org/wiki/Linear\_Algebra/Definition\_and\_Examples\_of\_Linear\_Independence}{Definition and Examples of Linear Independence, WikiBooks Linear
\end{itemize}
    Algebra}

    under a CC BY-SA license

\begin{itemize}
    \item \href{https://en.wikibooks.org/wiki/Linear\_Algebra/Subspaces\_and\_Spanning\_sets}{Subspaces and Spanning Sets, WikiBooks Linear
\end{itemize}
    Algebra}

    under a CC BY-SA license

\begin{itemize}
    \item \href{https://en.wikibooks.org/wiki/Linear\_Algebra/Subspaces}{Subspaces, WikiBooks Linear
\end{itemize}
    Algebra}

    under a CC BY-SA license

%%===============================================================
\section{Vectors and Matrices}
\label{sec:vectors-and-matrices}
%%===============================================================

\paragraph{Matrices}

A matrix is an array of numbers arranged into rows and columns. Some
examples of matrices are,
$$A=
\begin{bmatrix}
2 & 4 & 6 \\\\ 0 & -5 & 7.105\\\\ 1 & -3 & 2
\end{bmatrix}
,\quad
B=
\begin{bmatrix}
    -3 & -4 \\\\ \pi & \sqrt{2}
\end{bmatrix}
,\textrm{ and }\,\,\,
C=
\begin{bmatrix}
    -3 & -5 & 0 & 0 \\\\ -1 & 4.56 & 3.28 & 19
\end{bmatrix}
.$$

When describing matrices we indicate the number of rows first, then the
number of columns. For example, the matrix $C$ with two rows and four
columns is said to be a $2\times 4$ matrix.

It is standard notation to name matrices with capital letters and to use
lower case letters with subscripts to identify particular entries in a
matrix.

For example, to identify the entry in row 1 and column 3 of matrix $A$
we would write $a_{13}$. To indicate that this entry is a six we would
write the equation $a_{13}=6$.

Two matrices are considered to be \textbf{equal} only if they are the same
size and every pair of corresponding elements are equal.

A \textbf{column matrix} is a matrix with only one column. Similarly, a **row
matrix** has only one row.

\paragraph{Vectors}

A \textbf{vector} is an object often defined by a long list of properties.
However, for now we will avoid the more complicated definition, and just
say that a vector is an ordered list of numbers. Later we will see that
vectors can really be much more.

An ordered pair, $(x, y)$, that is used to identify a point in the plane
can be considered to be a vector.

Similarly, an ordered triple, $(x, y, z)$ is a vector.

Obviously, row and column matrices can also be considered to be a
vector.

It is common to name vectors using variables with arrows above.

For example, we might write $\vec{v}=(2, 3, 5, -4),\textrm{ or } \vec{w}=
\begin{bmatrix} 4 \\\\ 5 \\\\ 0 
\end{bmatrix}$.

For the most part, it will convenient to think of vectors as column
matrices.

"In mathematics, \textbf{Gaussian elimination}, also known as **row
reduction**, is an algorithm for solving systems of linear equations. It
consists of a sequence of operations performed on the corresponding
matrix of coefficients. This method can also be used to compute the rank
of a matrix, the determinant of a square matrix, and the inverse of an
invertible matrix. The method is named after Carl Friedrich Gauss
(1777–1855) although some special cases of the method—albeit presented
without proof—were known to Chinese mathematicians as early as circa
179 CE.

To perform row reduction on a matrix, one uses a sequence of elementary
row operations to modify the matrix until the lower left-hand corner of
the matrix is filled with zeros, as much as possible. There are three
types of elementary row operations:

\begin{itemize}
    \item Swapping two rows,
    \item Multiplying a row by a nonzero number,
    \item Adding a multiple of one row to another row. (subtraction can be
\end{itemize}
    achieved by multiplying one row with -1 and adding the result to
    another row)

Using these operations, a matrix can always be transformed into an upper
triangular matrix, and in fact one that is in row echelon form. Once all
of the leading coefficients (the leftmost nonzero entry in each row) are
1, and every column containing a leading coefficient has zeros
elsewhere, the matrix is said to be in reduced row echelon form. This
final form is unique; in other words, it is independent of the sequence
of row operations used. For example, in the following sequence of row
operations (where two elementary operations on different rows are done
at the first and third steps), the third and fourth matrices are the
ones in row echelon form, and the final matrix is the unique reduced row
echelon form.

$$\begin{bmatrix}
1 & 3 & 1 & 9 \\\\
1 & 1 & -1 & 1 \\\\
3 & 11 & 5 & 35
\end{bmatrix}\to
\begin{bmatrix}
1 & 3 & 1 & 9 \\\\
0 & -2 & -2 & -8 \\\\
0 & 2 & 2 & 8
\end{bmatrix}\to
\begin{bmatrix}
1 & 3 & 1 & 9 \\\\
0 & -2 & -2 & -8 \\\\
0 & 0 & 0 & 0
\end{bmatrix}\to
\begin{bmatrix}
1 & 0 & -2 & -3 \\\\
0 & 1 & 1 & 4 \\\\
0 & 0 & 0 & 0
\end{bmatrix}$$

Using row operations to convert a matrix into reduced row echelon form
is sometimes called \textbf{Gauss–Jordan elimination}. In this case, the term
\textit{Gaussian elimination} refers to the process until it has reached its
upper triangular, or (unreduced) row echelon form. For computational
reasons, when solving systems of linear equations, it is sometimes
preferable to stop row operations before the matrix is completely
reduced.

\paragraph{Definitions and example of algorithm [definitions\_and\_example\_of\_algorithm]}

The process of row reduction makes use of elementary row operations, and
can be divided into two parts. The first part (sometimes called forward
elimination) reduces a given system to \textit{row echelon form}, from which
one can tell whether there are no solutions, a unique solution, or
infinitely many solutions. The second part (sometimes called back
substitution) continues to use row operations until the solution is
found; in other words, it puts the matrix into \textit{reduced} row echelon
form.

Another point of view, which turns out to be very useful to analyze the
algorithm, is that row reduction produces a matrix decomposition of the
original matrix. The elementary row operations may be viewed as the
multiplication on the left of the original matrix by elementary
matrices. Alternatively, a sequence of elementary operations that
reduces a single row may be viewed as multiplication by a Frobenius
matrix. Then the first part of the algorithm computes an LU
decomposition, while the second part writes the original matrix as the
product of a uniquely determined invertible matrix and a uniquely
determined reduced row echelon matrix.

\paragraph{Row operations [row\_operations]}

There are three types of \textbf{elementary row operations} which may be
performed on the rows of a matrix:

1.  Swap the positions of two rows.
2.  Multiply a row by a non-zero scalar.
3.  Add to one row a scalar multiple of another.

If the matrix is associated to a system of linear equations, then these
operations do not change the solution set. Therefore, if one''''s goal is
to solve a system of linear equations, then using these row operations
could make the problem easier.

\paragraph{Echelon form [echelon\_form]}

For each row in a matrix, if the row does not consist of only zeros,
then the leftmost nonzero entry is called the \textit{leading coefficient} (or
\textit{pivot}) of that row. So if two leading coefficients are in the same
column, then a row operation of type 3 could be used to make one of
those coefficients zero. Then by using the row swapping operation, one
can always order the rows so that for every non-zero row, the leading
coefficient is to the right of the leading coefficient of the row above.
If this is the case, then matrix is said to be in \textbf{row echelon form}.
So the lower left part of the matrix contains only zeros, and all of the
zero rows are below the non-zero rows. The word ""echelon"" is used here
because one can roughly think of the rows being ranked by their size,
with the largest being at the top and the smallest being at the bottom.

For example, the following matrix is in row echelon form, and its
leading coefficients are shown in red:

:   $\begin{bmatrix}
      0 \& \color{red}{\mathbf{2}} \&                     1   \& -1 \\\\
            0 \&                     0   \& \color{red}{\mathbf{3}} \&  1 \\\\
            0 \&                     0   \&                     0   \&  0
    \end{bmatrix}.$

It is in echelon form because the zero row is at the bottom, and the
leading coefficient of the second row (in the third column), is to the
right of the leading coefficient of the first row (in the second
column).

A matrix is said to be in \textbf{reduced row echelon form} if furthermore
all of the leading coefficients are equal to 1 (which can be achieved by
using the elementary row operation of type 2), and in every column
containing a leading coefficient, all of the other entries in that
column are zero (which can be achieved by using elementary row
operations of type 3).

\paragraph{Example of the algorithm [example\_of\_the\_algorithm]}

Suppose the goal is to find and describe the set of solutions to the
following system of linear equations:

:   $\begin{alignedat}{4}
      2x \&{}+{}\& y \&{}-{}\&  z \&{}={}\&   8 \& \qquad (L_1) \\\\
           -3x \&{}-{}\& y \&{}+{}\& 2z \&{}={}\& -11 \& \qquad (L_2) \\\\
           -2x \&{}+{}\& y \&{}+{}\& 2z \&{}={}\&  -3 \& \qquad (L_3)
    \end{alignedat}$

The table below is the row reduction process applied simultaneously to
the system of equations and its associated augmented matrix. In
practice, one does not usually deal with the systems in terms of
equations, but instead makes use of the augmented matrix, which is more
suitable for computer manipulations. The row reduction procedure may be
summarized as follows: eliminate $x$ from all
equations below $L_1$, and then
eliminate $y$ from all equations below
$L_2$. This will put the system into
triangular form. Then, using back-substitution, each unknown can be
solved for.

:   {\| style=""background-color:white;"" class=""wikitable""

\|- ! System of equations !! Row operations !! Augmented matrix \|-
aligned=""center"" \| $\begin{alignedat}{4}
  2x \&{}+{}\& y \&{}-{}\&  z \&{}={}\&   8 \& \\\\
   -3x \&{}-{}\& y \&{}+{}\& 2z \&{}={}\& -11 \& \\\\
   -2x \&{}+{}\& y \&{}+{}\& 2z \&{}={}\&  -3 \&
\end{alignedat}$ \| \| $\left[\begin{array}{rrr|r}
  2 \&  1 \& -1 \&   8 \\\\
   -3 \& -1 \&  2 \& -11 \\\\
   -2 \&  1 \&  2 \&  -3
\end{array}\right]$ \|- aligned=""center"" \| $\begin{alignedat}{4}
 2x \&{}+{}\&          y \&{}-{}\&          z \&{}={}\& 8 \& \\\\
     \&     \& \tfrac12 y \&{}+{}\& \tfrac12 z \&{}={}\& 1 \& \\\\
     \&     \&         2y \&{}+{}\&          z \&{}={}\& 5 \&
\end{alignedat}$ \| $\begin{aligned}
 L_2 + \tfrac32 L_1 \&\to L_2 \\\\
  L_3 +          L_1 \&\to L_3
\end{aligned}$ \| $\left[\begin{array}{rrr|r}
 2 \&      1  \&     -1  \& 8 \\\\
  0 \& \frac12 \& \frac12 \& 1 \\\\
  0 \&      2  \&      1  \& 5
\end{array}\right]$ \|- aligned=""center"" \| $\begin{alignedat}{4}
 2x \&{}+{}\&          y \&{}-{}\&          z \&{}={}\& 8 \& \\\\
     \&     \& \tfrac12 y \&{}+{}\& \tfrac12 z \&{}={}\& 1 \& \\\\
     \&     \&            \&     \&         -z \&{}={}\& 1 \&
\end{alignedat}$ \| $L_3 + -4 L_2 \to L_3$ \| $\left[\begin{array}{rrr|r}
 2 \&      1  \&     -1  \& 8 \\\\
  0 \& \frac12 \& \frac12 \& 1 \\\\
  0 \&      0  \&     -1  \& 1
\end{array}\right]$ \|- \|colspan=3 aligned=""center""\| The matrix is now
in echelon form (also called triangular form) \|- aligned=""center"" \|
$\begin{alignedat}{4}
 2x \&{}+{}\&          y \&     \&   \&{}={}       7  \& \\\\
     \&     \& \tfrac12 y \&     \&   \&{}={} \tfrac32 \& \\\\
     \&     \&            \&{}-{}\& z \&{}={}       1  \&
\end{alignedat}$ \| $\begin{aligned}
 L_2 + \tfrac12 L_3 \&\to L_2 \\\\
  L_1 -          L_3 \&\to L_1
\end{aligned}$ \| $\left[\begin{array}{rrr|r}
 2 \&      1  \&  0 \&      7  \\\\
  0 \& \frac12 \&  0 \& \frac32 \\\\
  0 \&      0  \& -1 \&      1
\end{array}\right]$ \|- aligned=""center"" \| $\begin{alignedat}{4}
 2x \&{}+{}\& y \&\quad\&   \&{}={}\&  7 \& \\\\
     \&     \& y \&\quad\&   \&{}={}\&  3 \& \\\\
     \&     \&   \&\quad\& z \&{}={}\& -1 \&
\end{alignedat}$ \| $\begin{aligned}
 2 L_2 \&\to L_2 \\\\
   -L_3 \&\to L_3
\end{aligned}$ \| $\left[\begin{array}{rrr|r}
 2 \& 1 \& 0 \&  7 \\\\
  0 \& 1 \& 0 \&  3 \\\\
  0 \& 0 \& 1 \& -1
\end{array}\right]$ \|- aligned=""center"" \| $\begin{alignedat}{4}
 x \&\quad\&   \&\quad\&   \&{}={}\&  2 \& \\\\
    \&\quad\& y \&\quad\&   \&{}={}\&  3 \& \\\\
    \&\quad\&   \&\quad\& z \&{}={}\& -1 \&
\end{alignedat}$ \| $\begin{aligned}
          L_1 - L_2 \&\to L_1 \\\\
           \tfrac12 L_1       \&\to L_1
\end{aligned}$ \| $\left[\begin{array}{rrr|r}
 1 \& 0 \& 0 \&  2 \\\\
  0 \& 1 \& 0 \&  3 \\\\
  0 \& 0 \& 1 \& -1
\end{array}\right]$ \|}

The second column describes which row operations have just been
performed. So for the first step, the $x$ is
eliminated from $L_2$ by adding
$\tfrac{3}{2}$\textit{L}\textasciitilde{}1\textasciitilde{} to $L_2$. Next,
$x$ is eliminated from
$L_3$ by adding
$L_1$ to
$L_3$. These row operations are
labelled in the table as

:   $\begin{aligned}
     L_2 + \tfrac32 L_1 \&\to L_2, \\\\
          L_3 +          L_1 \&\to L_3.
    \end{aligned}$

Once $y$ is also eliminated from the third row, the
result is a system of linear equations in triangular form, and so the
first part of the algorithm is complete. From a computational point of
view, it is faster to solve the variables in reverse order, a process
known as back-substitution. One sees the solution is \textit{z} = -1, \textit{y} = 3,
and \textit{x} = 2. So there is a unique solution to the original system of
equations.

Instead of stopping once the matrix is in echelon form, one could
continue until the matrix is in \textit{reduced} row echelon form, as it is
done in the table. The process of row reducing until the matrix is
reduced is sometimes referred to as \textbf{Gauss–Jordan elimination}, to
distinguish it from stopping after reaching echelon form.

\paragraph{Applications}

Historically, the first application of the row reduction method is for
solving systems of linear equations. Here are some other important
applications of the algorithm.

\paragraph{Computing determinants [computing\_determinants]}

To explain how Gaussian elimination allows the computation of the
determinant of a square matrix, we have to recall how the elementary row
operations change the determinant:

\begin{itemize}
    \item Swapping two rows multiplies the determinant by -1
    \item Multiplying a row by a nonzero scalar multiplies the determinant by
\end{itemize}
    the same scalar
\begin{itemize}
    \item Adding to one row a scalar multiple of another does not change the
\end{itemize}
    determinant.

If Gaussian elimination applied to a square matrix
$A$ produces a row echelon matrix
$B$, let $d$ be the product of
the scalars by which the determinant has been multiplied, using the
above rules. Then the determinant of $A$ is the
quotient by $d$ of the product of the elements of
the diagonal of $B$:

$$\det(A) = \frac{\prod\operatorname{diag}(B)}{d}.$$

Computationally, for an $n'''''''' × ''''''''n''''''''$ matrix,
this method needs only `{{math|O(''''''''n''''''''3)$
arithmetic operations, while using Leibniz formula for determinants
requires `{{math|O(''''''''n''''''''!)$ operations (number of summands
in the formula), and recursive Laplace expansion requires
`{{math|O(2''''''''n'''''''')$ operations (number of
sub-determinants to compute, if none is computed twice). Even on the
fastest computers, these two methods are impractical or almost
impracticable for $n''''''''$ above 20.$$

Finding the inverse of a matrix

A variant of Gaussian elimination called Gauss–Jordan elimination can be
used for finding the inverse of a matrix, if it exists. If
$A''''''''$ is an $n'''''''' × ''''''''n''''''''$
square matrix, then one can use row reduction to compute its inverse
matrix, if it exists. First, the $n'''''''' × ''''''''n''''''''$
identity matrix is augmented to the right of
$A''''''''$, forming an
$n'''''''' × 2''''''''n''''''''$ block matrix
`{{math|[''''''''A'''''''' {{!$ *I*\]}}. Now through application of
elementary row operations, find the reduced echelon form of this
$n'''''''' × 2''''''''n''''''''$ matrix. The matrix
$A''''''''$ is invertible if and only if the left block
can be reduced to the identity matrix $I''''''''$; in
this case the right block of the final matrix is
$A''''''''-1$. If the algorithm is unable to
reduce the left block to $I''''''''$, then
$A''''''''$ is not invertible.

For example, consider the following matrix:

:   $A =
     \begin{bmatrix}
      2 & -1 &  0 \\\\
           -1 &  2 & -1 \\\\
            0 & -1 &  2
     \end{bmatrix}.$

To find the inverse of this matrix, one takes the following matrix
augmented by the identity and row-reduces it as a 3 × 6 matrix:

:   $[ A | I ] = 
     \left[\begin{array}{ccc|ccc}
       2 & -1 &  0 & 1 & 0 & 0 \\\\
             -1 &  2 & -1 & 0 & 1 & 0 \\\\
              0 & -1 &  2 & 0 & 0 & 1
     \end{array}\right].$

By performing row operations, one can check that the reduced row echelon
form of this augmented matrix is

:   $[ I | B ] = 
     \left[\begin{array}{rrr|rrr}
      1 & 0 & 0 & \frac34 & \frac12 & \frac14 \\\\
            0 & 1 & 0 & \frac12 &      1  & \frac12 \\\\
            0 & 0 & 1 & \frac14 & \frac12 & \frac34
     \end{array}\right].$

One can think of each row operation as the left product by an elementary
matrix. Denoting by $B''''''''$ the product of these
elementary matrices, we showed, on the left, that *BA* = *I*, and
therefore, *B* = *A*^-1^. On the right, we kept a record of *BI* = *B*,
which we know is the inverse desired. This procedure for finding the
inverse works for square matrices of any size.

Computing ranks and bases

The Gaussian elimination algorithm can be applied to any
$m'''''''' × ''''''''n''''''''$ matrix $A$. In
this way, for example, some 6 × 9 matrices can be transformed to a
matrix that has a row echelon form like

$$T=
\begin{bmatrix}
a \& \textit{ \& } \& \textit{\& } \& \textit{ \& } \& \textit{ \& } \\\\
0 \& 0 \& b \& \textit{ \& } \& \textit{ \& } \& \textit{ \& } \\\\
0 \& 0 \& 0 \& c \& \textit{ \& } \& \textit{ \& } \& * \\\\
0 \& 0 \& 0 \& 0 \& 0 \& 0 \& d \& \textit{ \& } \\\\
0 \& 0 \& 0 \& 0 \& 0 \& 0 \& 0 \& 0 \& e \\\\
0 \& 0 \& 0 \& 0 \& 0 \& 0 \& 0 \& 0 \& 0
\end{bmatrix},$$ where the stars are arbitrary entries, and
$a'''''''', ''''''''b'''''''', ''''''''c'''''''', ''''''''d'''''''', ''''''''e''''''''$ are nonzero
entries. This echelon matrix $T$ contains a wealth
of information about $A$: the rank of
$A$ is 5, since there are 5 nonzero rows in
$T$; the vector space spanned by the columns of
$A$ has a basis consisting of its columns 1, 3, 4,
7 and 9 (the columns with
$a'''''''', ''''''''b'''''''', ''''''''c'''''''', ''''''''d'''''''', ''''''''e''''''''$ in
$T$), and the stars show how the other columns of
$A$ can be written as linear combinations of the
basis columns. This is a consequence of the distributivity of the dot
product in the expression of a linear map as a matrix.

All of this applies also to the reduced row echelon form, which is a
particular row echelon format.

\paragraph{Computational efficiency [computational_efficiency]}

The number of arithmetic operations required to perform row reduction is
one way of measuring the algorithm''''s computational efficiency. For
example, to solve a system of $n''''''''$ equations for
$n''''''''$ unknowns by performing row operations on
the matrix until it is in echelon form, and then solving for each
unknown in reverse order, requires
$n''''''''(''''''''n'''''''' + 1)/2$ divisions,
`{{math|(2''''''''n''''''''3 + 3''''''''n''''''''2 - 5''''''''n'''''''')/6$
multiplications, and
`{{math|(2''''''''n''''''''3 + 3''''''''n''''''''2 - 5''''''''n'''''''')/6$
subtractions, for a total of approximately
`{{math|2''''''''n''''''''3/3$ operations. Thus it has
arithmetic complexity of `{{math|O(''''''''n''''''''3)$;
see Big O notation.

This arithmetic complexity is a good measure of the time needed for the
whole computation when the time for each arithmetic operation is
approximately constant. This is the case when the coefficients are
represented by floating-point numbers or when they belong to a finite
field. If the coefficients are integers or rational numbers exactly
represented, the intermediate entries can grow exponentially large, so
the bit complexity is exponential. However, there is a variant of
Gaussian elimination, called the Bareiss algorithm, that avoids this
exponential growth of the intermediate entries and, with the same
arithmetic complexity of `{{math|O(''''''''n''''''''3)$,
has a bit complexity of `{{math|O(''''''''n''''''''5)$.

This algorithm can be used on a computer for systems with thousands of
equations and unknowns. However, the cost becomes prohibitive for
systems with millions of equations. These large systems are generally
solved using iterative methods. Specific methods exist for systems whose
coefficients follow a regular pattern (see system of linear equations).

To put an $n'''''''' × ''''''''n''''''''$ matrix into reduced
echelon form by row operations, one needs
$n''''''''3$ arithmetic operations, which is
approximately 50\% more computation steps.

One possible problem is numerical instability, caused by the possibility
of dividing by very small numbers. If, for example, the leading
coefficient of one of the rows is very close to zero, then to row-reduce
the matrix, one would need to divide by that number. This means that any
error existed for the number that was close to zero would be amplified.
Gaussian elimination is numerically stable for diagonally dominant or
positive-definite matrices. For general matrices, Gaussian elimination
is usually considered to be stable, when using partial pivoting, even
though there are examples of stable matrices for which it is unstable.

\paragraph{Generalizations}

Gaussian elimination can be performed over any field, not just the real
numbers.

Buchberger''''s algorithm is a generalization of Gaussian elimination to
systems of polynomial equations. This generalization depends heavily on
the notion of a monomial order. The choice of an ordering on the
variables is already implicit in Gaussian elimination, manifesting as
the choice to work from left to right when selecting pivot positions.

Computing the rank of a tensor of order greater than 2 is NP-hard.
Therefore, if `{{math|P ? NP$, there cannot be a
polynomial time analog of Gaussian elimination for higher-order tensors
(matrices are array representations of order-2 tensors).

\paragraph{Pseudocode}

As explained above, Gaussian elimination transforms a given
$m'''''''' × ''''''''n''''''''$ matrix $A$
into a matrix in row-echelon form.

In the following pseudocode, `A[i, j]` denotes the entry of the matrix
$A$ in row $i$ and column
$j$ with the indices starting from 1. The
transformation is performed \textit{in place}, meaning that the original matrix
is lost for being eventually replaced by its row-echelon form.

`h := 1 /\textit{ `}`Initialization of the pivot row`\textit{` }/`  
`k := 1 /\textit{ `}`Initialization of the pivot column`\textit{` }/`  
  
\textbf{`while`}` h = m `\textbf{`and`}` k = n`  
`    /\textit{ `}`Find the k-th pivot:`\textit{` }/`  
`    i_max := argmax (i = h ... m, abs(A[i, k]))`  
`    `\textbf{`if`}` A[i_max, k] = 0`  
`        /\textit{ `}`No pivot in this column, pass to next column`\textit{` }/`  
`        k := k + 1`  
`    `\textbf{`else`}  
`         `\textbf{`swap rows`}`(h, i_max)`  
`         /\textit{ `}`Do for all rows below pivot:`\textit{` }/`  
`         `\textbf{`for`}` i = h + 1 ... m:`  
`             f := A[i, k] / A[h, k]`  
`             /\textit{ `}`Fill with zeros the lower part of pivot column:`\textit{` }/`  
`             A[i, k] := 0`  
`             /\textit{ `}`Do for all remaining elements in current row:`\textit{` }/`  
`             `\textbf{`for`}` j = k + 1 ... n:`  
`                 A[i, j] := A[i, j] - A[h, j] * f`  
`         /\textit{ `}`Increase pivot row and column`\textit{` }/`  
`         h := h + 1`  
`         k := k + 1`

This algorithm differs slightly from the one discussed earlier, by
choosing a pivot with largest absolute value. Such a \textit{partial pivoting}
may be required if, at the pivot place, the entry of the matrix is zero.
In any case, choosing the largest possible absolute value of the pivot
improves the numerical stability of the algorithm, when floating point
is used for representing numbers.

Upon completion of this procedure the matrix will be in row echelon form
and the corresponding system may be solved by back substitution.

\paragraph{Resources [resources_2]}

\begin{itemize}
    \item \href{https://openstax.org/books/college-algebra/pages/7-6-solving-systems-with-gaussian-elimination}{Solving Systems with Gaussian
\end{itemize}
    Elimination},
    OpenStax
\begin{itemize}
    \item \href{https://tutorial.math.lamar.edu/classes/alg/augmentedmatrix.aspx}{Augmented Matrices and Gauss-Jordan
\end{itemize}
    Elimination},
    Paul''''s Online Notes

\paragraph{Licensing [licensing_2]}

Content obtained and/or adapted from:

\begin{itemize}
    \item \href{https://en.wikipedia.org/wiki/Gaussian_elimination}{Gaussian elimination,
\end{itemize}
    Wikipedia} under
    a CC BY-SA license"
\paragraph{Resources [resources_1]}

\begin{itemize}
    \item \href{https://www.unf.edu/\textasciitilde{}mzhan/chapter3.pdf}{Brief Introduction to Vectors and
\end{itemize}
    Matrices}, University of
    North Florida
\begin{itemize}
    \item \href{https://www.usna.edu/Users/math/uhan/sm286a/lessons/02.pdf}{Introduction to Matrices and
\end{itemize}
    Vectors},
    United States Naval Academy

\paragraph{Licensing [licensing_1]}

Content obtained and/or adapted from:

\begin{itemize}
    \item \href{https://en.wikibooks.org/wiki/Linear_Algebra/Matrices_and_Vectors/}{Matrices and Vectors, Wikibooks: Linear
\end{itemize}
    Algebra}
    under a CC BY-SA license"

| \textbf{System of equations} | \textbf{Row operations} | \textbf{Augmented matrix} | 
|-------------------------|--------------------|----------------------|
| $\begin{alignedat}{4} 2x &{}+{}& y &{}-{}& z &{}={}& 8 & \\\\ -3x &{}-{}& y &{}+{}& 2z &{}={}& -11 & \\\\ -2x &{}+{}& y &{}+{}& 2z &{}={}& -3 & \end{alignedat}$ | | $\begin{array}{rrr|r} 2 & 1 & -1 & 8 \\\\ -3 & -1 & 2 & -11 \\\\ -2 & 1 & 2 & -3 \end{array}$ |
| $\begin{alignedat}{4} 2x &{}+{}& y &{}-{}& z &{}={}& 8 & \\\\ & & \tfrac{1}{2} y &{}+{}& \tfrac{1}{2} z &{}={}& 1 & \\\\ & & 2y &{}+{}& z &{}={}& 5 & \end{alignedat}$ | $\begin{aligned} L_2 + \tfrac{3}{2} L_1 \&\to L_2 \\\\ L_3 + L_1 \&\to L_3 \end{aligned}\) | \(\left[\begin{array}{rrr|r} 2 \& 1 \& -1 \& 8 \\\\ 0 \& \frac{1}{2} \& \frac{1}{2} \& 1 \\\\ 0 \& 2 \& 1 \& 5 \end{array}\right]\) |
| $\begin{alignedat}{4} 2x &{}+{}& y &{}-{}& z &{}={}& 8 & \\\\ & & \tfrac{1}{2} y &{}+{}& \tfrac{1}{2} z &{}={}& 1 & \\\\ & & &{}-{}& z &{}={}& 1 & \end{alignedat}$ | \(L_3 + -4 L_2 \to L_3\) | \(\left[\begin{array}{rrr|r} 2 \& 1 \& -1 \& 8 \\\\ 0 \& \frac{1}{2} \& \frac{1}{2} \& 1 \\\\ 0 \& 0 \& -1 \& 1 \end{array}\right]\) |
| \multicolumn{3}{|c|}{The matrix is now in echelon form (also called triangular form)} |
| $\begin{alignedat}{4} 2x &{}+{}& y & & &{}={} 7 & \\\\ & & \tfrac{1}{2} y & & &{}={} \tfrac{3}{2} & \\\\ & & &{}-{}& z &{}={} 1 & \end{alignedat}$ | $\begin{aligned} L_2 + \tfrac{1}{2} L_3 \&\to L_2 \\\\ L_1 - L_3 \&\to L_1 \end{aligned}\) | \(\left[\begin{array}{rrr|r} 2 \& 1 \& 0 \& 7 \\\\ 0 \& \frac{1}{2} \& 0 \& \frac{3}{2} \\\\ 0 \& 0 \& -1 \& 1 \end{array}\right]\) |
| $\begin{alignedat}{4} 2x &{}+{}& y &\quad& &{}={}& 7 & \\\\ & & y &\quad& &{}={}& 3 & \\\\ & & &\quad& z &{}={}& -1 & \end{alignedat}$ | $\begin{aligned} 2 L_2 \&\to L_2 \\\\ -L_3 \&\to L_3 \end{aligned}\) | \(\left[\begin{array}{rrr|r} 2 \& 1 \& 0 \& 7 \\\\ 0 \& 1 \& 0 \& 3 \\\\ 0 \& 0 \& 1 \& -1 \end{array}\right]\) |
| $\begin{alignedat}{4} x &\quad& &\quad& &{}={}& 2 & \\\\ &\quad& y &\quad& &{}={}& 3 & \\\\ &\quad& &\quad& z &{}={}& -1 & \end{alignedat}$ | $\begin{aligned} L_1 - L_2 \&\to L_1 \\\\ \tfrac{1}{2} L_1 \&\to L_1 \end{aligned}\) | \(\left[\begin{array}{rrr|r} 1 \& 0 \& 0 \& 2 \\\\ 0 \& 1 \& 0 \& 3 \\\\ 0 \& 0 \& 1 \& -1 \end{array}\right]\) |

| Left-aligned | Center-aligned | Right-aligned |
| :---         |     :---:      |          ---: |
| git status   | git status     | git status    |
| git diff     | git diff       | git diff      |

%%===============================================================
\section{Gauss-Jordan Elimination}
\label{sec:gauss-jordan-elimination}
%%===============================================================

In mathematics, \textbf{Gaussian elimination}, also known as **row
reduction**, is an algorithm for solving systems of linear equations. It
consists of a sequence of operations performed on the corresponding
matrix of coefficients. This method can also be used to compute the rank
of a matrix, the determinant of a square matrix, and the inverse of an
invertible matrix. The method is named after Carl Friedrich Gauss
(1777–1855) although some special cases of the method—albeit presented
without proof—were known to Chinese mathematicians as early as circa
179 CE.

To perform row reduction on a matrix, one uses a sequence of elementary
row operations to modify the matrix until the lower left-hand corner of
the matrix is filled with zeros, as much as possible. There are three
types of elementary row operations:

\begin{itemize}
    \item Swapping two rows,
    \item Multiplying a row by a nonzero number,
    \item Adding a multiple of one row to another row. (subtraction can be
\end{itemize}
    achieved by multiplying one row with -1 and adding the result to
    another row)

Using these operations, a matrix can always be transformed into an upper
triangular matrix, and in fact one that is in row echelon form. Once all
of the leading coefficients (the leftmost nonzero entry in each row) are
1, and every column containing a leading coefficient has zeros
elsewhere, the matrix is said to be in reduced row echelon form. This
final form is unique; in other words, it is independent of the sequence
of row operations used. For example, in the following sequence of row
operations (where two elementary operations on different rows are done
at the first and third steps), the third and fourth matrices are the
ones in row echelon form, and the final matrix is the unique reduced row
echelon form.

$$\begin{bmatrix}
1 & 3 & 1 & 9 \\\\
1 & 1 & -1 & 1 \\\\
3 & 11 & 5 & 35
\end{bmatrix}\to
\begin{bmatrix}
1 & 3 & 1 & 9 \\\\
0 & -2 & -2 & -8 \\\\
0 & 2 & 2 & 8
\end{bmatrix}\to
\begin{bmatrix}
1 & 3 & 1 & 9 \\\\
0 & -2 & -2 & -8 \\\\
0 & 0 & 0 & 0
\end{bmatrix}\to
\begin{bmatrix}
1 & 0 & -2 & -3 \\\\
0 & 1 & 1 & 4 \\\\
0 & 0 & 0 & 0
\end{bmatrix}$$

Using row operations to convert a matrix into reduced row echelon form
is sometimes called \textbf{Gauss–Jordan elimination}. In this case, the term
\textit{Gaussian elimination} refers to the process until it has reached its
upper triangular, or (unreduced) row echelon form. For computational
reasons, when solving systems of linear equations, it is sometimes
preferable to stop row operations before the matrix is completely
reduced.

\paragraph{Definitions and example of algorithm [definitions_and_example_of_algorithm]}

The process of row reduction makes use of elementary row operations, and
can be divided into two parts. The first part (sometimes called forward
elimination) reduces a given system to \textit{row echelon form}, from which
one can tell whether there are no solutions, a unique solution, or
infinitely many solutions. The second part (sometimes called back
substitution) continues to use row operations until the solution is
found; in other words, it puts the matrix into \textit{reduced} row echelon
form.

Another point of view, which turns out to be very useful to analyze the
algorithm, is that row reduction produces a matrix decomposition of the
original matrix. The elementary row operations may be viewed as the
multiplication on the left of the original matrix by elementary
matrices. Alternatively, a sequence of elementary operations that
reduces a single row may be viewed as multiplication by a Frobenius
matrix. Then the first part of the algorithm computes an LU
decomposition, while the second part writes the original matrix as the
product of a uniquely determined invertible matrix and a uniquely
determined reduced row echelon matrix.

\paragraph{Row operations [row_operations]}

There are three types of \textbf{elementary row operations} which may be
performed on the rows of a matrix:

1.  Swap the positions of two rows.
2.  Multiply a row by a non-zero scalar.
3.  Add to one row a scalar multiple of another.

If the matrix is associated to a system of linear equations, then these
operations do not change the solution set. Therefore, if one''''s goal is
to solve a system of linear equations, then using these row operations
could make the problem easier.

\paragraph{Echelon form [echelon_form]}

For each row in a matrix, if the row does not consist of only zeros,
then the leftmost nonzero entry is called the \textit{leading coefficient} (or
\textit{pivot}) of that row. So if two leading coefficients are in the same
column, then a row operation of type 3 could be used to make one of
those coefficients zero. Then by using the row swapping operation, one
can always order the rows so that for every non-zero row, the leading
coefficient is to the right of the leading coefficient of the row above.
If this is the case, then matrix is said to be in \textbf{row echelon form}.
So the lower left part of the matrix contains only zeros, and all of the
zero rows are below the non-zero rows. The word ""echelon"" is used here
because one can roughly think of the rows being ranked by their size,
with the largest being at the top and the smallest being at the bottom.

For example, the following matrix is in row echelon form, and its
leading coefficients are shown in red:

$\begin{bmatrix}
      0 \& \color{red}{\mathbf{2}} \&                     1   \& -1 \\\\
            0 \&                     0   \& \color{red}{\mathbf{3}} \&  1 \\\\
            0 \&                     0   \&                     0   \&  0
    \end{bmatrix}.$

It is in echelon form because the zero row is at the bottom, and the
leading coefficient of the second row (in the third column), is to the
right of the leading coefficient of the first row (in the second
column).

A matrix is said to be in \textbf{reduced row echelon form} if furthermore
all of the leading coefficients are equal to 1 (which can be achieved by
using the elementary row operation of type 2), and in every column
containing a leading coefficient, all of the other entries in that
column are zero (which can be achieved by using elementary row
operations of type 3).

\paragraph{Example of the algorithm [example_of_the_algorithm]}

Suppose the goal is to find and describe the set of solutions to the
following system of linear equations:

$\begin{alignedat}{4}
      2x \&{}+{}\& y \&{}-{}\&  z \&{}={}\&   8 \& \qquad (L_1) \\\\
           -3x \&{}-{}\& y \&{}+{}\& 2z \&{}={}\& -11 \& \qquad (L_2) \\\\
           -2x \&{}+{}\& y \&{}+{}\& 2z \&{}={}\&  -3 \& \qquad (L_3)
    \end{alignedat}$

The table below is the row reduction process applied simultaneously to
the system of equations and its associated augmented matrix. In
practice, one does not usually deal with the systems in terms of
equations, but instead makes use of the augmented matrix, which is more
suitable for computer manipulations. The row reduction procedure may be
summarized as follows: eliminate $x$ from all
equations below $L_1$, and then
eliminate $y$ from all equations below
$L_2$. This will put the system into
triangular form. Then, using back-substitution, each unknown can be
solved for.

{\| style=""background-color:white;"" class=""wikitable""

\|- ! System of equations !! Row operations !! Augmented matrix \|-
aligned=""center"" \| $\begin{alignedat}{4}
  2x \&{}+{}\& y \&{}-{}\&  z \&{}={}\&   8 \& \\\\
   -3x \&{}-{}\& y \&{}+{}\& 2z \&{}={}\& -11 \& \\\\
   -2x \&{}+{}\& y \&{}+{}\& 2z \&{}={}\&  -3 \&
\end{alignedat}$ \| \| $\left[\begin{array}{rrr|r}
  2 \&  1 \& -1 \&   8 \\\\
   -3 \& -1 \&  2 \& -11 \\\\
   -2 \&  1 \&  2 \&  -3
\end{array}\right]$ \|- aligned=""center"" \| $\begin{alignedat}{4}
 2x \&{}+{}\&          y \&{}-{}\&          z \&{}={}\& 8 \& \\\\
     \&     \& \tfrac12 y \&{}+{}\& \tfrac12 z \&{}={}\& 1 \& \\\\
     \&     \&         2y \&{}+{}\&          z \&{}={}\& 5 \&
\end{alignedat}$ \| $\begin{aligned}
 L_2 + \tfrac32 L_1 \&\to L_2 \\\\
  L_3 +          L_1 \&\to L_3
\end{aligned}$ \| $\left[\begin{array}{rrr|r}
 2 \&      1  \&     -1  \& 8 \\\\
  0 \& \frac12 \& \frac12 \& 1 \\\\
  0 \&      2  \&      1  \& 5
\end{array}\right]$ \|- aligned=""center"" \| $\begin{alignedat}{4}
 2x \&{}+{}\&          y \&{}-{}\&          z \&{}={}\& 8 \& \\\\
     \&     \& \tfrac12 y \&{}+{}\& \tfrac12 z \&{}={}\& 1 \& \\\\
     \&     \&            \&     \&         -z \&{}={}\& 1 \&
\end{alignedat}$ \| $L_3 + -4 L_2 \to L_3$ \| $\left[\begin{array}{rrr|r}
 2 \&      1  \&     -1  \& 8 \\\\
  0 \& \frac12 \& \frac12 \& 1 \\\\
  0 \&      0  \&     -1  \& 1
\end{array}\right]$ \|- \|colspan=3 aligned=""center""\| The matrix is now
in echelon form (also called triangular form) \|- aligned=""center"" \|
$\begin{alignedat}{4}
 2x \&{}+{}\&          y \&     \&   \&{}={}       7  \& \\\\
     \&     \& \tfrac12 y \&     \&   \&{}={} \tfrac32 \& \\\\
     \&     \&            \&{}-{}\& z \&{}={}       1  \&
\end{alignedat}$ \| $\begin{aligned}
 L_2 + \tfrac12 L_3 \&\to L_2 \\\\
  L_1 -          L_3 \&\to L_1
\end{aligned}$ \| $\left[\begin{array}{rrr|r}
 2 \&      1  \&  0 \&      7  \\\\
  0 \& \frac12 \&  0 \& \frac32 \\\\
  0 \&      0  \& -1 \&      1
\end{array}\right]$ \|- aligned=""center"" \| $\begin{alignedat}{4}
 2x \&{}+{}\& y \&\quad\&   \&{}={}\&  7 \& \\\\
     \&     \& y \&\quad\&   \&{}={}\&  3 \& \\\\
     \&     \&   \&\quad\& z \&{}={}\& -1 \&
\end{alignedat}$ \| $\begin{aligned}
 2 L_2 \&\to L_2 \\\\
   -L_3 \&\to L_3
\end{aligned}$ \| $\left[\begin{array}{rrr|r}
 2 \& 1 \& 0 \&  7 \\\\
  0 \& 1 \& 0 \&  3 \\\\
  0 \& 0 \& 1 \& -1
\end{array}\right]$ \|- aligned=""center"" \| $\begin{alignedat}{4}
 x \&\quad\&   \&\quad\&   \&{}={}\&  2 \& \\\\
    \&\quad\& y \&\quad\&   \&{}={}\&  3 \& \\\\
    \&\quad\&   \&\quad\& z \&{}={}\& -1 \&
\end{alignedat}$ \| $\begin{aligned}
          L_1 - L_2 \&\to L_1 \\\\
           \tfrac12 L_1       \&\to L_1
\end{aligned}$ \| $\left[\begin{array}{rrr|r}
 1 \& 0 \& 0 \&  2 \\\\
  0 \& 1 \& 0 \&  3 \\\\
  0 \& 0 \& 1 \& -1
\end{array}\right]$ \|}

The second column describes which row operations have just been
performed. So for the first step, the $x$ is
eliminated from $L_2$ by adding
$\tfrac{3}{2}$\textit{L}\textasciitilde{}1\textasciitilde{} to $L_2$. Next,
$x$ is eliminated from
$L_3$ by adding
$L_1$ to
$L_3$. These row operations are
labelled in the table as

$\begin{aligned}
     L_2 + \tfrac32 L_1 \&\to L_2, \\\\
          L_3 +          L_1 \&\to L_3.
    \end{aligned}$

Once $y$ is also eliminated from the third row, the
result is a system of linear equations in triangular form, and so the
first part of the algorithm is complete. From a computational point of
view, it is faster to solve the variables in reverse order, a process
known as back-substitution. One sees the solution is \textit{z} = -1, \textit{y} = 3,
and \textit{x} = 2. So there is a unique solution to the original system of
equations.

Instead of stopping once the matrix is in echelon form, one could
continue until the matrix is in \textit{reduced} row echelon form, as it is
done in the table. The process of row reducing until the matrix is
reduced is sometimes referred to as \textbf{Gauss–Jordan elimination}, to
distinguish it from stopping after reaching echelon form.

\paragraph{Applications}

Historically, the first application of the row reduction method is for
solving systems of linear equations. Here are some other important
applications of the algorithm.

\paragraph{Computing determinants [computing_determinants]}

To explain how Gaussian elimination allows the computation of the
determinant of a square matrix, we have to recall how the elementary row
operations change the determinant:

\begin{itemize}
    \item Swapping two rows multiplies the determinant by -1
    \item Multiplying a row by a nonzero scalar multiplies the determinant by
\end{itemize}
    the same scalar
\begin{itemize}
    \item Adding to one row a scalar multiple of another does not change the
\end{itemize}
    determinant.

If Gaussian elimination applied to a square matrix
$A$ produces a row echelon matrix
$B$, let $d$ be the product of
the scalars by which the determinant has been multiplied, using the
above rules. Then the determinant of $A$ is the
quotient by $d$ of the product of the elements of
the diagonal of $B$:

$$\det(A) = \frac{\prod\operatorname{diag}(B)}{d}.$$

Computationally, for an $n × n$ matrix,
this method needs only $O(n^3)$
arithmetic operations, while using Leibniz formula for determinants
requires $O(n!)$ operations (number of summands
in the formula), and recursive Laplace expansion requires
$O(2^n)$ operations (number of
sub-determinants to compute, if none is computed twice). Even on the
fastest computers, these two methods are impractical or almost
impracticable for $n$ above 20.$$

Finding the inverse of a matrix

A variant of Gaussian elimination called Gauss–Jordan elimination can be
used for finding the inverse of a matrix, if it exists. If
$A$ is an $n × n$
square matrix, then one can use row reduction to compute its inverse
matrix, if it exists. First, the $n × n$
identity matrix is augmented to the right of
$A$, forming an
$n × 2n$ block matrix
$[A | I]$. Now through application of
elementary row operations, find the reduced echelon form of this
$n × 2n$ matrix. The matrix
$A$ is invertible if and only if the left block
can be reduced to the identity matrix $I$; in
this case the right block of the final matrix is
$A^-1$. If the algorithm is unable to
reduce the left block to $I$, then
$A$ is not invertible.

For example, consider the following matrix:

$A =
     \begin{bmatrix}
      2 & -1 &  0 \\\\
           -1 &  2 & -1 \\\\
            0 & -1 &  2
     \end{bmatrix}.$

To find the inverse of this matrix, one takes the following matrix
augmented by the identity and row-reduces it as a 3 × 6 matrix:

$[ A | I ] = 
     \left[\begin{array}{ccc|ccc}
       2 & -1 &  0 & 1 & 0 & 0 \\\\
             -1 &  2 & -1 & 0 & 1 & 0 \\\\
              0 & -1 &  2 & 0 & 0 & 1
     \end{array}\right].$

By performing row operations, one can check that the reduced row echelon
form of this augmented matrix is

$[ I | B ] = 
     \left[\begin{array}{rrr|rrr}
      1 & 0 & 0 & \frac34 & \frac12 & \frac14 \\\\
            0 & 1 & 0 & \frac12 &      1  & \frac12 \\\\
            0 & 0 & 1 & \frac14 & \frac12 & \frac34
     \end{array}\right].$

One can think of each row operation as the left product by an elementary
matrix. Denoting by $B$ the product of these
elementary matrices, we showed, on the left, that *BA* = *I*, and
therefore, *B* = $A^{-1}$. On the right, we kept a record of *BI* = *B*,
which we know is the inverse desired. This procedure for finding the
inverse works for square matrices of any size.

Computing ranks and bases

The Gaussian elimination algorithm can be applied to any
$m × n$ matrix $A$. In
this way, for example, some 6 × 9 matrices can be transformed to a
matrix that has a row echelon form like

$$T=
\begin{bmatrix}
a \& \textit{ \& } \& \textit{\& } \& \textit{ \& } \& \textit{ \& } \\\\
0 \& 0 \& b \& \textit{ \& } \& \textit{ \& } \& \textit{ \& } \\\\
0 \& 0 \& 0 \& c \& \textit{ \& } \& \textit{ \& } \& * \\\\
0 \& 0 \& 0 \& 0 \& 0 \& 0 \& d \& \textit{ \& } \\\\
0 \& 0 \& 0 \& 0 \& 0 \& 0 \& 0 \& 0 \& e \\\\
0 \& 0 \& 0 \& 0 \& 0 \& 0 \& 0 \& 0 \& 0
\end{bmatrix},$$ where the stars are arbitrary entries, and
$a, b, c, d, e$ are nonzero
entries. This echelon matrix $T$ contains a wealth
of information about $A$: the rank of
$A$ is 5, since there are 5 nonzero rows in
$T$; the vector space spanned by the columns of
$A$ has a basis consisting of its columns 1, 3, 4,
7 and 9 (the columns with
$a, b, c, d, e$ in
$T$), and the stars show how the other columns of
$A$ can be written as linear combinations of the
basis columns. This is a consequence of the distributivity of the dot
product in the expression of a linear map as a matrix.

All of this applies also to the reduced row echelon form, which is a
particular row echelon format.

\paragraph{Computational efficiency [computational\_efficiency]}

The number of arithmetic operations required to perform row reduction is
one way of measuring the algorithm''''s computational efficiency. For
example, to solve a system of $n$ equations for
$n$ unknowns by performing row operations on
the matrix until it is in echelon form, and then solving for each
unknown in reverse order, requires
$n(n + 1)/2$ divisions,
$(2n^3 + 3n^2 - 5n)/6$
multiplications, and
$(2n^3 + 3n^2 - 5n)/6$
subtractions, for a total of approximately
$2n^3/3$ operations. Thus it has
arithmetic complexity of $O(n^3)$;
see Big O notation.

This arithmetic complexity is a good measure of the time needed for the
whole computation when the time for each arithmetic operation is
approximately constant. This is the case when the coefficients are
represented by floating-point numbers or when they belong to a finite
field. If the coefficients are integers or rational numbers exactly
represented, the intermediate entries can grow exponentially large, so
the bit complexity is exponential. However, there is a variant of
Gaussian elimination, called the Bareiss algorithm, that avoids this
exponential growth of the intermediate entries and, with the same
arithmetic complexity of $O(n^3)$,
has a bit complexity of $O(n^5)$.

This algorithm can be used on a computer for systems with thousands of
equations and unknowns. However, the cost becomes prohibitive for
systems with millions of equations. These large systems are generally
solved using iterative methods. Specific methods exist for systems whose
coefficients follow a regular pattern (see system of linear equations).

To put an $n × n$ matrix into reduced
echelon form by row operations, one needs
$n^3$ arithmetic operations, which is
approximately 50\% more computation steps.

One possible problem is numerical instability, caused by the possibility
of dividing by very small numbers. If, for example, the leading
coefficient of one of the rows is very close to zero, then to row-reduce
the matrix, one would need to divide by that number. This means that any
error existed for the number that was close to zero would be amplified.
Gaussian elimination is numerically stable for diagonally dominant or
positive-definite matrices. For general matrices, Gaussian elimination
is usually considered to be stable, when using partial pivoting, even
though there are examples of stable matrices for which it is unstable.

\paragraph{Generalizations}

Gaussian elimination can be performed over any field, not just the real
numbers.

Buchberger''''s algorithm is a generalization of Gaussian elimination to
systems of polynomial equations. This generalization depends heavily on
the notion of a monomial order. The choice of an ordering on the
variables is already implicit in Gaussian elimination, manifesting as
the choice to work from left to right when selecting pivot positions.

Computing the rank of a tensor of order greater than 2 is NP-hard.
Therefore, if $P \neq NP$, there cannot be a
polynomial time analog of Gaussian elimination for higher-order tensors
(matrices are array representations of order-2 tensors).

\paragraph{Pseudocode}

As explained above, Gaussian elimination transforms a given
$m × n$ matrix $A$
into a matrix in row-echelon form.

In the following pseudocode, `A[i, j]` denotes the entry of the matrix
$A$ in row $i$ and column
$j$ with the indices starting from 1. The
transformation is performed \textit{in place}, meaning that the original matrix
is lost for being eventually replaced by its row-echelon form.

`h := 1 /\textit{ `}`Initialization of the pivot row`\textit{` }/`  
`k := 1 /\textit{ `}`Initialization of the pivot column`\textit{` }/`  
  
\textbf{`while`}` h = m `\textbf{`and`}` k = n`  
`    /\textit{ `}`Find the k-th pivot:`\textit{` }/`  
`    i\_max := argmax (i = h ... m, abs(A[i, k]))`  
`    `\textbf{`if`}` A[i\_max, k] = 0`  
`        /\textit{ `}`No pivot in this column, pass to next column`\textit{` }/`  
`        k := k + 1`  
`    `\textbf{`else`}  
`         `\textbf{`swap rows`}`(h, i\_max)`  
`         /\textit{ `}`Do for all rows below pivot:`\textit{` }/`  
`         `\textbf{`for`}` i = h + 1 ... m:`  
`             f := A[i, k] / A[h, k]`  
`             /\textit{ `}`Fill with zeros the lower part of pivot column:`\textit{` }/`  
`             A[i, k] := 0`  
`             /\textit{ `}`Do for all remaining elements in current row:`\textit{` }/`  
`             `\textbf{`for`}` j = k + 1 ... n:`  
`                 A[i, j] := A[i, j] - A[h, j] * f`  
`         /\textit{ `}`Increase pivot row and column`\textit{` }/`  
`         h := h + 1`  
`         k := k + 1`

This algorithm differs slightly from the one discussed earlier, by
choosing a pivot with largest absolute value. Such a \textit{partial pivoting}
may be required if, at the pivot place, the entry of the matrix is zero.
In any case, choosing the largest possible absolute value of the pivot
improves the numerical stability of the algorithm, when floating point
is used for representing numbers.

Upon completion of this procedure the matrix will be in row echelon form
and the corresponding system may be solved by back substitution.

\paragraph{Resources [resources\_2]}

\begin{itemize}
    \item \href{https://openstax.org/books/college-algebra/pages/7-6-solving-systems-with-gaussian-elimination}{Solving Systems with Gaussian
\end{itemize}
    Elimination},
    OpenStax
\begin{itemize}
    \item \href{https://tutorial.math.lamar.edu/classes/alg/augmentedmatrix.aspx}{Augmented Matrices and Gauss-Jordan
\end{itemize}
    Elimination},
    Paul''''s Online Notes

\paragraph{Licensing [licensing\_2]}

Content obtained and/or adapted from:

\begin{itemize}
    \item \href{https://en.wikipedia.org/wiki/Gaussian\_elimination}{Gaussian elimination,
\end{itemize}
    Wikipedia} under
    a CC BY-SA license"
\paragraph{Resources [resources\_1]}

\begin{itemize}
    \item \href{https://www.unf.edu/\textasciitilde{}mzhan/chapter3.pdf}{Brief Introduction to Vectors and
\end{itemize}
    Matrices}, University of
    North Florida
\begin{itemize}
    \item \href{https://www.usna.edu/Users/math/uhan/sm286a/lessons/02.pdf}{Introduction to Matrices and
\end{itemize}
    Vectors},
    United States Naval Academy

\paragraph{Licensing [licensing\_1]}

Content obtained and/or adapted from:

\begin{itemize}
    \item \href{https://en.wikibooks.org/wiki/Linear\_Algebra/Matrices\_and\_Vectors/}{Matrices and Vectors, Wikibooks: Linear
\end{itemize}
    Algebra}
    under a CC BY-SA license"

%%===============================================================
\section{Solutions of Linear Systems}
\label{sec:solutions-of-linear-systems}
%%===============================================================

\paragraph{Solutions of linear systems [solutions\_of\_linear\_systems]}

A \textit{solution of a linear equation} is any n-tuple of values
$(s_1,s_2,....,s_n)$which satisfies the linear equation. For example,
$(-1,-1)$ is a solution of the linear equation $x + 3y = -4$ since
$-1 + (3 \times -1) = -1 + (-3) = -4$, but $(1,5)$ is not.

Similarly, a \textit{solution to a linear system} is any n-tuple of values
$(s_1,s_2,....,s_n)$which simultaneously satisfies all the linear
equations given in the system.

For example,

$$\begin{alignedat}{7}
3x && \; + \; &&             2y && \; - \; &&  z && \; = \; &&  1 & \\\\
2x && \; - \; &&             2y && \; + \; && 4z && \; = \; && -2 & \\\\
-x && \; + \; && \tfrac{1}{2} y && \; - \; &&  z && \; = \; &&  0 &
\end{alignedat}$$

has as its solution $(1,-2,-2)$. This can also be written as:

$\begin{alignedat}{2}
x \& = \& 1 \\\\
y \& = \& -2 \\\\
z \& = \& -2
\end{alignedat}$

We also refer to the collection of all possible solutions as the
\textit{solution set}.

In general, for any linear system of equations there are three
possibilities regarding solutions:

\textbf{A unique solution:} In this case only one specific solution set
exists. Geometrically this implies the n-planes specified by each
equation of the linear system all intersect at a unique point in the
space that is specified by the variables of the system.

\textbf{No solution:} The equations are termed inconsistent and specify
n-planes in space which do not intersect or overlap. It is not possible
to specify a solution set that satisfies all equations of the system.

\textbf{An infinite range of solutions:} The equations specify n-planes whose
intersection is an m-plane where $m\le n$. This being the case, it is
possible to show that an infinite set of solutions within a specific
range exists that satisfy the set of linear equations.

The following pictures illustrate these cases:

:   {\| border=0 cellpadding=5

\|- \|width=""150px"" aligned=""center""\|
 \|width=""150px""
aligned=""center""\|  \|width=""150px"" aligned=""center""\|
 \|- \|aligned=""center""\| A
unique solution \|aligned=""center""\| No solution \|aligned=""center""\| An
infinite range of solutions \|}

Why are there only these three cases and no others? Although a
justification shall be provided in the next chapter, it is a good
exercise for you to figure it out now.

A linear system is said to be \textit{inconsistent} if it has no solution. If
there exists at least one solution, then the system is said to be
\textit{consistent}.

\paragraph{Methods of solving linear systems [methods\_of\_solving\_linear\_systems]}

We know that linear equations in 2 or 3 variables can be solved using
techniques such as the elimination and the substitution method. However
these techniques are not appropriate for dealing with large systems
where there are a large number of variables. These techniques are
therefore generalized and a systematic procedure called Gaussian
elimination is usually used in actual practice. A variant of this
technique known as the Gauss Jordan method is also used.

Many times we are required to solve many linear systems where the only
difference in them are the constant terms. The coefficients of the
variables all remain the same. A technique called LU decomposition is
used in this case. A variant called Cholesky factorization is also used
when possible. We will study these techniques in later chapters.

\paragraph{Gauss Jordan Elimination [gauss\_jordan\_elimination]}

Performing row operations on a matrix is the method we use for solving a
system of equations. In order to solve the system of equations, we want
to convert the matrix to reduced row echelon form (""rref""), in which
there are ones down the main diagonal from the upper left corner to the
lower right corner, and zeros in every position below the main diagonal.
Here are some examples of matrices in rref:

$\begin{bmatrix}
1 \& 3 \& 1\\\\
0 \& 1 \& -1\\\\
0 \& 0 \& 1
\end{bmatrix}$, $\begin{bmatrix}
1 \& 3 \& 1\\\\
0 \& 1 \& -1\\\\
\end{bmatrix}$, $\begin{bmatrix}
1 \& 0 \& 0 \& 3\\\\
0 \& 1 \& 0 \& 0\\\\
0 \& 0 \& 1 \& 2\\\\
\end{bmatrix}$, $\begin{bmatrix}
1 \& 7 \\\\
0 \& 1 \\\\
\end{bmatrix}$, $\begin{bmatrix}
1 \& 3 \& 0\\\\
0 \& 1 \& 1\\\\
0 \& 0 \& 0
\end{bmatrix}$, $\begin{bmatrix}
1 \& 0 \& 0\\\\
0 \& 0 \& 0\\\\
0 \& 0 \& 0
\end{bmatrix}$

We use row operations corresponding to equation operations to obtain a
new matrix that is row-equivalent in a simpler form. Here are the
guidelines to obtaining row-echelon form.

\begin{itemize}
    \item In any nonzero row, the first nonzero number is a 1. It is called a
\end{itemize}
    leading 1.
\begin{itemize}
    \item Any all-zero rows are placed at the bottom on the matrix.
    \item Any leading 1 is below and to the right of a previous leading 1.
    \item Any column containing a leading 1 has zeros in all other positions
\end{itemize}
    in the column.

To solve a system of equations we can perform the following row
operations to convert the coefficient matrix to row-echelon form and do
back-substitution to find the solution.

\begin{itemize}
    \item Interchange a row $R_{i}$ with another row $R_{j}$. (Notation:
\end{itemize}
    $R_{i} \longleftrightarrow R_{j}$ )
\begin{itemize}
    \item Multiply a row by a constant. (Notation: $cR_{i}$)
    \item Add the product of a row multiplied by a constant to another row.
\end{itemize}
    (Notation: $R_{i} + cR_{j}$)

Each of the row operations corresponds to the operations we have already
learned to solve systems of equations in three variables. With these
operations, there are some key moves that will quickly achieve the goal
of writing a matrix in row-echelon form. To obtain a matrix in
row-echelon form for finding solutions, we use Gaussian (also called
Gauss-Jordan elimination, a method that uses row operations to obtain a
1 as the first entry so that row 1 can be used to convert the remaining
rows.

The steps to perform Gaussian elimination on an augmented matrix (and
thus convert the matrix to reduced row echelon form) are as follows.

1.  The first equation should have a leading coefficient of 1.
    Interchange rows or multiply by a constant, if necessary.
2.  Use row operations to obtain zeros down the first column below the
    first entry of 1.
3.  Use row operations to obtain a 1 in row 2, column 2.
4.  Use row operations to obtain zeros down column 2, below the entry of
    1.
5.  Use row operations to obtain a 1 in row 3, column 3.
6.  Continue this process for all rows until there is a 1 in every entry
    down the main diagonal and there are only zeros below.
7.  If any rows contain all zeros, place them at the bottom.

Example of Gauss-Jordan elimination:

$\begin{bmatrix}
1 \& 3 \& 1 \& 9 \\\\
1 \& 1 \& -1 \& 1 \\\\
3 \& 11 \& 5 \& 35
\end{bmatrix}\to
\begin{bmatrix}
1 \& 3 \& 1 \& 9 \\\\
0 \& -2 \& -2 \& -8 \\\\
0 \& 2 \& 2 \& 8
\end{bmatrix}\to
\begin{bmatrix}
1 \& 3 \& 1 \& 9 \\\\
0 \& -2 \& -2 \& -8 \\\\
0 \& 0 \& 0 \& 0
\end{bmatrix}\to
\begin{bmatrix}
1 \& 0 \& -2 \& -3 \\\\
0 \& 1 \& 1 \& 4 \\\\
0 \& 0 \& 0 \& 0
\end{bmatrix}$

First, we subtract 1 of the first row from the second row
($R_{2} - R_{1}$) and 3 of the first row from the third row
($R_{3} - 3R_{1}$). Then, we add the new second row to the new third row
($R_{3} + R_{2}$). Lastly, we multiply the new second row by
$-\tfrac{1}{2}$ to make the diagonal entries equal to 1
($-\tfrac{1}{2}R_{2})$. Our matrix is now in reduced row echelon form,
since all nonzero diagonal entries are equal to 1, the all-zero row is
at the bottom of the matrix, and all entries under the diagonal are
zero.

%%===============================================================
\section{Linear Transformations}
\label{sec:linear-transformations}
%%===============================================================

A \textbf{linear transformation} is an important concept in mathematics
because many real world phenomena can be approximated by linear models.

Unlike a linear function, a linear transformation works on vectors as
well as numbers.

\paragraph{Motivations and definitions [motivations\_and\_definitions]}

Say we have the vector $\begin{pmatrix} 1 \\\\ 0 \end{pmatrix}$ in
$\mathbb{R}^2$, and we rotate it through 90 degrees, to obtain the
vector $\begin{pmatrix} 0 \\\\ 1 \end{pmatrix}$.

Another example instead of rotating a vector, we stretch it, so a vector
$\mathbf{v}$ becomes $2\mathbf{v}$, for example.
$\begin{pmatrix} 2 \\\\ 3 \end{pmatrix}$ becomes
$\begin{pmatrix} 4 \\\\ 6 \end{pmatrix}$

Or, if we look at the \textit{projection} of one vector onto the \textit{x} axis -
extracting its \textit{x} component - , e.g. from
$\begin{pmatrix} 2 \\\\ 3 \end{pmatrix}$ we get
$\begin{pmatrix} 2 \\\\ 0 \end{pmatrix}$

These examples are all an example of a \textit{mapping} between two vectors,
and are all linear transformations. If the rule transforming the matrix
is called $T$, we often write $T\mathbf{v}$ for the mapping of the
vector $\mathbf{v}$ by the rule $T$. $T$ is often called the
transformation.

Note we do not always write brackets like when we write functions.
However we \textit{should} write brackets, especially when we want to express
the mapping of the sum or the product or the combination of many
vectors.

\paragraph{Definitions [definitions\_2]}

\paragraph{Linear Operators [linear\_operators]}

Suppose one has a field K, and let x be an element of that field. Let O
be a function taking values from K where O(x) is an element of a field
J. Define O to be a linear form if and only if:

1.  O(x+y)=O(x)+O(y)
2.  O($\lambda$x)=$\lambda$O(x)

\paragraph{Linear Forms [linear\_forms]}

Suppose one has a vector space V, and let x be an element of that vector
space. Let F be a function taking values from V where F(x) is an element
of a field K. Define F to be a linear form if and only if:

1.  F(x+y)=F(x)+F(y)
2.  F($\lambda$x)=$\lambda$F(x)

\paragraph{Linear Transformation [linear\_transformation]}

This time, instead of a field, let us consider functions from one vector
space into another vector space. Let T be a function taking values from
one vector space V where L(V) are elements of another vector space.
Define L to be a linear transformation when it:

1.  \textit{preserves scalar multiplication}: T($\lambda$\textbf{x}) = $\lambda$T\textbf{x}
2.  \textit{preserves addition}: T(\textbf{x}+\textbf{y}) = T\textbf{x} + T\textbf{y}

Note that not all transformations are linear. Many simple
transformations that are in the real world are also non-linear. Their
study is more difficult, and will not be done here. For example, the
transformation \textit{S} (whose input and output are both vectors in \textbf{$R^2$})
defined by

$S\mathbf{x} = S\begin{pmatrix} x \\\                                     y  \end{pmatrix}= \begin{pmatrix}  xy\\\                                                                        \cos(y)\end{pmatrix}$

We can learn about nonlinear transformations by studying easier, linear
ones.

We often \textit{describe} a transformation T in the following way

$$T : V \rightarrow W$$

This means that T, whatever transformation it may be, maps vectors in
the vector space V to a vector in the vector space W.

The actual transformation \textit{could} be written, for instance, as

$$T\begin{pmatrix} x \\\\ y\end{pmatrix} = \begin{pmatrix} x + y \\\\ x - y \end{pmatrix}$$

\paragraph{Examples and proofs [examples\_and\_proofs]}

Here are some examples of some linear transformations. At the same time,
let''''s look at how we can prove that a transformation we may find is
linear or not.

\paragraph{Projection}

Let us take the projection of vectors in \textbf{$R^2$} to vectors on the
\textit{x}-axis. Let''''s call this transformation T.

We know that T maps vectors from \textbf{$R^2$} to \textbf{$R^2$}, so we can say

$$T: \mathbb{R}^2 \rightarrow \mathbb{R}^2$$

and we can then write the transformation itself as

$$T\begin{pmatrix} x_0 \\\\ x_1 \end{pmatrix} = \begin{pmatrix} x_0 \\\\ 0 \end{pmatrix}$$

Clearly this is linear. (\textit{Can you see why, without looking below?})

Let''''s go through a proof that the conditions in the definitions are
established.

\paragraph{Scalar multiplication is preserved [scalar\_multiplication\_is\_preserved]}

We wish to show that for all vectors \textbf{v} and all scalars $\lambda$,
T($\lambda$\textbf{v})=$\lambda$T(\textbf{v}).

Let

$$\mathbf{v}=\begin{pmatrix} v_0 \\\\ v_1 \end{pmatrix}$$

Then

$$\lambda\mathbf{v}=\begin{pmatrix} \lambda v_0 \\\\ \lambda v_1 \end{pmatrix}$$
Now

$$T(\lambda\mathbf{v}) = T\begin{pmatrix} \lambda v_0 \\\\ \lambda v_1\end{pmatrix} =$$

$$\begin{pmatrix} \lambda v_0 \\\\ 0 \end{pmatrix}$$ If we work out
$\lambda$T(\textbf{v}) and find it is the same vector, we have proved our result.

$$\lambda T\mathbf{v}= \lambda \begin{pmatrix} v_0 \\\\ 0 \end{pmatrix}=$$

$$\begin{pmatrix} \lambda v_0 \\\\ 0 \end{pmatrix}$$ This is the same
vector as above, so under the transformation T, *scalar multiplication
is preserved*.

\paragraph{Addition is preserved [addition\_is\_preserved]}

We wish to show for all vectors \textbf{x} and \textbf{y},
T(\textbf{x}+\textbf{y})=T\textbf{x}+T\textbf{y}.

Let

$$\mathbf{x}=\begin{pmatrix} x_0 \\\\ x_1 \end{pmatrix}$$

and

$$\mathbf{y}=\begin{pmatrix} y_0 \\\\ y_1 \end{pmatrix}$$

Now

$$T(\mathbf{x}+\mathbf{y})=T\left(\begin{pmatrix} x_0 \\\\ x_1 \end{pmatrix}+\begin{pmatrix} y_0 \\\\ y_1 \end{pmatrix}\right)=$$
$$T\begin{pmatrix} x_0 + y_0 \\\\ x_1 + y_1 \end{pmatrix} =$$
$$\begin{pmatrix} x_0 + y_0 \\\\ 0 \end{pmatrix}$$

Now if we can show T\textbf{x}+T\textbf{y} is this vector above, we have proved
this result. Proceed, then,

$$T\begin{pmatrix} x_0 \\\\ x_1 \end{pmatrix} + T\begin{pmatrix} y_0 \\\\ y_1 \end{pmatrix}=\begin{pmatrix} x_0 \\\\ 0 \end{pmatrix} + \begin{pmatrix} y_0 \\\\0 \end{pmatrix}=$$

$$\begin{pmatrix} x_0 + y_0 \\\\ 0 \end{pmatrix}$$

So we have that the transformation T \textit{preserves addition}.

\paragraph{Zero vector is preserved [zero\_vector\_is\_preserved]}

Clearly we have

$$T\begin{pmatrix} 0 \\\\ 0 \end{pmatrix} = \begin{pmatrix} 0 \\\\ 0 \end{pmatrix}$$

\paragraph{Conclusion}

We have shown T preserves addition, scalar multiplication and the zero
vector. So T must be linear.

\paragraph{Disproof of linearity [disproof\_of\_linearity]}

When we want to \textit{disprove} linearity - that is, to \textit{prove} that a
transformation is \textit{not} linear, we need only find one counter-example.

If we can find just one case in which the transformation does not
preserve addition, scalar multiplication, or the zero vector, we can
conclude that the transformation is not linear.

For example, consider the transformation

$$T\begin{pmatrix} x \\\\ y\end{pmatrix} = \begin{pmatrix}  x^3 \\\\ y^2\end{pmatrix}$$

We suspect it is not linear. To prove it is not linear, take the vector

$$\mathbf{v} = \begin{pmatrix} 2 \\\\ 2 \end{pmatrix}$$

then

$$T(2\mathbf{v}) = \begin{pmatrix} 64 \\\\ 16 \end{pmatrix}$$

but

$$2T(\mathbf{v}) = \begin{pmatrix} 16 \\\\ 8 \end{pmatrix}$$

so we can immediately say T is not linear because it doesn''''t preserve
scalar multiplication.

\paragraph{Problem set [problem\_set]}

Given the above, determine whether the following transformations are in
fact linear or not. Write down each transformation in the form T:V -\>
W, and identify V and W. (Answers follow to even-numbered questions):

1.  $T\begin{pmatrix} v_0 \\\\ v_1 \end{pmatrix} = \begin{pmatrix} v_0^2 + v_1 \\\\ v_1 \end{pmatrix}$
2.  $T\begin{pmatrix} v_0 \\\\ v_1 \end{pmatrix} = \begin{pmatrix} 1 \\\\ v_0 \end{pmatrix}$
3.  $T\begin{pmatrix} v_0 \\\\ v_1 \end{pmatrix} = \mathbf{0}$
4.  $T\begin{pmatrix} v_0 \\\\ v_1 \\\\ v_2 \end{pmatrix} = \begin{pmatrix} v_0 - v_2 \\\\ v_1 \end{pmatrix}$

\paragraph{Answers}

$\quad$2\. No. A check whether the zero vector is preserved readily
    confirms this fact. T : \textbf{$R^2$} -\> \textbf{$R^2$}

$\quad$4\. Yes. T : \textbf{$R^3$} -\> \textbf{$R^2$}.

\paragraph{Images and kernels [images\_and\_kernels]}

We have some fundamental concepts underlying linear transformations,
such as the \textit{kernel} and the \textit{image} of a linear transformation, which
are analogous to the \textit{zeros} and \textit{range} of a function.

\paragraph{Kernel}

The \textit{kernel} of a linear transformation T: V -\> W is the set of all
vectors in V which are mapped to the zero vector in W, ie.,

$$\mathrm{ker}\ T = \{v \in V\ |\ T\mathbf{v} = \mathbf{0}\}$$

Coincidentally because of the matr to the matrix equation A\textbf{x}=\textbf{0}.

The kernel of a transform T: V-\>W is always a subspace of V. The
dimension of a transform or a matrix is called the \textit{nullity}..

\paragraph{Image}

The \textit{image} of a linear transformation T:V-\>W is the set of all vectors
in W which were mapped from vectors in V. For example with the trivial
mapping T:V-\>W such that T\textbf{x}=\textbf{0}, the image would be \textbf{0}. (*What
would the kernel be?*).

More formally, we say that the image of a transformation T:V-\>W is the
set

$$\mathrm{im}\ T = \{w \in W\ |\ w=T\mathbf{v}\ \mathrm{and}\ \mathbf{v}\in V\}$$

\paragraph{Isomorphism}

A linear transformation T:V -\> W is an isomorphic transformation if it
is:

\begin{itemize}
    \item one-to-one and onto.
    \item kernel(T) = {0} and the range(T) = W.
    \item an inverse of T exists.
    \item dim(V) = dim(W).
\end{itemize}

\paragraph{Licensing [licensing\_5]}

Content obtained and/or adapted from:

\begin{itemize}
    \item \href{https://en.wikibooks.org/wiki/Linear\_Algebra/Linear\_Transformations}{Linear Transformations, Wikibooks: Linear
\end{itemize}
    Algebra}
    under a CC BY-SA license"

%%===============================================================
\section{Matrix Products and Inverses}
\label{sec:matrix-products-and-inverses}
%%===============================================================

\paragraph{Matrix Operations [matrix\_operations]}

\paragraph{Adding and subtracting matrices [adding\_and\_subtracting\_matrices]}

In order to add or subtract two matrices, they must be of the same
dimension; that is, the two matrices must have the same number of rows
and the same number of columns. To add two matrices together, we simply
need to add every entry in one matrix to the entry in the same row and
same column in the other matrix. For example:

$$\begin{bmatrix}
a_{11} & a_{12} & a_{13}\\\\
a_{21} & a_{22} & a_{23}\\\\
a_{31} & a_{32} & a_{33}
\end{bmatrix} + 
\begin{bmatrix}
b_{11} & b_{12} & b_{13}\\\\
b_{21} & b_{22} & b_{23}\\\\
b_{31} & b_{32} & b_{33}
\end{bmatrix} = 
\begin{bmatrix}
a_{11} + b_{11} & a_{12} + b_{12} & a_{13} + b_{13}\\\\
a_{21} + b_{21} & a_{22} + b_{22} & a_{23} + b_{23}\\\\
a_{31} + b_{31} & a_{32} + b_{32} & a_{33} + b_{33}
\end{bmatrix}$$

$$\begin{bmatrix}
1 & 2 & 3\\\\
0 & 1 & 2\\\\
0 & 0 & 1
\end{bmatrix} + 
\begin{bmatrix}
5 & 2 & 1\\\\
3 & 3 & 1\\\\
4 & 2 & 0
\end{bmatrix} = 
\begin{bmatrix}
1 + 5 & 2 + 2 & 3 + 1\\\\
0 + 3 & 1 + 3 & 2 + 1\\\\
0 + 4 & 0 + 2 & 1 + 0
\end{bmatrix} = 
\begin{bmatrix}
6 & 4 & 4\\\\
3 & 4 & 3\\\\
4 & 2 & 1
\end{bmatrix}$$

\paragraph{Multiplying matrices by scalars [multiplying\_matrices\_by\_scalars]}

When multiplying a matrix by a scalar (or number), all we need to do is
multiply each entry of the matrix by the scalar. For example:

$$3\begin{bmatrix}
1 & 2 & 3\\\\
0 & 1 & 2\\\\
0 & 0 & 1
\end{bmatrix} = 
\begin{bmatrix}
3 & 6 & 9\\\\
0 & 3 & 6\\\\
0 & 0 & 3
\end{bmatrix}$$

$$-\frac{1}{2}\begin{bmatrix}
2 & -4 & 8\\\\
-2 & 1 & -6\\\\
0 & 0 & 7
\end{bmatrix} = 
\begin{bmatrix}
-1 & 2 & -4\\\\
1 & -\frac{1}{2} & 3\\\\
0 & 0 & -\frac{7}{2}
\end{bmatrix}$$

\paragraph{Multiplying matrices [multiplying\_matrices]}

Matrix multiplication is not commutative; that is, for most matrices $A$
and $B$, $AB \neq BA$. Two matrices can only be multiplied together if
the number of columns in the first matrix equals the number of rows in
the second matrix. For example, we can take the product of a 3-by-2
matrix times a 2-by-5 matrix, but not a 3-by-2 matrix times a 3-by-2
matrix. The product of two matrices has the same number of rows as the
first matrix and the same number of columns as the second matrix. For
example, a 2-by-3 matrix times a 3-by-2 matrix will result in a 2-by-2
matrix.

The product of two 3-by-3 matrices is $\begin{bmatrix}
a_{11} \& a_{12} \& a_{13}\\\\
a_{21} \& a_{22} \& a_{23}\\\\
a_{31} \& a_{32} \& a_{33}
\end{bmatrix} 
\begin{bmatrix}
b_{11} \& b_{12} \& b_{13}\\\\
b_{21} \& b_{22} \& b_{23}\\\\
b_{31} \& b_{32} \& b_{33}
\end{bmatrix} = 
\begin{bmatrix}
c_{11} \& c_{12} \& c_{13}\\\\
c_{21} \& c_{22} \& c_{23}\\\\
c_{31} \& c_{32} \& c_{33}
\end{bmatrix}$

where $c_{ij}$ is the dot product of the i-th row of the first matrix
and the j-th column of the second matrix; that is,
$c_{ij} = a_{i1}b_{1j} + a_{i2}b_{2j} + a_{i3}b_{3j}$.

Here is a 2-by-2 example of matrix multiplication:

$$\begin{bmatrix}
1 & 2\\\\
3 & 4
\end{bmatrix} 
\begin{bmatrix}
5 & 6\\\\
7 & 8
\end{bmatrix} = 
\begin{bmatrix}
1(5) + 2(7) & 1(6) + 2(8)\\\\
3(5) + 4(7) & 3(6) + 4(8)
\end{bmatrix} = 
\begin{bmatrix}
5 + 14 & 6 + 16\\\\
15 + 28 & 18 + 32
\end{bmatrix} \begin{bmatrix}
19 & 22\\\\
43 & 50
\end{bmatrix}$$

Note that there is no ""matrix division"". Rather, we multiply one matrix
by the inverse of another (for example, $AB^{-1}$, NOT $A/B$) to obtain
a ""quotient"" from two matrices.

\paragraph{Inverse of an n-by-n matrix [inverse\_of\_an\_n\_by\_n\_matrix]}

An n-by-n matrix A is the inverse of n-by-n matrix B (and B the inverse
of A) if BA = AB = I, where I is an identity matrix.

The inverse of an n-by-n matrix can be calculated by creating an n-by-2n
matrix which has the original matrix on the left and the identity matrix
on the right. Row reduce this matrix and the right half will be the
inverse. If the matrix does not row reduce completely (i.e., a row is
formed with all zeroes as its entries), it does not have an inverse.

\paragraph{Example 1 [example\_1]}

Let $\mathrm{A} = \begin{bmatrix}
  1 \& 4 \& 4\\\\
    2 \& 5 \& 8\\\\ 
  3 \& 6 \& 9
\end{bmatrix}$

We begin by expanding and partitioning A to include the identity matrix,
and then proceed to row reduce A until we reach the identity matrix on
the left-hand side.

$$\begin{bmatrix}
  1 & 4 & 4 & \big| & 1 & 0 & 0\\\\
    2 & 5 & 8 & \big| & 0 & 1 & 0\\\\ 
  3 & 6 & 9 & \big| & 0 & 0 & 1
\end{bmatrix} \rightsquigarrow \begin{bmatrix}   1 & 4 & 4 & \big| & 1 & 0 & 0\\\\
    0 & -3 & 0 & \big| & -2 & 1 & 0\\\\ 
  0 & -6 & -3 & \big| & -3 & 0 & 1
\end{bmatrix} \rightsquigarrow \begin{bmatrix}   1 & 4 & 4 & \big| & 1 & 0 & 0\\\\
    0 & 1 & 0 & \big| & 2/3 & -1/3 & 0\\\\ 
  0 & -6 & -3 & \big| & -3 & 0 & 1
\end{bmatrix} \rightsquigarrow$$

$$\begin{bmatrix}
  1 & 0 & 4 & \big| & -5/3 & 4/3 & 0\\\\
    0 & 1 & 0 & \big| & 2/3 & -1/3 & 0\\\\ 
  0 & 0 & -3 & \big| & 1 & -2 & 1
\end{bmatrix} \rightsquigarrow \begin{bmatrix}   1 & 0 & 4 & \big| & -5/3 & 4/3 & 0\\\\
    0 & 1 & 0 & \big| & 2/3 & -1/3 & 0\\\\ 
  0 & 0 & 1 & \big| & -1/3 & 2/3 & -1/3
\end{bmatrix} \rightsquigarrow \begin{bmatrix}   1 & 0 & 0 & \big| & -1/3 & -4/3 & 4/3\\\\
    0 & 1 & 0 & \big| & 2/3 & -1/3 & 0\\\\ 
  0 & 0 & 1 & \big| & -1/3 & 2/3 & -1/3
\end{bmatrix}$$  
  
The matrix $\mathrm{B} = \begin{bmatrix}
  -1/3 \& -4/3 \& 4/3\\\\
    2/3 \& -1/3 \& 0\\\\ 
  -1/3 \& 2/3 \& -1/3
\end{bmatrix}$  is then the inverse of the original matrix A.

\paragraph{Inverse of a Linear Transformation [inverse\_of\_a\_linear\_transformation]}

We now consider how to represent the inverse of a linear map. We start
by recalling some facts about function inverses. Some functions have no
inverse, or have an inverse on the left side or right side only.

\paragraph{Definitions}

Where $\pi:\mathbb{R}^3\to \mathbb{R}^2$ is the projection map

$$\begin{pmatrix} x \\\\ y \\\\ z \end{pmatrix} \mapsto
\begin{pmatrix} x \\\\ y \end{pmatrix}$$

and $\eta:\mathbb{R}^2\to \mathbb{R}^3$ is the embedding

$$\begin{pmatrix} x \\\\ y \end{pmatrix} \mapsto
\begin{pmatrix} x \\\\ y \\\\ 0 \end{pmatrix}$$

the composition $\pi\circ \eta$ is the identity map on $\mathbb{R}^2$.

$$\begin{pmatrix} x \\\\ y \end{pmatrix} \stackrel{\eta}{\longmapsto}
\begin{pmatrix} x \\\\ y \\\\ 0 \end{pmatrix} \stackrel{\pi}{\longmapsto}
\begin{pmatrix} x \\\\ y \end{pmatrix}$$

We say $\pi$ is a \textbf{left inverse map} of $\eta$ or, what is the same
thing, that $\eta$ is a \textbf{right inverse map} of $\pi$. However,
composition in the other order $\eta\circ \pi$ doesn''''t give the identity
map— here is a vector that is not sent to itself under $\eta\circ \pi$.

$$\begin{pmatrix} 0 \\\\ 0 \\\\ 1 \end{pmatrix} \stackrel{\pi}{\longmapsto}
\begin{pmatrix} 0 \\\\ 0 \end{pmatrix} \stackrel{\eta}{\longmapsto}
\begin{pmatrix} 0 \\\\ 0 \\\\ 0 \end{pmatrix}$$

In fact, the projection $\pi$ has no left inverse at all. For, if $f$$
were to be a left inverse of $\pi$ then we would have

$$\begin{pmatrix} x \\\\ y \\\\ z \end{pmatrix} \stackrel{\pi}{\longmapsto}
\begin{pmatrix} x \\\\ y \end{pmatrix} \stackrel{f}{\longmapsto}
\begin{pmatrix} x \\\\ y \\\\ z \end{pmatrix}$$

for all of the infinitely many $z$''''s. But no function $f$ can send a
single argument to more than one value.

(An example of a function with no inverse on either side is the zero
transformation on $\mathbb{R}^2$.)

Some functions have a **two-sided inverse map**, another function that
is the inverse of the first, both from the left and from the right. For
instance, the map given by $\vec{v}\mapsto 2\cdot \vec{v}$ has the
two-sided inverse $\vec{v}\mapsto (1/2)\cdot\vec{v}$. In this subsection
we will focus on two-sided inverses. The appendix shows that a function
has a two-sided inverse if and only if it is both one-to-one and onto.
The appendix also shows that if a function $f$ has a two-sided inverse
then it is unique, and so it is called ""the"" inverse, and is denoted
$f^{-1}$. So our purpose in this subsection is, where a linear map $h$
has an inverse, to find the relationship between ${\rm Rep}_{B,D}(h)$
and ${\rm Rep}_{D,B}(h^{-1})$.

A matrix $G$ is a **left inverse matrix** of the matrix $H$ if $GH$ is
the identity matrix. It is a **right inverse matrix** if $HG$ is the
identity. A matrix $H$ with a two-sided inverse is an **invertible
matrix**. That two-sided inverse is called **the inverse matrix** and is
denoted $H^{-1}$.

Because of the correspondence between linear maps and matrices,
statements about map inverses translate into statements about matrix
inverses.

Lemmas and Theorems

-   If a matrix has both a left inverse and a right inverse then the two
    are equal.

<!-- -->

-   A matrix is invertible if and only if it is nonsingular.
    -   Proof: *(For both results.)* Given a matrix $H$, fix spaces of
        appropriate dimension for the domain and codomain. Fix bases for
        these spaces. With respect to these bases, $H$ represents a map
        $h$. The statements are true about the map and therefore they
        are true about the matrix.

<!-- -->

-   A product of invertible matrices is invertible— if $G$ and $H$ are
    invertible and if $GH$ is defined then $GH$ is invertible and
    $(GH)^{-1}=H^{-1}G^{-1}$.
    -   Proof: *(This is just like the prior proof except that it
        requires two maps.)* Fix appropriate spaces and bases and
        consider the represented maps $h$ and $g$. Note that
        $h^{-1}g^{-1}$ is a two-sided map inverse of $gh$ since
        $(h^{-1}g^{-1})(gh) = h^{-1}(\textrm{id})h = h^{-1}h =\textrm{id}$
        and
        $(gh)(h^{-1}g^{-1}) = g(\textrm{id})g^{-1} = gg^{-1} = \textrm{id}$.
        This equality is reflected in the matrices representing the
        maps, as required.

Here is the arrow diagram giving the relationship between map inverses
and matrix inverses. It is a special case of the diagram for function
composition and matrix multiplication.

Linalg\_matinv\_arrow.png

Beyond its place in our general program of seeing how to represent map
operations, another reason for our interest in inverses comes from
solving linear systems. A linear system is equivalent to a matrix
equation, as here.

\begin{align}
x_1  \& +  \& x_2  \& =  \& 3  \\\\
2x_1  \& -  \& x_2  \& =  \& 2
\end{aligned} \quad \Longleftrightarrow \quad
\begin{pmatrix}
1  \& 1  \\\\
2  \& -1
\end{pmatrix} \begin{pmatrix} x_1 \\\\ x_2 \end{pmatrix} = \begin{pmatrix} 3 \\\\ 2 \end{pmatrix} \qquad\qquad (*)$$

By fixing spaces and bases (e.g., $\mathbb{R}^2, \mathbb{R}^2$ and
$\mathcal{E}_2,\mathcal{E}_2$), we take the matrix $H$ to represent some
map $h$. Then solving the system is the same as asking: what domain
vector $\vec{x}$ is mapped by $h$ to the result $\vec{d}\,$? If we could
invert $h$ then we could solve the system by multiplying
${\rm Rep}_{D,B}(h^{-1})\cdot{\rm Rep}_{D}(\vec{d})$ to get
${\rm Rep}_{B}(\vec{x})$.

Example 2 [example_2]

We can find a left inverse for the matrix just given

$$\begin{pmatrix}
m  \& n  \\\\
p  \& q
\end{pmatrix} \begin{pmatrix} 1  \& 1  \\\\
2  \& -1
\end{pmatrix} = \begin{pmatrix} 1  \& 0  \\\\
0  \& 1
\end{pmatrix}$$

by using Gauss'''' method to solve the resulting linear system.

$$\begin{array}{r}
m  \& +  \& 2n  \&    \&   \&   \&    \& =  \& 1     \\\\
m  \& -  \& n   \&    \&   \&   \&    \& =  \& 0     \\\\
\&   \&    \&    \& p  \& +  \& 2q  \& =  \& 0     \\\\
\&   \&    \&    \& p  \& -  \& q   \& =  \& 1
\end{array}$$

Answer: $m=1/3$, $n=1/3$, $p=2/3$, and $q=-1/3$. This matrix is actually
the two-sided inverse of $H$, as can easily be checked. With it we can
solve the system ($*$) above by applying the inverse.

$$\begin{pmatrix} x \\\\ y \end{pmatrix} =\begin{pmatrix}
1/3  \& 1/3  \\\\
2/3  \& -1/3
\end{pmatrix} \begin{pmatrix} 3 \\\\ 2 \end{pmatrix} = \begin{pmatrix} 5/3 \\\\ 4/3 \end{pmatrix}$$

Remark

Why solve systems this way, when Gauss'''' method takes less arithmetic
(this assertion can be made precise by counting the number of arithmetic
operations, as computer algorithm designers do)? Beyond its conceptual
appeal of fitting into our program of discovering how to represent the
various map operations, solving linear systems by using the matrix
inverse has at least two advantages.

First, once the work of finding an inverse has been done, solving a
system with the same coefficients but different constants is easy and
fast: if we change the entries on the right of the system ($*$) then we
get a related problem

$$\begin{pmatrix}
1  \& 1  \\\\
2  \& -1
\end{pmatrix} \begin{pmatrix} x \\\\ y \end{pmatrix} =
\begin{pmatrix} 5 \\\\ 1 \end{pmatrix}$$

with a related solution method.

$$\begin{pmatrix} x \\\\ y \end{pmatrix} =
\begin{pmatrix}
1/3  \& 1/3  \\\\
2/3  \& -1/3
\end{pmatrix} \begin{pmatrix} 5 \\\\ 1 \end{pmatrix} =
\begin{pmatrix} 2 \\\\ 3 \end{pmatrix}$$

In applications, solving many systems having the same matrix of
coefficients is common.

Another advantage of inverses is that we can explore a system''''s
sensitivity to changes in the constants. For example, tweaking the $3$
on the right of the system ($*$) to

$$\begin{pmatrix}
1  \& 1  \\\\
2  \& -1
\end{pmatrix} \begin{pmatrix} x_1 \\\\ x_2 \end{pmatrix} =
\begin{pmatrix} 3.01 \\\\ 2 \end{pmatrix}$$

can be solved with the inverse.

$$\begin{pmatrix}
1/3  \& 1/3  \\\\
2/3  \& -1/3
\end{pmatrix} \begin{pmatrix} 3.01 \\\\ 2 \end{pmatrix} =
\begin{pmatrix} (1/3)(3.01)+(1/3)(2) \\\\ (2/3)(3.01)-(1/3)(2) \end{pmatrix}$$

to show that $x_1$ changes by $1/3$ of the tweak while $x_2$ moves by
$2/3$ of that tweak. This sort of analysis is used, for example, to
decide how accurately data must be specified in a linear model to ensure
that the solution has a desired accuracy. }}

We finish by describing the computational procedure usually used to find
the inverse matrix.

-   A matrix is invertible if and only if it can be written as the
    product of elementary reduction matrices. The inverse can be
    computed by applying to the identity matrix the same row steps, in
    the same order, as are used to Gauss-Jordan reduce the invertible
    matrix.
    -   Proof: A matrix $H$ is invertible if and only if it is
        nonsingular and thus Gauss-Jordan reduces to the identity. This
        reduction can be done with elementary matrices

        $$R_{r}\cdot R_{r-1}\cdot\dots R_{1}\cdot H=I$$
        This equation gives the two halves of the result.

        First, elementary matrices are invertible and their inverses are
        also elementary. Applying $R_r^{-1}$ to the left of both sides
        of that equation, then $R_{r-1}^{-1}$, etc., gives $H$ as the
        product of elementary matrices
        $H=R_1^{-1}\cdots R_r^{-1}\cdot I$ (the $I$ is here to cover the
        trivial $r=0$ case).

        Second, matrix inverses are unique and so comparison of the
        above equation with $H^{-1}H=I$ shows that
        $H^{-1}=R_r\cdot R_{r-1}\dots R_1\cdot I$. Therefore, applying
        $R_1$ to the identity, followed by $R_2$, etc., yields the
        inverse of $H$.

Example 3 [example_3]

To find the inverse of

$$\begin{pmatrix}
1  \& 1  \\\\
2  \& -1
\end{pmatrix}$$

we do Gauss-Jordan reduction, meanwhile performing the same operations
on the identity. For clerical convenience we write the matrix and the
identity side-by-side, and do the reduction steps together.

$$\begin{array}{rcl}
\left(\begin{array}{cc|cc}
1  \& 1   \& 1  \& 0  \\\\
2  \& -1  \& 0  \& 1
\end{array}\right)
\& \xrightarrow[]{-2\rho_1+\rho_2}
\& \left(\begin{array}{cc|cc}
1  \& 1   \& 1  \& 0  \\\\
0  \& -3  \& -2 \& 1
\end{array}\right)                 \\\\
\& \xrightarrow[]{-1/3\rho_2}
\& \left(\begin{array}{cc|cc}
1  \& 1   \& 1   \& 0     \\\\
0  \& 1   \& 2/3 \& -1/3
\end{array}\right)                   \\\\
\& \xrightarrow[]{-\rho_2+\rho_1}
\& \left(\begin{array}{cc|cc}
1  \& 0   \& 1/3 \& 1/3  \\\\
0  \& 1   \& 2/3 \& -1/3
\end{array}\right)
\end{array}$$ This calculation has found the inverse.

$$\begin{pmatrix}
1  \& 1  \\\\
2  \& -1
\end{pmatrix}^{-1} = \begin{pmatrix} 1/3 \& 1/3  \\\\
2/3 \& -1/3
\end{pmatrix}$$

Example 4 [example_4]

This one happens to start with a row swap.

$$\begin{array}{rcl}
\left(\begin{array}{ccc|ccc}
0  \& 3  \& -1  \& 1  \& 0  \& 0  \\\\
1  \& 0  \& 1   \& 0  \& 1  \& 0  \\\\
1  \& -1 \& 0   \& 0  \& 0  \& 1
\end{array}\right)
\& \xrightarrow[]{\rho_1\leftrightarrow\rho_2}
\&   \left(\begin{array}{ccc|ccc}
1  \& 0  \& 1   \& 0  \& 1  \& 0  \\\\
0  \& 3  \& -1  \& 1  \& 0  \& 0  \\\\
1  \& -1 \& 0   \& 0  \& 0  \& 1
\end{array}\right)                             \\\\
\& \xrightarrow[]{-\rho_1+\rho_3}
\&   \left(\begin{array}{ccc|ccc}
1  \& 0  \& 1   \& 0  \& 1  \& 0  \\\\
0  \& 3  \& -1  \& 1  \& 0  \& 0  \\\\
0  \& -1 \& -1  \& 0  \& -1 \& 1
\end{array}\right)                             \\\\
\& \vdots                                  \\\\
\& \xrightarrow[]{}
\&   \left(\begin{array}{ccc|ccc}
1  \& 0  \& 0   \& 1/4  \& 1/4  \& 3/4  \\\\
0  \& 1  \& 0   \& 1/4  \& 1/4  \& -1/4 \\\\
0  \& 0  \& 1   \& -1/4 \& 3/4  \& -3/4
\end{array}\right)
\end{array}$$

Example 5 [example_5]

A non-invertible matrix is detected by the fact that the left half won''''t
reduce to the identity.

$$\left(\begin{array}{cc|cc}
1  \& 1   \& 1  \& 0  \\\\
2  \& 2   \& 0  \& 1
\end{array}\right)
\xrightarrow[]{-2\rho_1+\rho_2}
\left(\begin{array}{cc|cc}
1  \& 1   \& 1  \& 0  \\\\
0  \& 0   \& -2 \& 1
\end{array}\right)$$

This procedure will find the inverse of a general $n  \times n$
matrix. The $2 \times 2$ case is handy.

\paragraph{Resources [resources_3]}

\begin{itemize}
    \item \href{https://www.khanacademy.org/math/algebra-home/alg-matrices/alg-adding-and-subtracting-matrices/v/matrix-addition-and-subtraction-1}{Matrix Addition and
\end{itemize}
    Subtraction},
    Khan Academy
\begin{itemize}
    \item \href{https://www.khanacademy.org/math/algebra-home/alg-matrices/alg-multiplying-matrices-by-scalars/v/scalar-multiplication}{Multiplying Matrices by
\end{itemize}
    Scalars},
    Khan Academy
\begin{itemize}
    \item \href{https://www.khanacademy.org/math/algebra-home/alg-matrices/alg-multiplying-matrices-by-matrices/v/matrix-multiplication-intro}{Matrix
\end{itemize}
    Multiplication},
    Khan Academy
\begin{itemize}
    \item \href{https://openstax.org/books/college-algebra/pages/7-5-matrices-and-matrix-operations}{Matrices and Matrix
\end{itemize}
    Operations},
    OpenStax
\begin{itemize}
    \item \href{https://www.youtube.com/watch?v=0L90Kkn90J8}{Multiplying a 2-by-3 Matrix and 3-by-2
\end{itemize}
    Matrix}, patrickJMT

\paragraph{Licensing [licensing_3]}

Content obtained and/or adapted from:

\begin{itemize}
    \item \href{https://en.wikibooks.org/wiki/Linear_Algebra/Inverses}{Inverses, Wikibooks: Linear
\end{itemize}
    Algebra}
    under a CC BY-SA license
\begin{itemize}
    \item \href{https://en.wikibooks.org/wiki/Linear_Algebra/The_Inverse_of_a_Matrix}{The Inverse of a Matrix, Wikibooks: Linear
\end{itemize}
    Algebra}
    under a CC BY-SA license

%%===============================================================
\section{Column Space and Nullspace of a Linear Transformation}
\label{sec:column-space-and-nullspace-of-a-linear-transformation}
%%===============================================================

\paragraph{Nullspace}



Among the three important vector spaces associated with a matrix of

order m x n is the \textbf{Null Space}. Null spaces apply to linear

transformations.



\paragraph{Range}



Let T be a linear transformation from an m-dimension vector space X to

an n-dimensional vector space Y, and let $x_1, x_2, x_3, \ldots, x_m$ be

a basis for X and let $y_1, y_2, y_3, \ldots, y_n$ be a basis for Y, and

consider its corresponding n × m matrix,



$M=\begin{pmatrix}

a_{11} \& a_{12} \& a_{13} \& \ldots \& a_{1m} \\\a_{21} \& a_{22} \& a_{23} \& \ldots \& a_{2m} \\\a_{31} \& a_{32} \& a_{33} \& \ldots \& a_{3m} \\\\vdots \& \vdots \& \vdots \& \vdots \& \vdots \\\a_{n1} \& a_{n2} \& a_{n3} \& \ldots \& a_{nm} \\\\end{pmatrix}$.



The image of X, T(X), is called the range of T. T(A) is obviously a

subspace of Y.



Since any element x within X can be expressed as



$x=\sum_{i=1}^m c_i x_i$,



$T(x)=T(\sum_{i=1}^m c_i x_i)=\sum_{i=1}^n c_i T(x_i)$



implying that the range of T is the vector space spanned by the vectors

$T(x_i)$ which is indicated by the columns of the matrix. By a theorem

proven earlier, the dimension of the vector space spanned by those

vectors is equal to the maximum number of vectors that are linearly

independent. Since the linear dependence of columns in the matrix is the

same as the linear dependence of the vectors $T(x_i)$, the dimension is

equal to the maximum number of columns that are linearly independent,

which is equal to the rank. We have the following important conclusion:



The dimension of the range of a linear transformation is equal to the

rank of its corresponding matrix.



\paragraph{Null Space [null_space]}



For example, consider the matrix$$A = \begin{pmatrix}

  1 & 2 \\\    2 & 4 

\end{pmatrix}$$.



The null space of this matrix consists of the set:



$$\textrm{Null Space}(A) = \begin{Bmatrix}{\mathbf{\begin{pmatrix} -2r \\\\ r \end{pmatrix}} : r \in \mathbb{R}} \end{Bmatrix} $$



It may not be immediately obvious how we found this set but it can be

readily checked that any element of this set indeed gives the zero

vector on being multiplied by A. Clearly,



$\begin{pmatrix}

-2 \\\1

\end{pmatrix} \in \textrm{Null Space}(A)$ as



$\begin{pmatrix}

  1 \& 2 \\\    2 \& 4

\end{pmatrix} \begin{pmatrix}

  -2 \\\    1

\end{pmatrix} = \begin{pmatrix}

  0 \\\    0

\end{pmatrix}$.



\paragraph{Null Space as a vector space [null_space_as_a_vector_space]}



It is easy to show that the null space is in fact a vector space. If we

identify a n x 1 column matrix with an element of the n dimensional

Euclidean space then the null space becomes its subspace with the usual

operations. The null space may also be treated as a subspace of the

vector space of all n x 1 column matrices with matrix addition and

scalar multiplication of a matrix as the two operations.



To show that the null space is indeed a vector space it is sufficient to

show that



$x_1,x_2 \in \textrm{Null Space}(A)\Rightarrow x_1 + x_2\in \textrm{Null Space}(A)$



and



$x \in \textrm{Null Space}(A) \Rightarrow \alpha x \in \textrm{Null Space}(A)$



These are true due to the distributive law of matrices. The details of

the proof are left to the reader as an exercise.



\paragraph{Properties}



\paragraph{Null spaces of row equivalent matrices [null_spaces_of_row_equivalent_matrices]}



If A and B are two row equivalent matrices then they share the same null

space. This fact, which is in fact a little theorem, can be proved as

follows:



Suppose x is an element of the null space of A. Then Ax = 0. Also since

A is row equivalent to B so $B = E_1E_2\cdots E_kA$ where each $E_i$ is

an elementary matrix. (Recall that an elementary matrix is the matrix

obtained from I from performing any elementary row operation.) Now,



$Bx = E_1E_2\cdots E_kAx = 0$



and so x is in the null space of B as well. So the null space of A is

contained in that of B. Similarly the null space of B is contained in

that of A. It is now clear that A and B have the same null space.



\paragraph{Basis of Null Space [basis_of_null_space]}



As the null space of a matrix is a vector space, it is natural to wonder

what its basis will be. Of course, since the null space is a subspace of

$\mathbb{R}^n$, its basis can have at most n elements in it. The number

of elements in the basis of the null space is important and is called

the \textbf{nullity} of A. To find out the basis of the null space of A we

follow the following steps:



1.  First convert the given matrix into row echelon form say U.

2.  Next circle the first non zero entries in each row.

3.  Call the variable $x_1$ as a basic variable if the first column has

    a circled entry, and call it a free variable if the first column

    doesn''t have a circled entry. Similarly call the variable $x_2$

    basic if the second column has a non zero entry and free otherwise.

    In this way name n variables $x_1,x_2...,x_n$.

4.  If $x_i$ for any i, is a free variable, then let $v_1$ be the

    solution obtained by solving the system Ux = 0 where all the free

    variables are exactly 0, except for $x_i$ which is 1. If $x_i$ is

    not a free variable don''t do anything.

5.  Repeat the above step for all the free variables getting vectors

    $v_2,v_3$ etc in the process.

6.  The set $\{v_1,v_2...\}$ is the required basis.



The key point in the above algorithm was that A and U have the same null

space. For a complete proof of why the algorithm works we refer the

reader to the excellent text book given in the references by Hoffman and

Kunze.



Let us look at an example:



Suppose $A = \begin{pmatrix}

  1 \& 1 \& 0 \& 1 \& 5 \\\    1 \& 0 \& 0 \& 2 \& 2 \\\    0 \& 0 \& 1 \& 4 \& -1 \\\    0 \& 0 \& 0 \& 0 \& 0 

\end{pmatrix}$



The first step involves reducing A to its row echelon form U.



Now $U = \begin{pmatrix}

  1 \& 0 \& 0 \& 2 \& 2 \\\    0 \& 1 \& 0 \& -1 \& 3 \\\    0 \& 0 \& 1 \& 4 \& -1 \\\    0 \& 0 \& 0 \& 0 \& 0 

\end{pmatrix}$



We encircle the first non zero entries in each row by brackets:



$U = \begin{pmatrix}

  (1) \& 0 \& 0 \& 2 \& 2 \\\    0 \& (1) \& 0 \& -1 \& 3 \\\    0 \& 0 \& (1) \& 4 \& -1 \\\    0 \& 0 \& 0 \& 0 \& 0 

\end{pmatrix}$



Clearly the free variables are $x_4$ and $x_5$ and the rest $x_1,x_2$

and $x_3$ are basic variables. Now we shall solve the system Ux = 0 with

$x_4 = 1, x_5 = 0$ to get the vector $v_1$. Thus we need to solve,



$\begin{pmatrix}

  1 \& 0 \& 0 \& 2 \& 2 \\\    0 \& 1 \& 0 \& -1 \& 3 \\\    0 \& 0 \& 1 \& 4 \& -1 \\\    0 \& 0 \& 0 \& 0 \& 0

\end{pmatrix} \begin{pmatrix}

  x_1\\\    x_2\\\    x_3\\\    1\\\    0

\end{pmatrix} = \begin{pmatrix}

  0\\\    0\\\    0\\\    0

\end{pmatrix}$



This reduces to the following system on matrix multiplication:



$\begin{pmatrix}

  x_1 + 2\\\    x_2 - 1\\\    x_3 + 4\\\    0

\end{pmatrix} = \begin{pmatrix}

  0\\\    0\\\    0\\\    0

\end{pmatrix}$



It is clear from here that

$x_1 = -2,\ x_2 = 1,\ x_3  = -4,\ x_4 = 1,\ x_5 = 0$ is the solution.



Thus $v_1=(-2,\ 1,\ -4,\ 1,\ 0)$. Similarly $v_2$ is found to be

$(-2,\ -3,\ 1,\ 0,\ 1)$.



The set $\{v_1,v_2\}$ is the basis of the null space and the nullity of

the matrix A is 2. In fact this method gives us a way to describe the

null space as well which would be:

$\{\alpha v_1 + \beta v_2 : \alpha,\beta \in \mathbb{R} \}$ (Why? -

Because the linear combination of solutions is also a solution)



\paragraph{Implications of nullity being zero [implications_of_nullity_being_zero]}



The example given above gives no hint as to what happens when there are

no free variables in the row echelon form of A. All we said that in step

4 of our algorithm was that if $x_i$ is not a free variable then *don''t

do anything*. Following that logic, if no variable is free then we keep

on doing nothing, leading to the conclusion that *if no variable is free

then the basis of the null space is an empty set i.e. $\empty$*. In that

case we say that the nullity of the null space is 0. Note that the null

space itself is not empty and contains precisely one element which is

the zero vector.



Now suppose that A is any matrix of order m x n with columns

$c_1, c_2,...c_n$. Each $c_i$ is a vector in the m-dimensional space. If

the nullity of A is zero, then it follows that Ax=0 has only the zero

vector as the solution.



More precisely,



$Ax =  \begin{pmatrix}

  c_1 \& c_2 \& \ldots \& c_n

\end{pmatrix} \begin{pmatrix}

  x_1 \\\    x_2 \\\    \vdots \\\    x_n

\end{pmatrix} = x_1c_1 + x_2c_2 + \ldots + x_nc_n = 0$



has the trivial solution only. This implies that nullity being zero

makes it necessary for the columns of A to be linearly independent. By

retracing our steps we can show that the converse is true as well.



Let us examine the special case of a square matrix, i.e. when m = n. Now

if the nullity is zero then there is no free variable in the row reduced

echelon form of the matrix A, which is say U. Hence each row contains a

pivot, or a leading non zero entry. In that case U must be of the form,

$\begin{bmatrix}

  1 \& 0 \& 0 \& \cdots \& 0\\\    0 \& 1 \& 0 \& \cdots \& 0\\\    \vdots \& \vdots \& \vdots \& \vdots \& \vdots\\\\ 

  0 \& 0 \& 0 \& \cdots \& 1

\end{bmatrix}$



or U must precisely be the identity matrix I. Conversely, if A is row

equivalent to I then Ax = 0 and Ix = 0 have the same solutions, due to

their being equivalent. Since Ix = 0 has only the trivial solution x =

0, so does Ax = 0. It follows that the null space of A is merely {0} and

so the nullity of A is 0.



Thus nullity of A is 0 $\iff$ A is row equivalent to I.



Now if A is row equivalent to I then $A = E_1E_2\cdots E_k$ where each

$E_i$ is an elementary matrix. Since a product of invertible matrices is

invertible and each $E_i$ is invertible so A is invertible. Conversely

if A was invertible, and U its row reduced echelon form then

$U=E_1\cdots E_kA$ which is clearly invertible (by virtue of being a

product of invertible matrices). Now a matrix containing a zero row can

never be invertible (why?), so U has pivots in each row. It follows that

there are n pivots all equal to 1, with zeros above and below them and

so U = I. Thus A is row equivalent to I.



In summary, A is row equivalent to I $\iff$ A is invertible.



We can collect the entire argument in this section, to state the:



\textbf{Theorem}: For a square matrix of order n, the following are

equivalent:



1.  A is invertible.

2.  Nullity of A is 0.

3.  A is row equivalent to the identity matrix.

4.  Columns of A are linearly independent.

5.  The system Ax = 0 has only the trivial solution.

6.  A is a product of elementary matrices.



It will be a good exercise for the reader at this stage to try to

rewrite the proof of the theorem in detail.



\paragraph{Exercises}



1.  Evaluate null spaces and bases for:

    1.  $\begin{pmatrix}

          2 \& 1 \& -2 \\\                    1 \& -2 \& 1 \\\                    -3 \& 1 \& 1 \\\\ 

        \end{pmatrix}$

    2.  $\begin{pmatrix}

          -2 \& -1 \& 4 \& 2 \\\                    1 \& -2 \& 1 \& 1 \\\                    -3 \& 3 \& -5 \& -3 

        \end{pmatrix}$

    3.  $\begin{pmatrix}

          1 \& 1 \& 1 \\\                    1 \& 0 \& 1 \\\                    0 \& 1 \& 1 

        \end{pmatrix}$

2.  Show that null space of a matrix is a vector space.

3.  Prove the theorem regarding invertibility of a square matrix. Also

    by showing that A is invertible iff A$^T$ is, show that the

    condition that the rows are linearly independent can be added to the

    list.

4.  Is the solution set for Ax = b where b is a non zero vector (i.e.

    has at least one component non zero) a vector space? Give reasons.

5.  Let r be the number of basic variables associated with a n order

    matrix A (which is equal to those associated with its row echelon

    form). Show that A is invertible if and only if r = n.



\paragraph{Column and Row Spaces [column_and_row_spaces]}



The \textbf{column space} is an important vector space used in studying an m

x n matrix. If we consider multiplication by a matrix as a sort of

transformation that the vectors undergo, then the null space and the

column space are the two natural collections of vectors which need to be

studied to understand how this transformation works. While the null

space focussed on those vectors which vanished under action of the

matrix (i.e. the solutions of Ax = 0) the column space corresponds to

the transformed vectors themselves (i.e. all Ax). It is the totality of

all the vectors after the transformation. Those readers who have studied

abstract algebra can relate the null space to the kernel of a

homomorphism and the column space to the range.



Another important space associated with the matrix is the \textbf{row space}.

Like its name suggests it is built entirely out of the rows of the

matrix. We shall later see that the row space can be identified with the

column space in a particular sense. In the special case of an invertible

matrix, the row space and the column space are exactly equal.



\paragraph{The Column Space [the_column_space]}



Given an m x n matrix A the column space of A, denoted by C(A), is the

collection of all the vectors formed by linear combinations of the

columns of A.



More precisely if the columns of A are $c_1,c_2 \cdots c_n$, where each

$c_i$ is an m-dimensional vector then,



$$\textrm{C}(A) = \{\mathbf{v} \in \mathbb{R}^m : \mathbf{v} = \mathbf{\alpha_1 c_1 + \alpha_2 c_2 + \cdots \alpha_n c_n},\ \alpha_i\in \mathbb{R}\}$$



The column space is thus essentially built out of the columns of the

matrix. It is the linear span of the columns and in that capacity is

also a vector space with the standard operations. (Recall that span of

any collection of vectors is always a vector space). Also as the columns

are vectors in the m - dimensional space so the column space is

naturally a subspace of $\mathbb{R}^m$.



Clearly if we let,



$x = \begin{pmatrix}

  \alpha_1\\\    \alpha_2\\\    \vdots\\\    \alpha_n

\end{pmatrix}$



then we can say that,



$Ax = \begin{pmatrix}

  c_1 \& c_2 \& \ldots \& c_n

\end{pmatrix} \begin{pmatrix}

  \alpha_1 \\\    \alpha_2 \\\    \vdots \\\    \alpha_n

\end{pmatrix}

= \alpha_1c_1 + \alpha_2c_2 + \ldots + \alpha_nc_n$



Hence our definition can be modified to,



$$\textrm{C}(A) = \{\mathbf{v} \in \mathbb{R}^m : \mathbf{v} = \mathbf{Ax},\ x\in\mathbb{R}^n \}$$



This gives us the idea that the column space is precisely the set of

transformed vectors by the action of multiplication by the matrix.



\paragraph{Basis for the column space [basis_for_the_column_space]}



We shall now prove a theorem regarding the basis for the column space of

a matrix. This constructive proof will also allow us to state a very

important theorem called the \textbf{rank-nullity} theorem as a corollary.



\textbf{Theorem}: The columns of the m x n matrix A corresponding to the

basic variables form a basis of the column space of A.



\textbf{Proof}: First note that the column space of A is formed by the span

of all the columns, $c_1,c_2$...$c_n$. If we take the basis of the null

space of A, then as



$v_1 \in \textrm{Null Space}(A)$



so $Av_1$=0. But this shows us that the column associated with $v_1$ is

a linear combination of the columns associated with the basic variables.

(Note that the contribution of the other free variable columns is 0 and

that of the $v_1$ linked column is 1.) In this way all the free variable

associated columns are linear combinations of the basic variable

associated columns. So all the basic variable linked columns span the

entire column space.



Now suppose that the row echelon form of $A$ is $U$, and if we take the

basic variable columns the submatrix obtained from A is $A_1$ and that

obtained from U is $U_1$. Clearly $U_1$ is the row echelon form of $A_1$

(see exercises) and so the two share the same null space. Now $U_1$ has

only basic variables associated with it because we take only basic

variables from A to A1 and so its nullity is zero. This tells us that

$U_1 x = 0$ has only the trivial solution. Hence $A_1 x = 0$ also has

only the trivial solution and so its nullity is zero as well. It follows

that the columns of $A_1$ which were the basic variable linked columns

of A are linearly independent.



Thus the basic variable columns are both linearly independent and

spanning as a result of which they form a basis of the column space.

\textbf{Q.E.D}



Let us look at an example:



Suppose $A = \begin{pmatrix}

  1 \& 2 \& 0 \& 2 \& 5 \\\    -2 \& -5 \& 1 \& -1 \& -8 \\\    0 \& -3 \& 3 \& 4 \& 1 \\\    3 \& 6 \& 0 \& -7 \& 2 

\end{pmatrix}$



The first step involves reducing A to its row echelon form U.



Now $U = \begin{pmatrix}

  1 \& 0 \& 2 \& 0 \& 1 \\\    0 \& 1 \& -1 \& 0 \& 1 \\\    0 \& 0 \& 0 \& 1 \& 1 \\\    0 \& 0 \& 0 \& 0 \& 0 

\end{pmatrix}$



We encircle the first non zero entries in each row by brackets:



$U = \begin{pmatrix}

  (1) \& 0 \& 2 \& 0 \& 1 \\\    0 \& (1) \& -1 \& 0 \& 1 \\\    0 \& 0 \& 0 \& (1) \& 1 \\\    0 \& 0 \& 0 \& 0 \& 0 

\end{pmatrix}$



Clearly the basic variables are $x_1, x_2$ and $x_4$ and the free

variables are $x_3$ and $x_5$.



So by the theorem, the basis of the column space of A consists of the

columns $c_1, c_2$ and $c_4$ and is given by:



$\begin{Bmatrix}

 \begin{pmatrix}

  1\\\    -2\\\    0\\\    3

    \end{pmatrix},\begin{pmatrix}

  2\\\    -5\\\    -3\\\    6

    \end{pmatrix},\begin{pmatrix}

  2\\\    -1\\\    4\\\    -7 

    \end{pmatrix}

 \end{Bmatrix}$



Note that the nullity of A is 2 which is equal to the difference between

the total number of columns and the number of elements in the column

space.



\paragraph{The Row Space [the_row_space]}



The row space, as, the name suggests is the space built out by the rows

of the matrix. Given an m x n matrix A the row space of A, denoted by

RS(A), is defined as the collection of all the vectors formed by linear

combinations of the rows of A.



More precisely if the rows of A are $r_1,r_2$...$r_m$, where each $r_i$

is an n-dimensional vector then,



$$\textrm{RS}(A) = \{\mathbf{v} \in \mathbb{R}^n : \mathbf{v} = \mathbf{\alpha_1 r_1 + \alpha_2 r_2 + \cdots \alpha_m r_m},\ \alpha_i\in \mathbb{R}\}$$



The row space is thus the linear span of the rows and so is also a

vector space with the standard operations. It is a subspace of

$\mathbb{R}^n$.



\paragraph{Basis of the row space [basis_of_the_row_space]}



The basis of the row space of A consists of precisely the non zero rows

of U where U is the row echelon form of A. This fact is derived from

combining two results which are:



1.  R(A) = R(U) if U is the row echelon form of A.

2.  The non zero rows of a matrix in row echelon form are linearly

    independent.



The proofs are outlined in the exercises.



Let us look at an example:



Suppose $A = \begin{pmatrix}

  1 \& 2 \& 0 \& 2 \& 5 \\\    -2 \& -5 \& 1 \& -1 \& -8 \\\    0 \& -3 \& 3 \& 4 \& 1 \\\    3 \& 6 \& 0 \& -7 \& 2 

\end{pmatrix}$



If we take its row echelon form then we have,



$U = \begin{pmatrix}

  1 \& 0 \& 2 \& 0 \& 1 \\\    0 \& 1 \& -1 \& 0 \& 1 \\\    0 \& 0 \& 0 \& 1 \& 1 \\\    0 \& 0 \& 0 \& 0 \& 0 

\end{pmatrix}$



The first three rows are non zero and so the basis of the row space is:

$\{(1,\ 0,\ 2,\ 0,\ 1),(0,\ 1,\ -1,\ 0,\ 1),(0,\ 0,\ 0,\ 1,\ 1) \}$



\paragraph{Rank of a matrix [rank_of_a_matrix]}



Notice that in our example the basis of the row space has 3 elements

which is the same number as that in the basis of the column space which

we derived earlier. This is not a coincidence. It is in fact always true

that the number of elements in the basis of the row space and that in of

the column space is equal. The steps in the reasoning are as follows:



1.  The number of non zero rows in U is equal to the number of pivots

    (or leading non zero entries in each row) in U.

2.  The number of basic variables by their definition equals the number

    of pivot containing columns (and so equal the number of pivots).

3.  By the previous theorem the number of basic variables equals the

    number of elements in the basis of the column space.

4.  It follows that the bases have equal number of elements.



This unique number of elements in the two bases is called the **rank of

the matrix**. Clearly since the rows of a matrix are the columns of its

transpose, so \textit{a matrix and its transpose have the same rank}. The rank

of a matrix A is often written as Rank (A) just as the nullity is

written as Nullity (A).



\paragraph{The rank nullity theorem [the_rank_nullity_theorem]}



We now proceed to a very important theorem of linear algebra, called the

\textbf{rank nullity theorem}. Another form of this theorem will appear in

the chapter on linear transformations. The theorem, in our context, is :



For an m x n matrix A,



> Rank (A) + Nullity (A) = n = Number of columns of A



The logic behind this theorem is clear. The rank, as we have seen

earlier, corresponds to the number of basic variables and the nullity to

the number of free variables. Since any column is either basic or free

(but not both) obviously the theorem is true.



Readers acquainted with the first isomorphism theorem of groups can note

that this theorem can be related with it. We shall later see how this is

done.



\paragraph{Exercises [exercises_1]}



1\. Evaluate bases for the column and row spaces for:



>   1\. $\begin{pmatrix}

      2 \& 1 \& -2 \\\            1 \& -2 \& 1 \\\            -3 \& 1 \& 1

    \end{pmatrix}$



>   2\. $\begin{pmatrix}

      -2 \& -1 \& 4 \& 2 \\\            1 \& -2 \& 1 \& 1 \\\            -3 \& 3 \& -5 \& -3

    \end{pmatrix}$



>   3\. $\begin{pmatrix}

      1 \& 1 \& 1 \\\            1 \& 0 \& 1 \\\            0 \& 1 \& 1

    \end{pmatrix}$



2\. Show that if A is an m x n matrix then the row space of A coincides

with that of its row echelon form U.



>   \textit{Hint:} Prove that an elementary row transformation on any row of A

    doesn''t alter the fact that it is a linear combination of the rows

    of A.



3\. Show that the non zero rows of a matrix in row echelon form are

linearly independent.



>   \textit{Hint}: Suppose $\alpha_1r_1+\cdots \alpha_mr_m=0$ where $\alpha_i$

    are scalars and $r_i$ the non zero rows. Now if the first pivot

    occurs at the (1,j)th position of the matrix then the jth component

    of each row other then of $r_1$ is 0. So $\alpha_1$ is zero.

    Similarly all other $\alpha_i$''s are also necessarily zero.



4\. Combine the above two results to show that the non zero rows of a

matrix in its row echelon form U form a basis of both R(A) and R(U).



5\. Show that a n order square matrix is invertibile $\iff$ its rank is

n $\iff$ R(A) = C(A) = $\mathbb{R}^n$ $\iff$ Ax = b has a unique

solution for each m dimensional vector b.



6\. Show that if A is a m x n and B a n x p matrix then:



>   1\. C(AB) $\subseteq$ C(A).

>   2\. R(AB) $\subseteq$ R(B).



<!-- -->



>   \textit{Hint}: For (1) note that

    $\{ABx : x\in\mathbb{R}^p\}\subseteq\{Ay : y\in\mathbb{R}^n\}$ and

    for (2) note that

    $R(AB) = C((AB)^T) = C(B^TA^T) \subseteq  C(B^T) = R(B)$



7\. A \textbf{submatrix} of a matrix A is a matrix obtained by deleting some

rows and/or columns of A. Show that for every submatrix C of A, we have

Rank (C) $\le$ Rank (A).



>   \textit{Hint}: Consider a matrix B formed by deleting rows of A not in C.

    Then Rank (B) $\le$ Rank (A) and Rank (C) $\le$ Rank (B).



8\. Show that a m x n matrix A of rank r has at least one r x r

submatrix of rank r, that is, A has an invertible submatrix of order r.



>   \textit{Hint}: Let B be the matrix consisting of r linearly independent row

    vectors of A. Since Rank (B) = r so we can take r linearly

    independent vectors of B to get an r x r invertible submatrix C.



\paragraph{Licensing [licensing_7]}



Content obtained and/or adapted from:



\begin{itemize}
    \item \href{https://en.wikibooks.org/wiki/Linear_Algebra/Null_Spaces}{Nullspaces, Wikibooks: Linear
\end{itemize}
    Algebra}

    under a CC BY-SA license

\begin{itemize}
    \item \href{https://en.wikibooks.org/wiki/Linear_Algebra/Column_and_Row_Spaces}{Column and Row Spaces, Wikibooks: Linear
\end{itemize}
    Algebra}

    under a CC BY-SA license

%%===============================================================
\section{Linear Independence}
\label{sec:linear-independence}
%%===============================================================

\paragraph{Subspaces}

For any vector space, a \textbf{subspace} is a subset that is itself a vector
space, under the inherited operations.

\paragraph{Lemma 1}

For a nonempty subset $S$ of a vector space, under the inherited
operations, the following are equivalent statements.

1.  $S$ is a subspace of that vector space
2.  $S$ is closed under linear combinations of pairs of vectors: for any
    vectors $\vec{s}_1,\vec{s}_2\in S$ and scalars $r_1,r_2$ the vector
    $r_1\vec{s}_1+r_2\vec{s}_2$ is in $S$
3.  $S$ is closed under linear combinations of any number of vectors:
    for any vectors $\vec{s}_1,\ldots,\vec{s}_n\in S$ and scalars
    $r_1, \ldots,r_n$ the vector $r_1\vec{s}_1+\cdots+r_n\vec{s}_n$ is
    in $S$.

Briefly, the way that a subset gets to be a subspace is by being closed
under linear combinations.

>   Proof:
>>   "The following are equivalent" means that each pair of
        statements are equivalent.
>>    $$(1)\!\iff\!(2)
\qquad
(2)\!\iff\!(3)
\qquad
(3)\!\iff\!(1)$$

>>   We will show this equivalence by establishing that
        $(1)\implies (3)\implies (2)\implies (1)$. This strategy is
        suggested by noticing that $(1)\implies (3)$ and
        $(3)\implies (2)$ are easy and so we need only argue the single
        implication $(2)\implies (1)$.

>>   For that argument, assume that $S$ is a nonempty subset of a
        vector space $V$ and that $S$ is closed under combinations of
        pairs of vectors. We will show that $S$ is a vector space by
        checking the conditions.

>>   The first item in the vector space definition has five
        conditions. First, for closure under addition, if
        $\vec{s}_1,\vec{s}_2\in S$ then $\vec{s}_1+\vec{s}_2\in S$, as
        $\vec{s}_1+\vec{s}_2=1\cdot\vec{s}_1+1\cdot\vec{s}_2$.

>>   Second, for any $\vec{s}_1,\vec{s}_2\in S$, because addition is
        inherited from $V$, the sum $\vec{s}_1+\vec{s}_2$ in $S$ equals
        the sum $\vec{s}_1+\vec{s}_2$ in $V$, and that equals the sum
        $\vec{s}_2+\vec{s}_1$ in $V$ (because $V$ is a vector space, its
        addition is commutative), and that in turn equals the sum
        $\vec{s}_2+\vec{s}_1$ in $S$. The argument for the third
        condition is similar to that for the second.

>>   For the fourth, consider the zero vector of $V$ and note that
        closure of $S$ under linear combinations of pairs of vectors
        gives that (where $\vec{s}$ is any member of the nonempty set
        $S$) $0\cdot\vec{s}+0\cdot\vec{s}=\vec{0}$ is in $S$; showing
        that $\vec{0}$ acts under the inherited operations as the
        additive identity of $S$ is easy.

>>   The fifth condition is satisfied because for any $\vec{s}\in S$,
        closure under linear combinations shows that the vector
        $0\cdot\vec{0}+(-1)\cdot\vec{s}$ is in $S$; showing that it is
        the additive inverse of $\vec{s}$ under the inherited operations
        is routine.

We usually show that a subset is a subspace with $(2)\implies (1)$.

\paragraph{Example 1}

>   The plane
    $P=\{\begin{pmatrix} x \\\\
     y \\\\
     z \end{pmatrix}\,\big|\, x+y+z=0\}$
    is a subspace of $\mathbb{R}\^3$. As specified in the definition, the
    operations are the ones inherited from the larger space, that is,
    vectors add in $P$ as they add in $\mathbb{R}\^3$

>   $\begin{pmatrix} x_1 \\\\
     y_1 \\\\
     z_1 \end{pmatrix}+\begin{pmatrix} x_2 \\\\
     y_2 \\\\
     z_2 \end{pmatrix}
     = \begin{pmatrix} x_1+x_2 \\\\
         y_1+y_2 \\\\
         z_1+z_2 \end{pmatrix}$

>   and scalar multiplication is also the same as it is in
    $\mathbb{R}\^3$. To show that $P$ is a subspace, we need only note
    that it is a subset and then verify that it is a space. Checking
    that $P$ satisfies the conditions in the definition of a vector
    space is routine. For instance, for closure under addition, just
    note that if the summands satisfy that $x_1+y_1+z_1=0$ and
    $x_2+y_2+z_2=0$ then the sum satisfies that
    $(x_1+x_2)+(y_1+y_2)+(z_1+z_2)=(x_1+y_1+z_1)+(x_2+y_2+z_2)=0$.

\paragraph{Example 2}

>   The $x$-axis in $\mathbb{R}\^2$ is a subspace where the addition and
    scalar multiplication operations are the inherited ones.

>   $\begin{pmatrix} x_1 \\\\
     0 \end{pmatrix}
        +
        \begin{pmatrix} x_2 \\\\
         0 \end{pmatrix}
        =
        \begin{pmatrix} x_1+x_2 \\\\
         0 \end{pmatrix}
        \qquad
        r\cdot\begin{pmatrix} x \\\\
         0 \end{pmatrix}
        =\begin{pmatrix} rx \\\\
         0 \end{pmatrix}$

>   As above, to verify that this is a subspace, we simply note that it
    is a subset and then check that it satisfies the conditions in
    definition of a vector space. For instance, the two closure
    conditions are satisfied: (1) adding two vectors with a second
    component of zero results in a vector with a second component of
    zero, and (2) multiplying a scalar times a vector with a second
    component of zero results in a vector with a second component of
    zero.

\paragraph{Example 3}

>   Another subspace of $\mathbb{R}\^2$ is

>   $\{\begin{pmatrix} 0 \\\\
     0 \end{pmatrix}\}$

>   which is its \textbf{trivial subspace}.

>   Any vector space has a trivial subspace $\{\vec{0}\,\}$.

At the opposite extreme, any vector space has itself for a subspace.
`{{anchor|improper}}`{=mediawiki}These two are the \textbf{improper}
subspaces. `{{anchor|proper}}`{=mediawiki}Other subspaces are
\textbf{proper}.

\paragraph{Example 4}

The condition in the definition requiring that the addition and scalar
multiplication operations must be the ones inherited from the larger
space is important. Consider the subset $\{1\}$ of the vector space
$\mathbb{R}\^1$. Under the operations $1+1=1$ and $r\cdot 1=1$ that set
is a vector space, specifically, a trivial space. But it is not a
subspace of $\mathbb{R}\^1$ because those aren''t the inherited
operations, since of course $\mathbb{R}\^1$ has $1+1=2$.

\paragraph{Example 5}

>   All kinds of vector spaces, not just $\mathbb{R}\^n$''s, have
    subspaces. The vector space of cubic polynomials
    $\{a+bx+cx\^2+dx\^3\,\big|\, a,b,c,d\in\mathbb{R}\}$ has a subspace
    comprised of all linear polynomials
    $\{m+nx\,\big|\, m,n\in\mathbb{R}\}$.

\paragraph{Example 6}

>   This is a subspace of the $2 \times 2$ matrices

>   $$L=\begin{Bmatrix} \begin{pmatrix}
a  & 0  \\\\
b  & c
\end{pmatrix}
\,\big|\, a+b+c=0\ \end{Bmatrix}$$

>   (checking that it is nonempty and closed under linear combinations
    is easy).
>   To parametrize, express the condition as $a=-b-c$.

>   $$L=\begin{Bmatrix} \begin{pmatrix}
-b-c  & 0  \\\\
b & c
\end{pmatrix}
\,\big|\, b,c \in \mathbb{R} \end{Bmatrix} = \begin{Bmatrix} \begin{pmatrix}
-1 & 0  \\\\
1 & 0
\end{pmatrix}
+c \begin{pmatrix}
-1 & 0  \\\\
0 & 1
\end{pmatrix} \,\big|\, b,c \in \mathbb{R} \end{Bmatrix}$$

>   As above, we''ve described the subspace as a collection of
    unrestricted linear combinations (by coincidence, also of two
    elements).

\paragraph{Span}

The \textbf{span}(or \textbf{linear closure}) of a nonempty subset $S$ of a vector
space is the set of all linear combinations of vectors from $S$.

$$[S] =\{ c_1\vec{s}_1+\cdots+c_n\vec{s}_n
\,\big|\, c_1,\ldots, c_n\in\mathbb{R}
\text{ and } \vec{s}_1,\ldots,\vec{s}_n\in S \}$$ The span of the empty
subset of a vector space is the trivial subspace. No notation for the
span is completely standard. The square brackets used here are common,
but so are "$\textrm{span}(S)$" and "$\textrm{sp}(S)$".

\paragraph{Lemma 2}

In a vector space, the span of any subset is a subspace.

Proof:

>   Call the subset $S$. If $S$ is empty then by definition its span is
    the trivial subspace. If $S$ is not empty, then we need only check
    that the span $[S]$ is closed under linear combinations. For a pair
    of vectors from that span,
    $\vec{v}=c_1\vec{s}_1+\cdots+c_n\vec{s}_n$ and
    $\vec{w}=c_{n+1}\vec{s}_{n+1}+\cdots+c_m\vec{s}_m$, a linear
    combination

>>   $p \cdot (c_1\vec{s}_1+\cdots+c_n\vec{s}_n)+
        r\cdot(c_{n+1}\vec{s}_{n+1}+\cdots+c_m\vec{s}_m)$
>>   $=
        pc_1\vec{s}_1+\cdots+pc_n\vec{s}_n
        +rc_{n+1}\vec{s}_{n+1}+\cdots+rc_m\vec{s}_m$

($p$, $r$ scalars) is a linear combination of elements of $S$ and so
    is in $[S]$ (possibly some of the $\vec{s}_i$''s forming $\vec{v}$
    equal some of the $\vec{s}_j$''s from $\vec{w}$, but it does not
    matter).

\paragraph{Example 7}

The span of this set is all of $\mathbb{R}\^2$.

$$\{\begin{pmatrix} 1 \\\\
 1 \end{pmatrix},\begin{pmatrix} 1 \\\\
 -1 \end{pmatrix}\}$$

To check this we must show that any member of $\mathbb{R}\^2$ is a linear
combination of these two vectors. So we ask: for which vectors (with
real components $x$ and $y$) are there scalars $c_1$ and $c_2$ such that
this holds?

$$c_1\begin{pmatrix} 1 \\\\
 1 \end{pmatrix}+c_2\begin{pmatrix} 1 \\\\
 -1 \end{pmatrix}=\begin{pmatrix} x \\\\
 y \end{pmatrix}$$

Gauss'' method

$$\begin{array}{r}
\begin{array}{r}
c_1  &+  &c_2  &=  &x  \\\\
c_1  &-  &c_2  &=  &y
\end{array}
&\xrightarrow[]{-\rho_1+\rho_2}
&\begin{array}{r}
c_1  &+  &c_2    &=  &x  \\\\
&   &-2c_2  &=  &-x+y
\end{array}
\end{array}$$ with back substitution gives $c_2=(x-y)/2$ and
$c_1=(x+y)/2$. These two equations show that for any $x$ and $y$ that we
start with, there are appropriate coefficients $c_1$ and $c_2$ making
the above vector equation true. For instance, for $x=1$ and $y=2$ the
coefficients $c_2=-1/2$ and $c_1=3/2$ will do. That is, any vector in
$\mathbb{R}\^2$ can be written as a linear combination of the two given
vectors.

\paragraph{Linear Independence}

We first characterize when a vector can be removed from a set without
changing the span of that set.

\paragraph{Lemma 3}

Where $S$ is a subset of a vector space $V$,

$$[S]=[S\cup\{\vec{v}\}]
\quad\text{if and only if}\quad
\vec{v}\in[S]$$ for any $\vec{v}\in V$.

>   Proof: The left to right implication is easy. If
    $[S]=[S\cup\{\vec{v}\}]$ then, since $\vec{v}\in[S\cup\{\vec{v}\}]$,
    the equality of the two sets gives that $\vec{v}\in[S]$.
>   For the right to left implication assume that $\vec{v}\in [S]$ to
    show that $[S]=[S\cup\{\vec{v}\}]$ by mutual inclusion. The
    inclusion is obvious. For the other inclusion
    $[S]\supseteq[S\cup\{\vec{v}\}]$, write an element of
    $[S\cup\{\vec{v}\}]$ as $d_0\vec{v}+d_1\vec{s}_1+\dots+d_m\vec{s}_m$
    and substitute $\vec{v}$''s expansion as a linear combination of
    members of the same set
    $d_0(c_0\vec{t}_0+\dots+c_k\vec{t}_k)+d_1\vec{s}_1+\dots+d_m\vec{s}_m$.
    This is a linear combination of linear combinations and so
    distributing $d_0$ results in a linear combination of vectors from
    $S$. Hence each member of $[S\cup\{\vec{v}\}]$ is also a member of
    $[S]$.
>   Example: In $\mathbb{R}\^3$, where

>   $\vec{v}_1=\begin{pmatrix} 1 \\\\
     0 \\\\
     0 \end{pmatrix}\quad
        \vec{v}_2=\begin{pmatrix} 0 \\\\
         1 \\\\
         0 \end{pmatrix}\quad
        \vec{v}_3=\begin{pmatrix} 2 \\\\
         1 \\\\
         0 \end{pmatrix}$

>   the spans $[\{\vec{v}_1,\vec{v}_2\}]$ and
    $[\{\vec{v}_1,\vec{v}_2,\vec{v}_3\}]$ are equal since $\vec{v}_3$ is
    in the span $[\{\vec{v}_1,\vec{v}_2\}]$.

Lemma 2 says that if we have a spanning set then we can remove a
$\vec{v}$ to get a new set $S$ with the same span if and only if
$\vec{v}$ is a linear combination of vectors from $S$. Thus, under the
second sense described above, a spanning set is minimal if and only if
it contains no vectors that are linear combinations of the others in
that set. We have a term for this important property.

\paragraph{Definition of Linear Independence}

A subset of a vector space is \textbf{linearly independent} if none of its
elements is a linear combination of the others. Otherwise it is
\textbf{linearly dependent}. }}

Here is an important observation:

$$\vec{s}_0=c_1\vec{s}_1+c_2\vec{s}_2+\cdots +c_n\vec{s}_n$$

although this way of writing one vector as a combination of the others
visually sets $\vec{s}_0$ off from the other vectors, algebraically
there is nothing special in that equation about $\vec{s}_0$. For any
$\vec{s}_i$ with a coefficient $c_i$ that is nonzero, we can rewrite the
relationship to set off $\vec{s}_i$.

$$\vec{s}_i=(1/c_i)\vec{s}_0+(-c_1/c_i)\vec{s}_1
+\dots+(-c_n/c_i)\vec{s}_n$$

When we don''t want to single out any vector by writing it alone on one
side of the equation, we will instead say that
$\vec{s}_0,\vec{s}_1,\dots,\vec{s}_n$ are in a \textbf{linear relationship}
and write the relationship with all of the vectors on the same side. The
next result rephrases the linear independence definition in this style.
It gives what is usually the easiest way to compute whether a finite set
is dependent or independent.

\paragraph{Lemma 4}

A subset $S$ of a vector space is linearly independent if and only if
for any distinct $\vec{s}_1,\dots,\vec{s}_n\in S$ the only linear
relationship among those vectors

$$c_1\vec{s}_1+\dots+c_n\vec{s}_n=\vec{0}
\qquad c_1,\dots,c_n\in\mathbb{R}$$

is the trivial one: $c_1=0,\dots,\,c_n=0$. Proof: This is a direct
consequence of the observation above.

>   If the set $S$ is linearly independent then no vector $\vec{s}_i$
    can be written as a linear combination of the other vectors from $S$
    so there is no linear relationship where some of the $\vec{s}\,$''s
    have nonzero coefficients. If $S$ is not linearly independent then
    some $\vec{s}_i$ is a linear combination
    $\vec{s}_i=c_1\vec{s}_1+\dots+c_{i-1}\vec{s}_{i-1} +c_{i+1}\vec{s}_{i+1}+\dots+c_n\vec{s}_n$
    of other vectors from $S$, and subtracting $\vec{s}_i$ from both
    sides of that equation gives a linear relationship involving a
    nonzero coefficient, namely the $-1$ in front of $\vec{s}_i$.

\paragraph{Example 8}

In the vector space of two-wide row vectors, the two-element set
$\{ \begin{pmatrix} 40 &15 \end{pmatrix},\begin{pmatrix} -50 &25 \end{pmatrix}\}$
is linearly independent. To check this, set

$$c_1\cdot\begin{pmatrix} 40 &15 \end{pmatrix}+c_2\cdot\begin{pmatrix} -50 &25 \end{pmatrix}=\begin{pmatrix} 0 &0 \end{pmatrix}$$

and solving the resulting system

$$\begin{array}{r}
40c_1 &- &50c_2 &= &0 \\\\
15c_1 &+ &25c_2 &= &0
\end{array}
\;\xrightarrow[]{-(15/40)\rho_1+\rho_2}\;
\begin{array}{r}
40c_1 &- &50c_2    &= &0 \\\\
& &(175/4)c_2 &= &0
\end{array}$$

shows that both $c_1$ and $c_2$ are zero. So the only linear
relationship between the two given row vectors is the trivial
relationship.

In the same vector space,
$\{ \begin{pmatrix} 40 &15 \end{pmatrix},\begin{pmatrix} 20 &7.5 \end{pmatrix}\}$
is linearly dependent since we can satisfy

$$c_1\begin{pmatrix} 40 &15 \end{pmatrix}+c_2\cdot\begin{pmatrix} 20 &7.5 \end{pmatrix}=\begin{pmatrix} 0 &0 \end{pmatrix}$$

with $c_1=1$ and $c_2=-2$.

\paragraph{Example 9}

The set $\{1+x,1-x\}$ is linearly independent in $\mathcal{P}_2$, the
space of quadratic polynomials with real coefficients, because

$$0+0x+0x^2 = c_1(1+x)+c_2(1-x) =(c_1+c_2)+(c_1-c_2)x+0x^2$$ 

gives

$$\begin{array}{r}
\begin{array}{r}
c_1 &+ &c_2 &= &0 \\\\
c_1 &- &c_2 &= &0
\end{array}
&\xrightarrow[]{-\rho_1+\rho_2}
&\begin{array}{r}
c_1 &+ &c_2 &= &0 \\\\
&  &2c_2 &= &0
\end{array}
\end{array}$$ since polynomials are equal only if their coefficients are
equal. Thus, the only linear relationship between these two members of
$\mathcal{P}_2$ is the trivial one.

\paragraph{Example 10}

In $\mathbb{R}\^3$, where

$$\vec{v}_1=\begin{pmatrix} 3 \\\\
 4 \\\\
 5 \end{pmatrix}
\quad
\vec{v}_2=\begin{pmatrix} 2 \\\\
 9 \\\\
 2 \end{pmatrix}
\quad
\vec{v}_3=\begin{pmatrix} 4 \\\\
 18 \\\\
 4 \end{pmatrix}$$

the set $S=\{\vec{v}_1,\vec{v}_2,\vec{v}_3\}$ is linearly dependent
because this is a relationship

$$0\cdot\vec{v}_1
+2\cdot\vec{v}_2
-1\cdot\vec{v}_3
=\vec{0}$$

where not all of the scalars are zero (the fact that some of the scalars
are zero doesn''t matter).

\paragraph{Resources}

\begin{itemize}
    \item \href{https://textbooks.math.gatech.edu/ila/subspaces.html}{Subspaces},
\end{itemize}
    Interactive Linear Algebra from Georgia Tech

\paragraph{Licensing}

Content obtained and/or adapted from:

\begin{itemize}
    \item \href{https://en.wikibooks.org/wiki/Linear_Algebra/Definition_and_Examples_of_Linear_Independence}{Definition and Examples of Linear Independence, WikiBooks Linear
\end{itemize}
    Algebra}
    under a CC BY-SA license
\begin{itemize}
    \item \href{https://en.wikibooks.org/wiki/Linear_Algebra/Subspaces_and_Spanning_sets}{Subspaces and Spanning Sets, WikiBooks Linear
\end{itemize}
    Algebra}
    under a CC BY-SA license
\begin{itemize}
    \item \href{https://en.wikibooks.org/wiki/Linear_Algebra/Subspaces}{Subspaces, WikiBooks Linear
\end{itemize}
    Algebra}
    under a CC BY-SA license

%%===============================================================
\section{Introduction to Vector Spaces}
\label{sec:introduction-to-vector-spaces}
%%===============================================================

A \textbf{vector space} (over $\R$) consists of a set $V$ along with two
operations "$+$" and "$\cdot$" subject to these conditions.

1.  For any $\vec v,\vec w\in V:\vec v+\vec w\in V$.
2.  For any $\vec v,\vec w\in V:\vec v+\vec w=\vec w+\vec v$.
3.  For any
   $\vec u,\vec v,\vec w\in V:(\vec v+\vec w)+\vec u=\vec v+(\vec w+\vec u)$.
4.  There is a \textbf{zero vector} $\vec0\in V$ such that
   $\vec v+\vec0=\vec v$ for all $\vec v\in V$.
5.  Each $\vec v\in V$ has an \textbf{additive inverse} $\vec w\in V$ such
   that $\vec w+\vec v=\vec0$.
6.  If $r$ is a \textbf{scalar}, that is, a member of $\R$ and $\vec v\in V$
   then the \textbf{scalar multiple} $r\cdot\vec v$ is in $V$.
7.  If $r,s\in\R$ and $\vec v\in V$ then
   $(r+s)\cdot\vec v=r\cdot\vec v+s\cdot\vec v$.
8.  If $r\in\R$ and $\vec v,\vec w\in V$ , then
   $r\cdot(\vec v+\vec w)=r\cdot\vec v+r\cdot\vec w$.
9.  If $r,s\in\R$ and $\vec v\in V$ , then
   $(rs)\cdot\vec v=r\cdot(s\cdot\vec v)$.
10. For any $\vec v\in V$ , $1\cdot\vec v=\vec v$.

Remark: Because it involves two kinds of addition and two kinds of
multiplication, that definition may seem confused. For instance, in
condition 7 "$(r+s)\cdot\vec v=r\cdot\vec v+s\cdot\vec v$", the first
"$+$" is the real number addition operator while the "$+$" to the right
of the equals sign represents vector addition in the structure $V$ .
These expressions aren''t ambiguous because, e.g., $r$ and $s$ are real
numbers so "$r+s$" can only mean real number addition.

Lemma 1.17: In any vector space $V$, for any $\vec{v}\in V$ and
$r\in\mathbb{R}$, we have

1.  $0\cdot\vec{v}=\vec{0}$, and
2.  $(-1\cdot\vec{v})+\vec{v}=\vec{0}$, and
3.  $r\cdot\vec{0}=\vec{0}$.

Proof: For 1, note that
$\vec{v}=(1+0)\cdot\vec{v}=\vec{v}+(0\cdot\vec{v})$. Add to both sides
the additive inverse of $\vec{v}$, the vector $\vec{w}$ such that
$\vec{w}+\vec{v}=\vec{0}$.

$$\begin{array}{rl}
\vec{w}+\vec{v} & =\vec{w}+\vec{v}+0\cdot\vec{v} \\\\
\vec{0} & =\vec{0}+0\cdot\vec{v} \\\\
\vec{0} & =0\cdot\vec{v}
\end{array}$$

The second item is easy:
$(-1\cdot\vec{v})+\vec{v}=(-1+1)\cdot\vec{v}=0\cdot\vec{v}=\vec{0}$
shows that we can write "$-\vec{v}\,$" for the additive inverse of
$\vec{v}$ without worrying about possible confusion with
$(-1)\cdot\vec{v}$.

For 3, this
$r\cdot\vec{0}=r\cdot(0\cdot\vec{0})=(r\cdot 0)\cdot\vec{0}=\vec{0}$
will do.

\paragraph{Example 1 [example_1_2]}

The set $\R^2$ is a vector space if the operations "$+$" and "$\cdot$"
have their usual meaning.

   $$\binom{x_1}{x_2}+\binom{y_1}{y_2}=\binom{x_1+y_1}{x_2+y_2}\qquad r\cdot\binom{x_1}{x_2}=\binom{rx_1}{rx_2}$$

We shall check all of the conditions.

There are five conditions in item 1. For 1, closure of addition, note
that for any $v_1,v_2,w_1,w_2\in\R$ the result of the sum

   $$\binom{v_1}{v_2}+\binom{w_1}{w_2}=\binom{v_1+w_1}{v_2+w_2}$$

is a column array with two real entries, and so is in $\R^2$ . For 2,
that addition of vectors commutes, take all entries to be real numbers
and compute

   $$\binom{v_1}{v_2}+\binom{w_1}{w_2}=\binom{v_1+w_1}{v_2+w_2}=\binom{w_1+v_1}{w_2+v_2}=\binom{w_1}{w_2}+\binom{v_1}{v_2}$$

(the second equality follows from the fact that the components of the
vectors are real numbers, and the addition of real numbers is
commutative). Condition 3, associativity of vector addition, is similar.

   \begin{align}
   \Bigg(\binom{v_1}{v_2}+\binom{w_1}{w_2}\Bigg)+\binom{u_1}{u_2}
   &=\binom{(v_1+w_1)+u_1}{(v_2+w_2)+u_2}\\\   &=\binom{v_1+(w_1+u_1)}{v_2+(w_2+u_2)}\\\   &=\binom{v_1}{v_2}+\Bigg(\binom{w_1}{w_2}+\binom{u_1}{u_2}\Bigg)
   \end{align}

For the fourth condition we must produce a zero element — the vector of
zeroes is it.

   $$\binom{v_1}{v_2}+\binom00=\binom{v_1}{v_2}$$

For 5, to produce an additive inverse, note that for any $v_1,v_2\in\R$
we have

   $$\binom{-v_1}{-v_2}+\binom{v_1}{v_2}=\binom00$$

so the first vector is the desired additive inverse of the second.

The checks for the five conditions having to do with scalar
multiplication are just as routine. For 6, closure under scalar
multiplication, where $r,v_1,v_2\in\R$ ,

   $$r\cdot\binom{v_1}{v_2}=\binom{rv_1}{rv_2}$$

is a column array with two real entries, and so is in $\R^2$ . Next,
this checks 7.

   $$(r+s)\cdot\binom{v_1}{v_2}=\binom{(r+s)v_1}{(r+s)v_2}=
   \binom{rv_1+sv_1}{rv_2+sv_2}=r\cdot\binom{v_1}{v_2}+s\cdot\binom{v_1}{v_2}$$

For 8, that scalar multiplication distributes from the left over vector
addition, we have this.

   $$r\cdot\Bigg(\binom{v_1}{v_2}+\binom{w_1}{w_2}\Bigg)=\binom{r(v_1+w_1)}{r(v_2+w_2)}
   =\binom{rv_1+rw_1}{rv_2+rw_2}=r\cdot\binom{v_1}{v_2}+r\cdot\binom{w_1}{w_2}$$

The ninth

   $$(rs)\cdot\binom{v_1}{v_2}=\binom{(rs)v_1}{(rs)v_2}=\binom{r(sv_1)}{r(sv_2)}=r\cdot\left(s\cdot\binom{v_1}{v_2}\right)$$

and tenth conditions are also straightforward.

   $$1\cdot\binom{v_1}{v_2}=\binom{1v_1}{1v_2}=\binom{v_1}{v_2}$$

In a similar way, each $\mathbb{R}^n$ is a vector space with the usual
operations of vector addition and scalar multiplication. (In
$\mathbb{R}^1$, we usually do not write the members as column vectors,
i.e., we usually do not write "$(\pi)$". Instead we just write "$\pi$".)

\paragraph{Example 2 [example\_2\_2]}

The set $F=\{ a\cos\theta+b\sin\theta \,\big|\, a,b\in\mathbb{R}\}$ of
real-valued functions of the real variable $\theta$ is a vector space
under the operations

$$(a_1\cos\theta+b_1\sin\theta)+(a_2\cos\theta+b_2\sin\theta) =(a_1+a_2)\cos\theta+(b_1+b_2)\sin\theta$$

and

$$r\cdot (a\cos\theta+b\sin\theta)=(ra)\cos\theta+(rb)\sin\theta$$

inherited from the space in the prior example. (We can think of $F$ as
"the same" as $\mathbb{R}^2$ in that $a\cos\theta+b\sin\theta$
corresponds to the vector with components $a$ and $b$.)

\paragraph{Resources [resources\_4]}

\begin{itemize}
    \item \href{https://www.math.uh.edu/\textasciitilde{}jiwenhe/math2331/lectures/sec4\_1.pdf}{Linear Algebra: Vector Spaces and
\end{itemize}
   Subspaces},
   University of Houston

\paragraph{Licensing [licensing\_6]}

Content obtained and/or adapted from:

\begin{itemize}
    \item \href{https://en.wikibooks.org/wiki/Linear\_Algebra/Definition\_and\_Examples\_of\_Vector\_Spaces}{Definition and Examples of Vector Spaces,
\end{itemize}
   WikiBooks}
   under a CC BY-SA license

%%===============================================================
\section{The Dimension of a Vector Space}
\label{sec:the-dimension-of-a-vector-space}
%%===============================================================

In the prior subsection we defined the basis of a vector space, and we
saw that a space can have many different bases. For example, following
the definition of a basis, we saw three different bases for
$\mathbb{R}^2$. So we cannot talk about "the" basis for a vector space.
True, some vector spaces have bases that strike us as more natural than
others, for instance, $\mathbb{R}^2$''s basis $\mathcal{E}_2$ or
$\mathbb{R}^3$''s basis $\mathcal{E}_3$ or $\mathcal{P}_2$''s basis
$\langle 1,x,x^2 \rangle$. But, for example in the space
$\{a_2x^2+a_1x+a_0\,\big|\, 2a_2-a_0=a_1\}$, no particular basis leaps
out at us as the most natural one. We cannot, in general, associate with
a space any single basis that best describes that space.

We can, however, find something about the bases that is uniquely
associated with the space. This subsection shows that any two bases for
a space have the same number of elements. So, with each space we can
associate a number, the number of vectors in any of its bases.

This brings us back to when we considered the two things that could be
meant by the term "minimal spanning set". At that point we defined
"minimal" as linearly independent, but we noted that another reasonable
interpretation of the term is that a spanning set is "minimal" when it
has the fewest number of elements of any set with the same span. At the
end of this subsection, after we have shown that all bases have the same
number of elements, then we will have shown that the two senses of
"minimal" are equivalent.

Before we start, we first limit our attention to spaces where at least
one basis has only finitely many members.

> Definition 2.1:
>>   A vector space is \textbf{finite-dimensional} if it has a basis with
>>     only finitely many vectors.

(One reason for sticking to finite-dimensional spaces is so that the
representation of a vector with respect to a basis is a finitely-tall
vector, and so can be easily written.) From now on we study only
finite-dimensional vector spaces. We shall take the term "vector space"
to mean "finite-dimensional vector space". Other spaces are interesting
and important, but they lie outside of our scope.

To prove the main theorem we shall use a technical result.

> Lemma 2.2 (Exchange Lemma):
>>  Assume that $B=\langle \vec{\beta}_1,\dots,\vec{\beta}_n \rangle$
>     is a basis for a vector space, and that for the vector $\vec{v}$
>     the relationship
>     $\vec{v}=c_1\vec{\beta}_1+c_2\vec{\beta}_2+\cdots +c_n\vec{\beta}_n$
>     has $c_i\neq 0$. Then exchanging $\vec{\beta}_i$ for $\vec{v}$
>     yields another basis for the space.

Proof: Call the outcome of the exchange
$\hat{B}=\langle \vec{\beta}_1,\dots,\vec{\beta}_{i-1},\vec{v},\vec{\beta}_{i+1},\dots,\vec{\beta}_n \rangle$.

We first show that $\hat{B}$ is linearly independent. Any relationship
$d_1\vec{\beta}_1+\dots+d_i\vec{v}+\dots+d_n\vec{\beta}_n=\vec{0}$ among
the members of $\hat{B}$, after substitution for $\vec{v}$,

$$d_1\vec{\beta}_1+\dots
+d_i\cdot(c_1\vec{\beta}_1+\dots+c_i\vec{\beta}_i+\dots+c_n\vec{\beta}_n)
+\dots+d_n\vec{\beta}_n
=\vec{0}
\qquad\qquad(*)$$

gives a linear relationship among the members of $B$. The basis $B$ is
linearly independent, so the coefficient $d_ic_i$ of $\vec{\beta}_i$ is
zero. Because $c_i$ is assumed to be nonzero, $d_i=0$. Using this in
equation $(*)$ above gives that all of the other $d$''s are also zero.
Therefore $\hat{B}$ is linearly independent.

We finish by showing that $\hat{B}$ has the same span as $B$. Half of
this argument, that $[{\hat{B}}]\subseteq[B]$, is easy; any member
$d_1\vec{\beta}_1+\dots+d_i\vec{v}+\dots+d_n\vec{\beta}_n$ of
$[{\hat{B}}]$ can be written
$d_1\vec{\beta}_1+\dots+d_i\cdot(c_1\vec{\beta}_1+\dots+c_n\vec{\beta}_n)+\dots+d_n\vec{\beta}_n$,
which is a linear combination of linear combinations of members of $B$,
and hence is in $[B]$. For the $[B]\subseteq[{\hat{B}}]$ half of the
argument, recall that when
$\vec{v}=c_1\vec{\beta}_1+\dots+c_n\vec{\beta}_n$ with $c_i\neq 0$, then
the equation can be rearranged to
$\vec{\beta}_i=(-c_1/c_i)\vec{\beta}_1+\dots+(1/c_i)\vec{v}+\dots+(-c_n/c_i)\vec{\beta}_n$.
Now, consider any member
$d_1\vec{\beta}_1+\dots+d_i\vec{\beta}_i+\dots+d_n\vec{\beta}_n$ of
$[B]$, substitute for $\vec{\beta}_i$ its expression as a linear
combination of the members of $\hat{B}$, and recognize (as in the first
half of this argument) that the result is a linear combination of linear
combinations, of members of $\hat{B}$, and hence is in $[{\hat{B}}]$.

> Theorem 2.3:
>>   In any finite-dimensional vector space, all of the bases have the
>     same number of elements.

Proof: Fix a vector space with at least one finite basis. Choose, from
among all of this space''s bases, one
$B=\langle \vec{\beta}_1,\dots,\vec{\beta}_n \rangle$ of minimal size.
We will show that any other basis
$D={\langle \vec{\delta}_1,\vec{\delta}_2,\ldots \rangle }$ also has the
same number of members, $n$. Because $B$ has minimal size, $D$ has no
fewer than $n$ vectors. We will argue that it cannot have more than $n$
vectors.

The basis $B$ spans the space and $\vec{\delta}_1$ is in the space, so
$\vec{\delta}_1$ is a nontrivial linear combination of elements of $B$.
By the Exchange Lemma, $\vec{\delta}_1$ can be swapped for a vector from
$B$, resulting in a basis $B_1$, where one element is $\vec{\delta}$ and
all of the $n-1$ other elements are $\vec{\beta}$''s.

The prior paragraph forms the basis step for an induction argument. The
inductive step starts with a basis $B_k$ (for $ 1 \leq k < n $) containing
$k$ members of $D$ and $n-k$ members of $B$. We know that $D$ has at
least $n$ members so there is a $\vec{\delta}_{k+1}$. Represent it as a
linear combination of elements of $B_k$. The key point: in that
representation, at least one of the nonzero scalars must be associated
with a $\vec{\beta}_i$ or else that representation would be a nontrivial
linear relationship among elements of the linearly independent set $D$.
Exchange $\vec{\delta}_{k+1}$ for $\vec{\beta}_i$ to get a new basis
$B_{k+1}$ with one $\vec{\delta}$ more and one $\vec{\beta}$ fewer than
the previous basis $B_k$.

Repeat the inductive step until no $\vec{\beta}$''s remain, so that $B_n$
contains $\vec{\delta}_1,\dots,\vec{\delta}_n$. Now, $D$ cannot have
more than these $n$ vectors because any $\vec{\delta}_{n+1}$ that
remains would be in the span of $B_n$ (since it is a basis) and hence
would be a linear combination of the other $\vec{\delta}$''s,
contradicting that $D$ is linearly independent.

> Definition 2.4:
>>   The \textbf{dimension} of a vector space is the number of vectors in
>     any of its bases.

> Example 2.5: Any basis for $\mathbb{R}^n$ has $n$ vectors since the
> standard basis $\mathcal{E}_n$ has $n$ vectors. Thus, this definition
> generalizes the most familiar use of term, that $\mathbb{R}^n$ is
> $n$-dimensional.

> Example 2.6: The space $\mathcal{P}_n$ of polynomials of degree at
> most $n$ has dimension $n+1$. We can show this by exhibiting any
> basis— $\langle 1,x,\dots,x^n \rangle$ comes to mind— and counting its
> members.

> Example 2.7: A trivial space is zero-dimensional since its basis is
> empty.

Again, although we sometimes say "finite-dimensional" as a reminder, in
the rest of this book all vector spaces are assumed to be
finite-dimensional. An instance of this is that in the next result the
word "space" should be taken to mean "finite-dimensional vector space".

> Corollary 2.8:
>>   No linearly independent set can have a size greater than the
>     dimension of the enclosing space.

Proof: Inspection of the above proof shows that it never uses that $D$
spans the space, only that $D$ is linearly independent.

> Example 2.9: Recall the subspace diagram from the prior section
> showing the subspaces of $\mathbb{R}^3$. Each subspace shown is
> described with a minimal spanning set, for which we now have the term
> "basis". The whole space has a basis with three members, the plane
> subspaces have bases with two members, the line subspaces have bases
> with one member, and the trivial subspace has a basis with zero
> members. When we saw that diagram we could not show that these are the
> only subspaces that this space has. We can show it now. The prior
> corollary proves that the only subspaces of $\mathbb{R}^3$ are either
> three-, two-, one-, or zero-dimensional. Therefore, the diagram
> indicates all of the subspaces. There are no subspaces somehow, say,
> between lines and planes.

> Corollary 2.10:
>>   Any linearly independent set can be expanded to make a basis.

Proof: If a linearly independent set is not already a basis then it must
not span the space. Adding to it a vector that is not in the span
preserves linear independence. Keep adding, until the resulting set does
span the space, which the prior corollary shows will happen after only a
finite number of steps.

> Corollary 2.11:
>>   Any spanning set can be shrunk to a basis.

Proof: Call the spanning set $S$. If $S$ is empty then it is already a
basis (the space must be a trivial space). If $S=\{\vec{0}\}$ then it
can be shrunk to the empty basis, thereby making it linearly
independent, without changing its span.

Otherwise, $S$ contains a vector $\vec{s}_1$ with $\vec{s}_1\neq\vec{0}$
and we can form a basis $B_1=\langle \vec{s}_1 \rangle$. If $[B_1]=[S]$
then we are done.

If not then there is a $\vec{s}_2\in[S]$ such that
$\vec{s}_2\not\in[B_1]$. Let $B_2=\langle \vec{s}_1,\vec{s_2} \rangle$;
if $[B_2]=[S]$ then we are done.

We can repeat this process until the spans are equal, which must happen
in at most finitely many steps.

> Corollary 2.12:
>>   In an $n$-dimensional space, a set of $n$ vectors is linearly
>     independent if and only if it spans the space.

Proof: First we will show that a subset with $n$ vectors is linearly
independent if and only if it is a basis. "If" is trivially true— bases
are linearly independent. "Only if" holds because a linearly independent
set can be expanded to a basis, but a basis has $n$ elements, so this
expansion is actually the set that we began with.

To finish, we will show that any subset with $n$ vectors spans the space
if and only if it is a basis. Again, "if" is trivial. "Only if" holds
because any spanning set can be shrunk to a basis, but a basis has $n$
elements and so this shrunken set is just the one we started with.

The main result of this subsection, that all of the bases in a
finite-dimensional vector space have the same number of elements, is the
single most important result in this book because it describes what
vector spaces and subspaces there can be.

> Remark 2.13: The case of infinite-dimensional vector spaces is
> somewhat controversial. The statement "any infinite-dimensional vector
> space has a basis" is known to be equivalent to a statement called the
> Axiom of Choice. Mathematicians differ philosophically on whether to
> accept or reject this statement as an axiom on which to base
> mathematics (although, the great majority seem to accept it).
> Consequently the question about infinite-dimensional vector spaces is
> still somewhat up in the air.

\paragraph{Licensing [licensing\_8]}

Content obtained and/or adapted from:

\begin{itemize}
    \item \href{https://en.wikibooks.org/wiki/Linear\_Algebra/Dimension}{Dimension, Wikibooks: Linear
\end{itemize}
    Algebra}
    under a CC BY-SA license

%%===============================================================
\section{Dot products and orthogonality}
\label{sec:dot-products-and-orthogonality}
%%===============================================================

We first consider orthogonal projection onto a line. To orthogonally
project a vector $\vec{v}$ onto a line $\ell$, mark the point on the
line at which someone standing on that point could see $\vec{v}$ by
looking straight up or down (from that person''s point of view).

Linalg\_projection\_3.png

The picture shows someone who has walked out on the line until the tip
of $\vec{v}$ is straight overhead. That is, where the line is described
as the span of some nonzero vector
$\ell=\{c\cdot\vec{s}\,\big|\, c\in\mathbb{R}\}$, the person has walked
out to find the coefficient $c_{\vec{p}}\,$ with the property that
$\vec{v}-c_{\vec{p}}\cdot\vec{s}\,$ is orthogonal to
$c_{\vec{p}}\cdot\vec{s}$.

Linalg\_projection\_4.png

We can solve for this coefficient by noting that because
$\vec{v}-c_{\vec{p}}\vec{s}\,$ is orthogonal to a scalar multiple of
$\vec{s}\,$ it must be orthogonal to $\vec{s}\,$ itself, and then the
consequent fact that the dot product
$(\vec{v}-c_{\vec{p}}\vec{s})\cdot \vec{s}\,$ is zero gives that
$c_{\vec{p}}=\vec{v}\cdot \vec{s}/\vec{s}\cdot \vec{s}$.

Definition 1.1

The **orthogonal projection of $\vec{v}$ onto the line spanned by a
nonzero $\vec{s}\,$** is this vector.

$$\textrm{proj}_{[\vec{s}\,]}({\vec{v}})=
\frac{ \vec{v}\cdot\vec{s} }{ \vec{s}\cdot\vec{s} }\vec{s}$$

Remark 1.2:

The wording of that definition says "spanned by $\vec{s}\,$" instead the
more formal "the span of the set $\{\vec{s}\,\}$". This casual first
phrase is common.

Example 1.3:

To orthogonally project the vector $\binom{2}{3}$ onto the line $y=2x$,
we first pick a direction vector for the line. For instance,

$$\vec{s}=\begin{pmatrix} 1 \\\\ 2 \end{pmatrix}$$

will do. Then the calculation is routine.

Linalg\_projection\_5.png

$\frac{ \begin{pmatrix} 2 \\\\ 3 \end{pmatrix}\cdot\begin{pmatrix} 1 \\\\ 2 \end{pmatrix} }{
\begin{pmatrix} 1 \\\\ 2 \end{pmatrix}\cdot\begin{pmatrix} 1 \\\\ 2 \end{pmatrix} }
\cdot\begin{pmatrix} 1 \\\\ 2 \end{pmatrix}
={\displaystyle \frac{8}{5}}\cdot\begin{pmatrix} 1 \\\\ 2 \end{pmatrix}=\begin{pmatrix} 8/5 \\\\ 16/5 \end{pmatrix}$

Example 1.4:

In $\mathbb{R}^3$, the orthogonal projection of a general vector

$$\begin{pmatrix} x \\\\ y \\\\ z \end{pmatrix}$$

onto the $y$-axis is

$$\frac{ \begin{pmatrix} x \\\\ y \\\\ z \end{pmatrix}\cdot\begin{pmatrix} 0 \\\\ 1 \\\\ 0 \end{pmatrix} }{
\begin{pmatrix} 0 \\\\ 1 \\\\ 0 \end{pmatrix}\cdot\begin{pmatrix} 0 \\\\ 1 \\\\ 0 \end{pmatrix} }
\cdot\begin{pmatrix} 0 \\\\ 1 \\\\ 0 \end{pmatrix}=\begin{pmatrix} 0 \\\\ y \\\\ 0 \end{pmatrix}$$

which matches our intuitive expectation.

The picture above with the stick figure walking out on the line until
$\vec{v}$''s tip is overhead is one way to think of the orthogonal
projection of a vector onto a line. We finish this subsection with two
other ways.

Example 1.5

A railroad car left on an east-west track without its brake is pushed by
a wind blowing toward the northeast at fifteen miles per hour; what
speed will the car reach?

Linalg\_railroad\_wind.png

For the wind we use a vector of length $15$ that points toward the
northeast.

$$\vec{v}=\begin{pmatrix} 15\sqrt{1/2} \\\\ 15\sqrt{1/2} \end{pmatrix}$$

The car can only be affected by the part of the wind blowing in the
east-west direction— the part of $\vec{v}$ in the direction of the
$x$-axis is this (the picture has the same perspective as the railroad
car picture above).

Linalg\_railroad\_wind\_direction.png

$\displaystyle \vec{p}=\begin{pmatrix} 15\sqrt{1/2} \\\\ 0 \end{pmatrix}$

So the car will reach a velocity of $15\sqrt{1/2}$ miles per hour toward
the east.

Thus, another way to think of the picture that precedes the definition
is that it shows $\vec{v}$ as decomposed into two parts, the part with
the line (here, the part with the tracks, $\vec{p}$), and the part that
is orthogonal to the line (shown here lying on the north-south axis).
These two are "not interacting" or "independent", in the sense that the
east-west car is not at all affected by the north-south part of the wind
(see \href{Problem 5} ). So the orthogonal
projection of $\vec{v}$ onto the line spanned by $\vec{s}\,$ can be
thought of as the part of $\vec{v}\,$ that lies in the direction of
$\vec{s}$.

Finally, another useful way to think of the orthogonal projection is to
have the person stand not on the line, but on the vector that is to be
projected to the line. This person has a rope over the line and pulls it
tight, naturally making the rope orthogonal to the line.

Linalg\_projection\_6.png

That is, we can think of the projection $\vec{p}\,$ as being the vector
in the line that is closest to $\vec{v}$ (see \href{Problem
11} ).

Example 1.6:

A submarine is tracking a ship moving along the line $y=3x+2$. Torpedo
range is one-half mile. Can the sub stay where it is, at the origin on
the chart below, or must it move to reach a place where the ship will
pass within range?

Linalg\_submarine\_direction.png

The formula for projection onto a line does not immediately apply
because the line doesn''t pass through the origin, and so isn''t the span
of any $\vec{s}$. To adjust for this, we start by shifting the entire
map down two units. Now the line is $y=3x$, which is a subspace, and we
can project to get the point $\vec{p}$ of closest approach, the point on
the line through the origin closest to

$$\vec{v}=\begin{pmatrix} 0 \\\\ -2 \end{pmatrix}$$

the sub''s shifted position.

$$\vec{p}=\frac{\begin{pmatrix} 0 \\\\ -2 \end{pmatrix}\cdot\begin{pmatrix} 1 \\\\ 3 \end{pmatrix}}{
\begin{pmatrix} 1 \\\\ 3 \end{pmatrix}\cdot\begin{pmatrix} 1 \\\\ 3 \end{pmatrix}}
\cdot \begin{pmatrix} 1 \\\\ 3 \end{pmatrix}=\begin{pmatrix} -3/5 \\\\ -9/5 \end{pmatrix}$$

The distance between $\vec{v}$ and $\vec{p}$ is approximately $0.63$
miles and so the sub must move to get in range.

This subsection has developed a natural projection map: orthogonal
projection onto a line. As suggested by the examples, it is often called
for in applications. The next subsection shows how the definition of
orthogonal projection onto a line gives us a way to calculate especially
convienent bases for vector spaces, again something that is common in
applications. The final subsection completely generalizes projection,
orthogonal or not, onto any subspace at all.

\paragraph{Licensing [licensing\_9]}

Content obtained and/or adapted from:

\begin{itemize}
    \item \href{https://en.wikibooks.org/wiki/Linear\_Algebra/Orthogonal\_Projection\_Onto\_a\_Line}{Orthogonal Projection Onto a Line, Wikibooks: Linear
\end{itemize}
    Algebra}
    under a CC BY-SA license

%%===============================================================
\section{Orthonormal Bases and the Gram-Schmidt Process}
\label{sec:orthonormal-bases-and-the-gram-schmidt-process}
%%===============================================================

"The prior subsection suggests that projecting onto the line spanned by
$\vec{s}$ decomposes a vector $\vec{v}$ into two parts

Linalg\_projection\_and\_orthog.png

$\displaystyle \vec{v}=\textrm{proj}_{[\vec{s}\,]}({\vec{v}})
\,+\,\left(\vec{v}-\textrm{proj}_{[\vec{s}\,]}({\vec{v}})\right)$

that are orthogonal and so are "not interacting". We will now develop
that suggestion.

> \textbf{Definition 2.1}
>
>    Vectors $\vec{v}_1,\dots,\vec{v}_k\in\mathbb{R}^n$ are **mutually
>     orthogonal** when any two are orthogonal: if $i\neq j$ then the
>     dot product $\vec{v}_i\cdot\vec{v}_j$ is zero.

> \textbf{Theorem 2.2}
>
>    If the vectors in a set
>     $\{\vec{v}_1,\dots,\vec{v}_k\}\subset\mathbb{R}^n$ are mutually
>     orthogonal and nonzero then that set is linearly independent.

Proof: Consider a linear relationship
$c_1\vec{v}_1+c_2\vec{v}_2+\dots+c_k\vec{v}_k=\vec{0}$. If $i\in [1..k]$
then taking the dot product of $\vec{v}_i$ with both sides of the
equation

$$\begin{array}{rl}
\vec{v}_i\cdot (c_1\vec{v}_1+c_2\vec{v}_2+\dots+c_k\vec{v}_k)
&=\vec{v}_i\cdot\vec{0}   \\\c_i\cdot(\vec{v}_i\cdot\vec{v}_i)
&=0
\end{array}$$

shows, since $\vec{v}_i$ is nonzero, that $c_i$ is zero.

> \textbf{Corollary 2.3}
>
>    If the vectors in a size $k$ subset of a $k$ dimensional space are
>     mutually orthogonal and nonzero then that set is a basis for the
>     space.

Proof: Any linearly independent size $k$ subset of a $k$ dimensional
space is a basis.

Of course, the converse of Corollary 2.3 does not hold— not every basis
of every subspace of $\mathbb{R}^n$ is made of mutually orthogonal
vectors. However, we can get the partial converse that for every
subspace of $\mathbb{R}^n$ there is at least one basis consisting of
mutually orthogonal vectors.

> \textbf{Example 2.4}: The members $\vec{\beta}_1$ and $\vec{\beta}_2$ of
> this basis for $\mathbb{R}^2$ are not orthogonal.
>
> 
> 
> 
> 
>
> $\displaystyle
> B=\left\langle
> \begin{pmatrix} 4 \\\\ 2 \end{pmatrix},
> \begin{pmatrix} 1 \\\\ 3 \end{pmatrix}  \right\rangle$
>
> 
>
> 
>  title="Linalg\_non\_orthog\_basis\_R2.png" height="200" />
> Linalg\_non\_orthog\_basis\_R2.png
> 
>
> 
> 
>
> However, we can derive from $B$ a new basis for the same space that
> does have mutually orthogonal members. For the first member of the new
> basis we simply use $\vec{\beta}_1$.
>
> $$\vec{\kappa}_1=\begin{pmatrix} 4 \\\\ 2 \end{pmatrix}$$
>
> For the second member of the new basis, we take away from
> $\vec{\beta}_2$ its part in the direction of $\vec{\kappa}_1$,
>
> 
> 
> 
> 
>
> $\displaystyle \vec{\kappa}_2=\begin{pmatrix} 1 \\\\ 3 \end{pmatrix}
> -\textrm{proj}_{\scriptstyle[\vec{\kappa}_1]}({\begin{pmatrix} 1 \\\\ 3 \end{pmatrix}})
> =\begin{pmatrix} 1 \\\\ 3 \end{pmatrix}-\begin{pmatrix} 2 \\\\ 1 \end{pmatrix}=\begin{pmatrix} -1 \\\\ 2 \end{pmatrix}$
>
> 
>
> 
>  title="Linalg\_non\_orthog\_basis\_R2\_2.png" height="200" />
> Linalg\_non\_orthog\_basis\_R2\_2.png
> 
>
> 
> 
>
> which leaves the part, $\vec{\kappa}_2$ pictured above, of
> $\vec{\beta}_2$ that is orthogonal to $\vec{\kappa}_1$ (it is
> orthogonal by the definition of the projection onto the span of
> $\vec{\kappa}_1$). Note that, by the corollary,
> $\{\vec{\kappa}_1,\vec{\kappa}_2\}$ is a basis for $\mathbb{R}^2$.

> \textbf{Definition 2.5}:
>
>    An \textbf{orthogonal basis} for a vector space is a basis of mutually
>     orthogonal vectors.

> \textbf{Example 2.6}: To turn this basis for $\mathbb{R}^3$
>
> $$\left\langle
> \begin{pmatrix} 1 \\\\ 1 \\\\ 1 \end{pmatrix},
> \begin{pmatrix} 0 \\\\ 2 \\\\ 0 \end{pmatrix},
> \begin{pmatrix} 1 \\\\ 0 \\\\ 3 \end{pmatrix}
> \right\rangle$$
>
> into an orthogonal basis, we take the first vector as it is given.
>
> $$\vec{\kappa}_1=\begin{pmatrix} 1 \\\\ 1 \\\\ 1 \end{pmatrix}$$
>
> We get $\vec{\kappa}_2$ by starting with the given second vector
> $\vec{\beta}_2$ and subtracting away the part of it in the direction
> of $\vec{\kappa}_1$.
>
> $$\vec{\kappa}_2=\begin{pmatrix} 0 \\\\ 2 \\\\ 0 \end{pmatrix}
> -\textrm{proj}_{[\vec{\kappa}_1]}({\begin{pmatrix} 0 \\\\ 2 \\\\ 0 \end{pmatrix}})
> =\begin{pmatrix} 0 \\\\ 2 \\\\ 0 \end{pmatrix}-\begin{pmatrix} 2/3 \\\\ 2/3 \\\\ 2/3 \end{pmatrix}
> =\begin{pmatrix} -2/3 \\\\ 4/3 \\\\ -2/3 \end{pmatrix}$$
>
> Finally, we get $\vec{\kappa}_3$ by taking the third given vector and
> subtracting the part of it in the direction of $\vec{\kappa}_1$, and
> also the part of it in the direction of $\vec{\kappa}_2$.
>
> $$\vec{\kappa}_3=\begin{pmatrix} 1 \\\\ 0 \\\\ 3 \end{pmatrix}
> -\textrm{proj}_{[\vec{\kappa}_1]}({\begin{pmatrix} 1 \\\\ 0 \\\\ 3 \end{pmatrix}})
> -\textrm{proj}_{[\vec{\kappa}_2]}({\begin{pmatrix} 1 \\\\ 0 \\\\ 3 \end{pmatrix}})
> =\begin{pmatrix} -1 \\\\ 0 \\\\ 1 \end{pmatrix}$$
>
> Again the corollary gives that
>
> $$\left\langle
> \begin{pmatrix} 1 \\\\ 1 \\\\ 1 \end{pmatrix},
> \begin{pmatrix} -2/3 \\\\ 4/3 \\\\ -2/3 \end{pmatrix},
> \begin{pmatrix} -1 \\\\ 0 \\\\ 1 \end{pmatrix}
> \right\rangle$$
>
> is a basis for the space.

The next result verifies that the process used in those examples works
with any basis for any subspace of an $\mathbb{R}^n$ (we are restricted
to $\mathbb{R}^n$ only because we have not given a definition of
orthogonality for other vector spaces).

> \textbf{Theorem 2.7 (Gram-Schmidt orthogonalization)}:
>
>    If $\left\langle \vec{\beta}_1,\ldots\vec{\beta}_k \right\rangle$
>
> is a basis for a subspace of $\mathbb{R}^n$ then, where
>
> $$\begin{array}{rl}
> \vec{\kappa}_1
> &=\vec{\beta}_1    \\\> \vec{\kappa}_2
> &=\vec{\beta}_2-\textrm{proj}_{[\vec{\kappa}_1]}({\vec{\beta}_2})    \\\> \vec{\kappa}_3
> &=\vec{\beta}_3
> -\textrm{proj}_{[\vec{\kappa}_1]}({\vec{\beta}_3})
> -\textrm{proj}_{[\vec{\kappa}_2]}({\vec{\beta}_3})    \\\> &\vdots           \\\> \vec{\kappa}_k
> &=\vec{\beta}_k
> -\textrm{proj}_{[\vec{\kappa}_1]}({\vec{\beta}_k})-\cdots
> -\textrm{proj}_{[\vec{\kappa}_{k-1}]}({\vec{\beta}_k})
> \end{array}$$
>
> the $\vec{\kappa}\,$''s form an orthogonal basis for the same subspace.

Proof: We will use induction to check that each $\vec{\kappa}_i$ is
nonzero, is in the span of
$\left\langle \vec{\beta}_1,\ldots\vec{\beta}_i \right\rangle$ and is
orthogonal to all preceding vectors:
$\vec{\kappa}_1\cdot\vec{\kappa}_i= \cdots
=\vec{\kappa}_{i-1}\cdot\vec{\kappa}_{i}=0$. With those, and with
[Corollary 2.3], we will have
that $\left\langle \vec{\kappa}_1,\ldots\vec{\kappa}_k \right\rangle$ is
a basis for the same space as
$\left\langle \vec{\beta}_1,\ldots\vec{\beta}_k \right\rangle$.

We shall cover the cases up to $i=3$, which give the sense of the
argument. Completing the details is [Problem
15].

The $i=1$ case is trivial— setting $\vec{\kappa}_1$ equal to
$\vec{\beta}_1$ makes it a nonzero vector since $\vec{\beta}_1$ is a
member of a basis, it is obviously in the desired span, and the
"orthogonal to all preceding vectors" condition is vacuously met.

For the $i=2$ case, expand the definition of $\vec{\kappa}_2$.

$$\vec{\kappa}_2=\vec{\beta}_2
-\textrm{proj}_{[\vec{\kappa}_1]}({\vec{\beta}_2})=
\vec{\beta}_2
-\frac{\vec{\beta}_2\cdot\vec{\kappa}_1}{
\vec{\kappa}_1\cdot\vec{\kappa}_1}
\cdot\vec{\kappa}_1 =
\vec{\beta}_2
-\frac{\vec{\beta}_2\cdot\vec{\kappa}_1}{
\vec{\kappa}_1\cdot\vec{\kappa}_1}
\cdot\vec{\beta}_1$$

This expansion shows that $\vec{\kappa}_2$ is nonzero or else this would
be a non-trivial linear dependence among the $\vec{\beta}$''s (it is
nontrivial because the coefficient of $\vec{\beta}_2$ is $1$) and also
shows that $\vec{\kappa}_2$ is in the desired span. Finally,
$\vec{\kappa}_2$ is orthogonal to the only preceding vector

$$\vec{\kappa}_1\cdot\vec{\kappa}_2=
\vec{\kappa}_1\cdot(\vec{\beta}_2
-\textrm{proj}_{[\vec{\kappa}_1]}({\vec{\beta}_2}))=0$$

because this projection is orthogonal.

The $i=3$ case is the same as the $i=2$ case except for one detail. As
in the $i=2$ case, expanding the definition

$$\begin{array}{rl}
\vec{\kappa}_3
&=\vec{\beta}_3
-\frac{\vec{\beta}_3\cdot\vec{\kappa}_1}{
\vec{\kappa}_1\cdot\vec{\kappa}_1}
\cdot\vec{\kappa}_1
-\frac{\vec{\beta}_3\cdot\vec{\kappa}_2}{
\vec{\kappa}_2\cdot\vec{\kappa}_2}
\cdot\vec{\kappa}_2  \\\&=\vec{\beta}_3
-\frac{\vec{\beta}_3\cdot\vec{\kappa}_1}{
\vec{\kappa}_1\cdot\vec{\kappa}_1}
\cdot\vec{\beta}_1
-\frac{\vec{\beta}_3\cdot\vec{\kappa}_2}{
\vec{\kappa}_2\cdot\vec{\kappa}_2}
\cdot\bigl(\vec{\beta}_2
-\frac{\vec{\beta}_2\cdot\vec{\kappa}_1}{
\vec{\kappa}_1\cdot\vec{\kappa}_1}
\cdot\vec{\beta}_1\bigr)
\end{array}$$

shows that $\vec{\kappa}_3$ is nonzero and is in the span. A calculation
shows that $\vec{\kappa}_3$ is orthogonal to the preceding vector
$\vec{\kappa}_1$.

$$\begin{array}{rl}
\vec{\kappa}_1\cdot\vec{\kappa}_3
&=\vec{\kappa}_1\cdot\bigl(\vec{\beta}_3
-\textrm{proj}_{[\vec{\kappa}_1]}({\vec{\beta}_3})
-\textrm{proj}_{[\vec{\kappa}_2]}({\vec{\beta}_3})\bigr) \\\&=\vec{\kappa}_1\cdot\bigl(\vec{\beta}_3
-\textrm{proj}_{[\vec{\kappa}_1]}({\vec{\beta}_3})\bigr)
-\vec{\kappa}_1\cdot\textrm{proj}_{[\vec{\kappa}_2]}({\vec{\beta}_3}) \\\&=0
\end{array}$$

(Here''s the difference from the $i=2$ case— the second line has two
kinds of terms. The first term is zero because this projection is
orthogonal, as in the $i=2$ case. The second term is zero because
$\vec{\kappa}_1$ is orthogonal to $\vec{\kappa}_2$ and so is orthogonal
to any vector in the line spanned by $\vec{\kappa}_2$.) The check that
$\vec{\kappa}_3$ is also orthogonal to the other preceding vector
$\vec{\kappa}_2$ is similar.

Beyond having the vectors in the basis be orthogonal, we can do more; we
can arrange for each vector to have length one by dividing each by its
own length (we can \textbf{normalize} the lengths).

> \textbf{Example 2.8}: Normalizing the length of each vector in the
> orthogonal basis of Example 2.6 produces this \textbf{orthonormal basis}.
>
> $$\left\langle
> \begin{pmatrix} 1/\sqrt{3} \\\\ 1/\sqrt{3} \\\\ 1/\sqrt{3} \end{pmatrix},
> \begin{pmatrix} -1/\sqrt{6} \\\\ 2/\sqrt{6} \\\\ -1/\sqrt{6} \end{pmatrix},
> \begin{pmatrix} -1/\sqrt{2} \\\\ 0 \\\\ 1/\sqrt{2} \end{pmatrix}
> \right\rangle$$

Besides its intuitive appeal, and its analogy with the standard basis
$\mathcal{E}_n$ for $\mathbb{R}^n$, an orthonormal basis also simplifies
some computations.

\paragraph{Licensing [licensing\_10]}

Content obtained and/or adapted from:

\begin{itemize}
    \item \href{https://en.wikibooks.org/wiki/Linear\_Algebra/Gram-Schmidt\_Orthogonalization}{Gram-Schmidt Orthogonalization, Wikibooks: Linear
\end{itemize}
    Algebra}
    under a CC BY-SA license

%%===============================================================
\section{Orthogonal Transformations and Orthogonal Matrices}
\label{sec:orthogonal-transformations-and-orthogonal-matrices}
%%===============================================================

Orthogonal Transformations

In linear algebra, an \textbf{orthogonal transformation} is a linear
transformation \textit{T} : \textit{V} ? \textit{V} on a real inner product space \textit{V}, that
preserves the inner product. That is, for each pair \textit{u}, \textit{v} of elements
of \textit{V}, we have

   $\langle u,v \rangle = \langle Tu,Tv \rangle \, .$

Since the lengths of vectors and the angles between them are defined
through the inner product, orthogonal transformations preserve lengths
of vectors and angles between them. In particular, orthogonal
transformations map orthonormal bases to orthonormal bases.

Orthogonal transformations are injective: if $Tv = 0$ then
$0 = \langle Tv,Tv \rangle = \langle v,v \rangle$, hence $v = 0$, so the
kernel of $T$ is trivial.

Orthogonal transformations in two- or three-dimensional Euclidean space
are stiff rotations, reflections, or combinations of a rotation and a
reflection (also known as improper rotations). Reflections are
transformations that reverse the direction front to back, orthogonal to
the mirror plane, like (real-world) mirrors do. The matrices
corresponding to proper rotations (without reflection) have a
determinant of +1. Transformations with reflection are represented by
matrices with a determinant of -1. This allows the concept of rotation
and reflection to be generalized to higher dimensions.

In finite-dimensional spaces, the matrix representation (with respect to
an orthonormal basis) of an orthogonal transformation is an orthogonal
matrix. Its rows are mutually orthogonal vectors with unit norm, so that
the rows constitute an orthonormal basis of \textit{V}. The columns of the
matrix form another orthonormal basis of \textit{V}.

If an orthogonal transformation is invertible (which is always the case
when \textit{V} is finite-dimensional) then its inverse is another orthogonal
transformation. Its matrix representation is the transpose of the matrix
representation of the original transformation.

\paragraph{Examples}

Consider the inner-product space
$(\mathbb{R}^2,\langle\cdot,\cdot\rangle)$ with the standard euclidean
inner product and standard basis. Then, the matrix transformation

$$T = \begin{bmatrix}
\cos(\theta) & -\sin(\theta) \\\\sin(\theta) & \cos(\theta)
\end{bmatrix} : \mathbb{R}^2 \to \mathbb{R}^2$$ is orthogonal. To see
this, consider

\begin{align}
Te_1 = \begin{bmatrix}\cos(\theta) \\\\ \sin(\theta)\end{bmatrix} && Te_2 = \begin{bmatrix}-\sin(\theta) \\\\ \cos(\theta)\end{bmatrix}
\end{align} Then,

\begin{align}
&\langle Te_1,Te_1\rangle = \begin{bmatrix} \cos(\theta) & \sin(\theta) \end{bmatrix} \cdot \begin{bmatrix} \cos(\theta) \\\\ \sin(\theta) \end{bmatrix} = \cos^2(\theta) + \sin^2(\theta) = 1 \\\\
&\langle Te_1,Te_2\rangle = \begin{bmatrix} \cos(\theta) & \sin(\theta) \end{bmatrix} \cdot \begin{bmatrix} -\sin(\theta) \\\\ \cos(\theta) \end{bmatrix} = \sin(\theta)\cos(\theta) - \sin(\theta)\cos(\theta) = 0 \\\\
&\langle Te_2,Te_2\rangle = \begin{bmatrix} -\sin(\theta) & \cos(\theta)  \end{bmatrix} \cdot \begin{bmatrix} -\sin(\theta) \\\\ \cos(\theta) \end{bmatrix} = \sin^2(\theta) + \cos^2(\theta) = 1\\\\
\end{align}

The previous example can be extended to construct all orthogonal
transformations. For example, the following matrices define orthogonal
transformations on $(\mathbb{R}^3,\langle\cdot,\cdot\rangle)$:

$$\begin{bmatrix}
\cos(\theta) & -\sin(\theta) & 0 \\\\sin(\theta) & \cos(\theta) & 0 \\\0 & 0 & 1
\end{bmatrix},
\begin{bmatrix}
\cos(\theta) & 0 & -\sin(\theta) \\\0 & 1 & 0 \\\\sin(\theta) & 0 & \cos(\theta)
\end{bmatrix},
\begin{bmatrix}
1 & 0 & 0 \\\0 & \cos(\theta) & -\sin(\theta) \\\0 & \sin(\theta) & \cos(\theta) 
\end{bmatrix}$$

\paragraph{Orthogonal Matrix [orthogonal\_matrix]}

In linear algebra, an \textbf{orthogonal matrix}, or \textbf{orthonormal matrix},
is a real square matrix whose columns and rows are orthonormal vectors.

One way to express this is $Q^\mathrm{T} Q = Q Q^\mathrm{T} = I,$ where
$Q^T$ is the transpose of
$Q$ and $I$ is the identity
matrix.

This leads to the equivalent characterization: a matrix
$Q$ is orthogonal if its transpose is equal to its
inverse: $Q^\mathrm{T}=Q^{-1},$ where
$Q^{-1}$ is the inverse of
$Q$.

An orthogonal matrix $Q$ is necessarily invertible
(with inverse
$1=Q^{-1} = Q^T$),
unitary ($1=Q^{-1} = Q^\*$), where
$1=Q^\*$ is the Hermitian adjoint
(conjugate transpose) of $Q$, and therefore normal
($1=Q^\*Q = QQ^\*$)
over the real numbers. The determinant of any orthogonal matrix is
either +1 or -1. As a linear transformation, an orthogonal matrix
preserves the inner product of vectors, and therefore acts as an
isometry of Euclidean space, such as a rotation, reflection or
rotoreflection. In other words, it is a unitary transformation.

The set of $n × n$ orthogonal matrices
forms a group, $O(n)$, known as the orthogonal
group. The subgroup $SO(n)$ consisting of
orthogonal matrices with determinant +1 is called the special orthogonal
group, and each of its elements is a \textbf{special orthogonal matrix}. As a
linear transformation, every special orthogonal matrix acts as a
rotation.

\paragraph{Overview}

An orthogonal matrix is the real specialization of a unitary matrix, and
thus always a normal matrix. Although we consider only real matrices
here, the definition can be used for matrices with entries from any
field. However, orthogonal matrices arise naturally from dot products,
and for matrices of complex numbers that leads instead to the unitary
requirement. Orthogonal matrices preserve the dot product, so, for
vectors $''u''$ and
$''v''$ in an
$n$-dimensional real Euclidean space
${\mathbf u} \cdot {\mathbf v} = \left(Q {\mathbf u}\right) \cdot \left(Q {\mathbf v}\right)$
where $Q$ is an orthogonal matrix. To see the inner
product connection, consider a vector $''v''$ in
an $n$-dimensional real Euclidean space. Written
with respect to an orthonormal basis, the squared length of
$''v''$ is
$''v''^T''v''$. If a linear
transformation, in matrix form $Q''v''$,
preserves vector lengths, then
${\mathbf v}^\mathrm{T}{\mathbf v} = (Q{\mathbf v})^\mathrm{T}(Q{\mathbf v}) = {\mathbf v}^\mathrm{T} Q^\mathrm{T} Q {\mathbf v} .$

Thus finite-dimensional linear isometries—rotations, reflections, and
their combinations—produce orthogonal matrices. The converse is also
true: orthogonal matrices imply orthogonal transformations. However,
linear algebra includes orthogonal transformations between spaces which
may be neither finite-dimensional nor of the same dimension, and these
have no orthogonal matrix equivalent.

Orthogonal matrices are important for a number of reasons, both
theoretical and practical. The $n × n$
orthogonal matrices form a group under matrix multiplication, the
orthogonal group denoted by $O(n)$, which—with
its subgroups—is widely used in mathematics and the physical sciences.
For example, the point group of a molecule is a subgroup of O(3).
Because floating point versions of orthogonal matrices have advantageous
properties, they are key to many algorithms in numerical linear algebra,
such as $QR$ decomposition. As another example,
with appropriate normalization the discrete cosine transform (used in
MP3 compression) is represented by an orthogonal matrix.

\paragraph{Examples [examples\_1]}

Below are a few examples of small orthogonal matrices and possible
interpretations.

\begin{itemize}
    \item $\begin{bmatrix}
\end{itemize}
    1 \& 0 \\\    0 \& 1 \\\    \end{bmatrix}$    (identity transformation)
\begin{itemize}
    \item $\begin{bmatrix}
\end{itemize}
    \cos \theta \& -\sin \theta \\\    \sin \theta \& \cos \theta \\\    \end{bmatrix}$   
\begin{itemize}
    \item $\begin{bmatrix}
\end{itemize}
    \cos 16.26^\circ \& -\sin 16.26^\circ \\\    \sin 16.26^\circ \& \cos 16.26^\circ \\\    \end{bmatrix} \approx \left[\begin{array}{rr}
    0.96 \& -0.28 \\\    0.28 \& 0.96 \\\    \end{array}\right]$(rotation by 16.26°)
\begin{itemize}
    \item $\begin{bmatrix}
\end{itemize}
    1 \& 0 \\\    0 \& -1 \\\    \end{bmatrix}$    (reflection across \textit{x}-axis)
\begin{itemize}
    \item $\begin{bmatrix}
\end{itemize}
    0 \& 0 \& 0 \& 1 \\\    0 \& 0 \& 1 \& 0 \\\    1 \& 0 \& 0 \& 0 \\\    0 \& 1 \& 0 \& 0
    \end{bmatrix}$    (permutation of coordinate axes)

\paragraph{Elementary constructions [elementary\_constructions]}

\paragraph{Lower dimensions [lower\_dimensions]}

The simplest orthogonal matrices are the 1 × 1 matrices \[1\] and
\[-1\], which we can interpret as the identity and a reflection of the
real line across the origin.

The 2 × 2 matrices have the form $\begin{bmatrix}
p \& t\\\q \& u
\end{bmatrix},$ which orthogonality demands satisfy the three equations
$\begin{aligned}
1 \& = p^2+t^2, \\\1 \& = q^2+u^2, \\\0 \& = pq+tu.
\end{aligned}$

In consideration of the first equation, without loss of generality let
$1=p = cos ?$,
$1=q = sin ?$; then either
$1=t = -q$,
$1=u = p$ or
$1=t = q$,
$1=u = -p$. We can interpret the first case
as a rotation by $?$ (where
$1=? = 0$ is the identity), and the second as a
reflection across a line at an angle of $\tfrac{\theta}{2}$.

$\begin{bmatrix}
\cos \theta \& -\sin \theta \\\\sin \theta \& \cos \theta \\\\end{bmatrix}\textrm{ (rotation), }\qquad
\begin{bmatrix}
\cos \theta \& \sin \theta \\\\sin \theta \& -\cos \theta \\\\end{bmatrix}\textrm{ (reflection)}$

The special case of the reflection matrix with
$1=? = 90°$ generates a reflection about the
line at 45° given by $1=y = x$ and
therefore exchanges $x$ and
$y$; it is a permutation matrix, with a single 1 in
each column and row (and otherwise 0): $\begin{bmatrix}
0 \& 1\\\1 \& 0
\end{bmatrix}.$

The identity is also a permutation matrix.

A reflection is its own inverse, which implies that a reflection matrix
is symmetric (equal to its transpose) as well as orthogonal. The product
of two rotation matrices is a rotation matrix, and the product of two
reflection matrices is also a rotation matrix.

\paragraph{Higher dimensions [higher\_dimensions]}

Regardless of the dimension, it is always possible to classify
orthogonal matrices as purely rotational or not, but for 3 × 3 matrices
and larger the non-rotational matrices can be more complicated than
reflections. For example, $\begin{bmatrix}
-1 \& 0 \& 0\\\0 \& -1 \& 0\\\0 \& 0 \& -1
\end{bmatrix}\textrm{ and }
\begin{bmatrix}
0 \& -1 \& 0\\\1 \& 0 \& 0\\\0 \& 0 \& -1
\end{bmatrix}$

represent an \textit{inversion} through the origin and a \textit{rotoinversion},
respectively, about the $z$-axis.

Rotations become more complicated in higher dimensions; they can no
longer be completely characterized by one angle, and may affect more
than one planar subspace. It is common to describe a 3 × 3 rotation
matrix in terms of an axis and angle, but this only works in three
dimensions. Above three dimensions two or more angles are needed, each
associated with a plane of rotation.

However, we have elementary building blocks for permutations,
reflections, and rotations that apply in general.

\paragraph{Primitives}

The most elementary permutation is a transposition, obtained from the
identity matrix by exchanging two rows. Any
$n × n$ permutation matrix can be
constructed as a product of no more than
$n - 1$ transpositions.

A Householder reflection is constructed from a non-null vector
$''v''$ as
$Q = I - 2 \frac{{\mathbf v}{\mathbf v}^\mathrm{T}}{{\mathbf v}^\mathrm{T}{\mathbf v}} .$

Here the numerator is a symmetric matrix while the denominator is a
number, the squared magnitude of $''v''$. This is
a reflection in the hyperplane perpendicular to
$''v''$ (negating any vector component parallel
to $''v''$). If $''v''$ is
a unit vector, then
$1=Q = I - 2''vv''^T$ suffices.
A Householder reflection is typically used to simultaneously zero the
lower part of a column. Any orthogonal matrix of size \textit{n} × \textit{n} can be
constructed as a product of at most $n$ such
reflections.

A Givens rotation acts on a two-dimensional (planar) subspace spanned by
two coordinate axes, rotating by a chosen angle. It is typically used to
zero a single subdiagonal entry. Any rotation matrix of size
$n \times n$ can be constructed as a product of at most
$\tfrac{n(n-1)}{2}$ such rotations. In the case of 3 × 3 matrices, three
such rotations suffice; and by fixing the sequence we can thus describe
all 3 × 3 rotation matrices (though not uniquely) in terms of the three
angles used, often called Euler angles.

A Jacobi rotation has the same form as a Givens rotation, but is used to
zero both off-diagonal entries of a 2 × 2 symmetric submatrix.

\paragraph{Properties [properties\_1]}

\paragraph{Matrix properties [matrix\_properties]}

A real square matrix is orthogonal if and only if its columns form an
orthonormal basis of the Euclidean space
$''R''^n$ with the ordinary
Euclidean dot product, which is the case if and only if its rows form an
orthonormal basis of $''R''^n$. It
might be tempting to suppose a matrix with orthogonal (not orthonormal)
columns would be called an orthogonal matrix, but such matrices have no
special interest and no special name; they only satisfy
$1=M^TM = D$, with
$D$ a diagonal matrix.

The determinant of any orthogonal matrix is +1 or -1. This follows from
basic facts about determinants, as follows:
$1=\det(I)=\det\left(Q^\mathrm{T}Q\right)=\det\left(Q^\mathrm{T}\right)\det(Q)=\bigl(\det(Q)\bigr)^2 .$

The converse is not true; having a determinant of ±1 is no guarantee of
orthogonality, even with orthogonal columns, as shown by the following
counterexample. $\begin{bmatrix}
2 \& 0 \\\0 \& \frac{1}{2}
\end{bmatrix}$

With permutation matrices the determinant matches the signature, being
+1 or -1 as the parity of the permutation is even or odd, for the
determinant is an alternating function of the rows.

Stronger than the determinant restriction is the fact that an orthogonal
matrix can always be diagonalized over the complex numbers to exhibit a
full set of eigenvalues, all of which must have (complex) modulus 1.

\paragraph{Group properties [group\_properties]}

The inverse of every orthogonal matrix is again orthogonal, as is the
matrix product of two orthogonal matrices. In fact, the set of all
$n × n$ orthogonal matrices satisfies all
the axioms of a group. It is a compact Lie group of dimension
$\tfrac{n(n-1)}{2}$, called the orthogonal group and denoted by
$O(n)$.

The orthogonal matrices whose determinant is +1 form a path-connected
normal subgroup of $O(n)$ of index 2, the
special orthogonal group $SO(n)$ of rotations.
The quotient group $O(n)/SO(n)$ is
isomorphic to $O(1)$, with the projection map
choosing \[+1\] or \[-1\] according to the determinant. Orthogonal
matrices with determinant -1 do not include the identity, and so do not
form a subgroup but only a coset; it is also (separately) connected.
Thus each orthogonal group falls into two pieces; and because the
projection map splits, $O(n)$ is a semidirect
product of $SO(n)$ by
$O(1)$. In practical terms, a comparable statement
is that any orthogonal matrix can be produced by taking a rotation
matrix and possibly negating one of its columns, as we saw with 2 × 2
matrices. If $n$ is odd, then the semidirect
product is in fact a direct product, and any orthogonal matrix can be
produced by taking a rotation matrix and possibly negating all of its
columns. This follows from the property of determinants that negating a
column negates the determinant, and thus negating an odd (but not even)
number of columns negates the determinant.

Now consider $(n + 1) × (n + 1)$ orthogonal
matrices with bottom right entry equal to 1. The remainder of the last
column (and last row) must be zeros, and the product of any two such
matrices has the same form. The rest of the matrix is an
$n × n$ orthogonal matrix; thus
$O(n)$ is a subgroup of
$O(n + 1)$ (and of all higher groups).

$\begin{bmatrix}
  \& \& \& 0\\\  \& \mathrm{O}(n) \& \& \vdots\\\  \& \& \& 0\\\  0 \& \cdots \& 0 \& 1
\end{bmatrix}$

Since an elementary reflection in the form of a Householder matrix can
reduce any orthogonal matrix to this constrained form, a series of such
reflections can bring any orthogonal matrix to the identity; thus an
orthogonal group is a reflection group. The last column can be fixed to
any unit vector, and each choice gives a different copy of
$O(n)$ in $O(n + 1)$;
in this way $O(n + 1)$ is a bundle over the
unit sphere $S^n$ with fiber
$O(n)$.

Similarly, $SO(n)$ is a subgroup of
$SO(n + 1)$; and any special orthogonal matrix
can be generated by Givens plane rotations using an analogous procedure.
The bundle structure persists:
$SO(n) ? SO(n + 1) ? S^n$.
A single rotation can produce a zero in the first row of the last
column, and series of $n - 1$ rotations will
zero all but the last row of the last column of an
$n × n$ rotation matrix. Since the planes
are fixed, each rotation has only one degree of freedom, its angle. By
induction, $SO(n)$ therefore has
$(n-1) + (n-2) + \cdots + 1 = \frac{n(n-1)}{2}$ degrees of freedom, and
so does $O(n)$.

Permutation matrices are simpler still; they form, not a Lie group, but
only a finite group, the order $n!$ symmetric
group $Sn$. By the same kind of
argument, $Sn$ is a subgroup of
$Sn + 1$. The even permutations
produce the subgroup of permutation matrices of determinant +1, the
order $\frac{n!}{2}$ alternating group.

\paragraph{Canonical form [canonical\_form]}

More broadly, the effect of any orthogonal matrix separates into
independent actions on orthogonal two-dimensional subspaces. That is, if
$Q$ is special orthogonal then one can always find
an orthogonal matrix $P$, a (rotational) change of
basis, that brings $Q$ into block diagonal form:

$P^\mathrm{T}QP = \begin{bmatrix}
R_1 \& \& \\\\ \& \ddots \& \\\\ \& \& R_k
\end{bmatrix}\ (n\textrm{ even}),
\ P^\mathrm{T}QP = \begin{bmatrix}
R_1 \& \& \& \\\\ \& \ddots \& \& \\\\ \& \& R_k \& \\\\ \& \& \& 1
\end{bmatrix}\ (n\textrm{ odd}).$

where the matrices
$R1, ..., Rk$ are
2 × 2 rotation matrices, and with the remaining entries zero.
Exceptionally, a rotation block may be diagonal,
$±I$. Thus, negating one column if necessary,
and noting that a 2 × 2 reflection diagonalizes to a +1 and -1, any
orthogonal matrix can be brought to the form
$P^\mathrm{T}QP = \begin{bmatrix}
\begin{matrix}R_1 \& \& \\\\ \& \ddots \& \\\\ \& \& R_k\end{matrix} \& 0 \\\0 \& \begin{matrix}\pm 1 \& \& \\\\ \& \ddots \& \\\\ \& \& \pm 1\end{matrix} \\\\end{bmatrix},$

The matrices
$R1, ..., Rk$
give conjugate pairs of eigenvalues lying on the unit circle in the
complex plane; so this decomposition confirms that all eigenvalues have
absolute value 1. If $n$ is odd, there is at least
one real eigenvalue, +1 or -1; for a 3 × 3 rotation, the eigenvector
associated with +1 is the rotation axis.

\paragraph{Lie algebra [lie\_algebra]}

Suppose the entries of $Q$ are differentiable
functions of $t$, and that
$1=t = 0$ gives
$1=Q = I$. Differentiating the
orthogonality condition $Q^\mathrm{T} Q = I$ yields
$\dot{Q}^\mathrm{T} Q + Q^\mathrm{T} \dot{Q} = 0$

Evaluation at $1=t = 0$
($1=Q = I$) then implies
$\dot{Q}^\mathrm{T} = -\dot{Q} .$

In Lie group terms, this means that the Lie algebra of an orthogonal
matrix group consists of skew-symmetric matrices. Going the other
direction, the matrix exponential of any skew-symmetric matrix is an
orthogonal matrix (in fact, special orthogonal).

For example, the three-dimensional object physics calls angular velocity
is a differential rotation, thus a vector in the Lie algebra
$\mathfrak{so}(3)$ tangent to $SO(3)$. Given
$1=''?'' = (x?, y?, z?)$, with
$1=''v'' = (x, y, z)$ being a unit
vector, the correct skew-symmetric matrix form of
$''?''$ is $\Omega = \begin{bmatrix}
0 \& -z\theta \& y\theta \\\z\theta \& 0 \& -x\theta \\\-y\theta \& x\theta \& 0
\end{bmatrix} .$

The exponential of this is the orthogonal matrix for rotation around
axis $''v''$ by angle $?$;
setting $c = cos \tfrac{\theta}{2}$, $s = sin \tfrac{\theta}{2}$,
$\exp(\Omega) = \begin{bmatrix}
1  -  2s^2  +  2x^2 s^2  \&  2xy s^2  -  2z sc  \&  2xz s^2  +  2y sc\\\2xy s^2  +  2z sc  \&  1  -  2s^2  +  2y^2 s^2  \&  2yz s^2  -  2x sc\\\2xz s^2  -  2y sc  \&  2yz s^2  +  2x sc  \&  1  -  2s^2  +  2z^2 s^2   
\end{bmatrix}.$

\paragraph{Numerical linear algebra [numerical\_linear\_algebra]}

\paragraph{Benefits}

Numerical analysis takes advantage of many of the properties of
orthogonal matrices for numerical linear algebra, and they arise
naturally. For example, it is often desirable to compute an orthonormal
basis for a space, or an orthogonal change of bases; both take the form
of orthogonal matrices. Having determinant ±1 and all eigenvalues of
magnitude 1 is of great benefit for numeric stability. One implication
is that the condition number is 1 (which is the minimum), so errors are
not magnified when multiplying with an orthogonal matrix. Many
algorithms use orthogonal matrices like Householder reflections and
Givens rotations for this reason. It is also helpful that, not only is
an orthogonal matrix invertible, but its inverse is available
essentially free, by exchanging indices.

Permutations are essential to the success of many algorithms, including
the workhorse Gaussian elimination with partial pivoting (where
permutations do the pivoting). However, they rarely appear explicitly as
matrices; their special form allows more efficient representation, such
as a list of $n$ indices.

Likewise, algorithms using Householder and Givens matrices typically use
specialized methods of multiplication and storage. For example, a Givens
rotation affects only two rows of a matrix it multiplies, changing a
full multiplication of order $n^3$ to
a much more efficient order $n$. When uses of these
reflections and rotations introduce zeros in a matrix, the space vacated
is enough to store sufficient data to reproduce the transform, and to do
so robustly. (Following Stewart (1976), we do \textit{not} store a rotation
angle, which is both expensive and badly behaved.)

\paragraph{Decompositions}

A number of important matrix decompositions (Golub \& Van Loan 1996)
involve orthogonal matrices, including especially:

$QR$ decomposition : $1=M = QR$, $Q$ orthogonal, $R$ upper triangular  
Singular value decomposition : $1=M = USV^T$, $U$ and $V$ orthogonal, $S$ diagonal matrix  
Eigendecomposition of a symmetric matrix (decomposition according to the spectral theorem) : $1=S = Q?Q^T$, $S$ symmetric, $Q$ orthogonal, $?$ diagonal  
Polar decomposition : $1=M = QS$, $Q$ orthogonal, $S$ symmetric positive-semidefinite

\paragraph{Examples [examples\_2]}

Consider an overdetermined system of linear equations, as might occur
with repeated measurements of a physical phenomenon to compensate for
experimental errors. Write
$1=A''x'' = ''b''$, where
$A$ is $m × n$,
$m > n$. A $QR$
decomposition reduces $A$ to upper triangular
$R$. For example, if $A$ is 5
× 3 then $R$ has the form $R = \begin{bmatrix}
\cdot \& \cdot \& \cdot \\\0 \& \cdot \& \cdot \\\0 \& 0 \& \cdot \\\0 \& 0 \& 0 \\\0 \& 0 \& 0
\end{bmatrix}.$

The linear least squares problem is to find the
$''x''$ that minimizes \|\|\textit{A\textbf{}x} - \textbf{b}\|\|,
which is equivalent to projecting $''b''$ to the
subspace spanned by the columns of $A$. Assuming
the columns of $A$ (and hence
$R$) are independent, the projection solution is
found from
$1=A^TA''x'' = A^T''b''$.
Now $A^TA$ is square
($n × n$) and invertible, and also equal to
$R^TR$. But the lower rows of
zeros in $R$ are superfluous in the product, which
is thus already in lower-triangular upper-triangular factored form, as
in Gaussian elimination (Cholesky decomposition). Here orthogonality is
important not only for reducing
$1=A^TA = (R^TQ^T)QR$
to $R^TR$, but also for allowing
solution without magnifying numerical problems.

In the case of a linear system which is underdetermined, or an otherwise
non-invertible matrix, singular value decomposition (SVD) is equally
useful. With $A$ factored as
$USV^T$, a satisfactory solution
uses the Moore-Penrose pseudoinverse,
$VS^+U^T$, where
$S^+$ merely replaces each non-zero
diagonal entry with its reciprocal. Set $''x''$
to $VS^+U^T''b''$.

The case of a square invertible matrix also holds interest. Suppose, for
example, that $A$ is a 3 × 3 rotation matrix which
has been computed as the composition of numerous twists and turns.
Floating point does not match the mathematical ideal of real numbers, so
$A$ has gradually lost its true orthogonality. A
Gram–Schmidt process could orthogonalize the columns, but it is not the
most reliable, nor the most efficient, nor the most invariant method.
The polar decomposition factors a matrix into a pair, one of which is
the unique \textit{closest} orthogonal matrix to the given matrix, or one of
the closest if the given matrix is singular. (Closeness can be measured
by any matrix norm invariant under an orthogonal change of basis, such
as the spectral norm or the Frobenius norm.) For a near-orthogonal
matrix, rapid convergence to the orthogonal factor can be achieved by a
""Newton''s method"" approach due to Higham (1986) (1990), repeatedly
averaging the matrix with its inverse transpose. Dubrulle (1994) has
published an accelerated method with a convenient convergence test.

For example, consider a non-orthogonal matrix for which the simple
averaging algorithm takes seven steps
$\begin{bmatrix}3 \& 1\\\\7 \& 5\end{bmatrix}
\rightarrow
\begin{bmatrix}1.8125 \& 0.0625\\\\3.4375 \& 2.6875\end{bmatrix}
\rightarrow \cdots \rightarrow
\begin{bmatrix}0.8 \& -0.6\\\\0.6 \& 0.8\end{bmatrix}$ and which
acceleration trims to two steps (with $?$ =
0.353553, 0.565685).

$\begin{bmatrix}3 \& 1\\\\7 \& 5\end{bmatrix}
\rightarrow
\begin{bmatrix}1.41421 \& -1.06066\\\\1.06066 \& 1.41421\end{bmatrix}
\rightarrow
\begin{bmatrix}0.8 \& -0.6\\\\0.6 \& 0.8\end{bmatrix}$

Gram-Schmidt yields an inferior solution, shown by a Frobenius distance
of 8.28659 instead of the minimum 8.12404.

$\begin{bmatrix}3 \& 1\\\\7 \& 5\end{bmatrix}
\rightarrow
\begin{bmatrix}0.393919 \& -0.919145\\\\0.919145 \& 0.393919\end{bmatrix}$

\paragraph{Randomization}

Some numerical applications, such as Monte Carlo methods and exploration
of high-dimensional data spaces, require generation of uniformly
distributed random orthogonal matrices. In this context, ""uniform"" is
defined in terms of Haar measure, which essentially requires that the
distribution not change if multiplied by any freely chosen orthogonal
matrix. Orthogonalizing matrices with independent uniformly distributed
random entries does not result in uniformly distributed orthogonal
matrices, but the $QR$ decomposition of independent
normally distributed random entries does, as long as the diagonal of
$R$ contains only positive entries (Mezzadri 2006).
Stewart (1980) replaced this with a more efficient idea that Diaconis \&
Shahshahani (1987) later generalized as the ""subgroup algorithm"" (in
which form it works just as well for permutations and rotations). To
generate an $(n + 1) × (n + 1)$ orthogonal
matrix, take an $n × n$ one and a uniformly
distributed unit vector of dimension \textit{n} + 1. Construct a Householder
reflection from the vector, then apply it to the smaller matrix
(embedded in the larger size with a 1 at the bottom right corner).

\paragraph{Nearest orthogonal matrix [nearest\_orthogonal\_matrix]}

The problem of finding the orthogonal matrix $Q$
nearest a given matrix $M$ is related to the
Orthogonal Procrustes problem. There are several different ways to get
the unique solution, the simplest of which is taking the singular value
decomposition of $M$ and replacing the singular
values with ones. Another method expresses the $R$
explicitly but requires the use of a matrix square root:
$Q = M \left(M^\mathrm{T} M\right)^{-\frac 1 2}$

This may be combined with the Babylonian method for extracting the
square root of a matrix to give a recurrence which converges to an
orthogonal matrix quadratically:
$Q_{n + 1} = 2 M \left(Q_n^{-1} M + M^\mathrm{T} Q_n\right)^{-1}$ where
$1=Q0 = M$.

These iterations are stable provided the condition number of
$M$ is less than three.

Using a first-order approximation of the inverse and the same
initialization results in the modified iteration:

$N_{n} = Q_n^\mathrm{T} Q_n$ $P_{n} = \frac 1 2 Q_n N_{n}$
$Q_{n + 1} = 2 Q_n + P_n N_n - 3 P_n$

\paragraph{Spin and pin [spin\_and\_pin]}

A subtle technical problem afflicts some uses of orthogonal matrices.
Not only are the group components with determinant +1 and -1 not
connected to each other, even the +1 component,
$SO(n)$, is not simply connected (except for
SO(1), which is trivial). Thus it is sometimes advantageous, or even
necessary, to work with a covering group of SO(\textit{n}), the spin group,
$Spin(n)$. Likewise,
$O(n)"

%%===============================================================
\section{Least Squares}
\label{sec:least-squares}
%%===============================================================

\textbf{Ordinary least squares} (\textbf{OLS}) is a type of linear least squares
method for estimating the unknown parameters in a linear regression
model. OLS chooses the parameters of a linear function of a set of
explanatory variables by the principle of least squares: minimizing the
sum of the squares of the differences between the observed dependent
variable (values of the variable being observed) in the given dataset
and those predicted by the linear function of the independent variable.

Geometrically, this is seen as the sum of the squared distances,
parallel to the axis of the dependent variable, between each data point
in the set and the corresponding point on the regression surface—the
smaller the differences, the better the model fits the data. The
resulting estimator can be expressed by a simple formula, especially in
the case of a simple linear regression, in which there is a single
regressor on the right side of the regression equation.

The OLS estimator is consistent when the regressors are exogenous,
and—by the Gauss–Markov theorem-optimal in the class of linear unbiased
estimators when the errors are homoscedastic and serially uncorrelated.
Under these conditions, the method of OLS provides minimum-variance
mean-unbiased estimation when the errors have finite variances. Under
the additional assumption that the errors are normally distributed, OLS
is the maximum likelihood estimator.

\paragraph{Linear model [linear_model]}

Suppose the data consists of $n$ observations
$\begin{Bmatrix} \mathbf{x}\_i, y\_i \end{Bmatrix}_{i=1} ^n$. Each observation $i$
includes a scalar response $y\_i$ and a column vector $\mathbf{x}\_i$ of
$p$ parameters (regressors), i.e.,
$\mathbf{x}\_i=\left[x\_{i1}, x\_{i2}, \dots, x\_{ip}\right]^\mathsf{T}$. In
a linear regression model, the response variable, $y\_i$, is a linear
function of the regressors:

$$y\_i = \beta\_1\ x\_{i1} + \beta\_2\ x\_{i2} + \cdots + \beta\_p\ x\_{ip} + \varepsilon\_i,$$

or in vector form,

   $y_i = \mathbf{x}_i^\mathsf{T} \boldsymbol{\beta} + \varepsilon_i, \,$

where $\mathbf{x}_i$, as introduced previously, is a column vector of
the $i$-th observation of all the explanatory variables;
$\boldsymbol{\beta}$ is a $p \times 1$ vector of unknown parameters; and
the scalar $\varepsilon_i$ represents unobserved random variables
(\href{errors_and_residuals_in_statistics ""wikilink""}{errors}) of the $i$-th
observation. $\varepsilon_i$ accounts for the influences upon the
responses $y_i$ from sources other than the explanators $\mathbf{x}_i$.
This model can also be written in matrix notation as

   $\mathbf{y} = \mathrm{X} \boldsymbol{\beta} + \boldsymbol{\varepsilon}, \,$

where $\mathbf{y}$ and $\boldsymbol{\varepsilon}$ are $n \times 1$
vectors of the response variables and the errors of the $n$
observations, and $\mathrm{X}$ is an $n \times p$ matrix of regressors,
also sometimes called the design matrix, whose row $i$ is
$\mathbf{x}_i^\mathsf{T}$ and contains the $i$-th observations on all
the explanatory variables.

As a rule, the constant term is always included in the set of regressors
$\mathrm{X}$, say, by taking $x_{i1}=1$ for all $i=1, \dots, n$. The
coefficient $\beta_1$ corresponding to this regressor is called the
\textit{intercept}.

Regressors do not have to be independent: there can be any desired
relationship between the regressors (so long as it is not a linear
relationship). For instance, we might suspect the response depends
linearly both on a value and its square; in which case we would include
one regressor whose value is just the square of another regressor. In
that case, the model would be \textit{quadratic} in the second regressor, but
none-the-less is still considered a \textit{linear} model because the model
\textit{is} still linear in the parameters ($\boldsymbol{\beta}$).

\paragraph{Matrix/vector formulation [matrixvector_formulation]}

Consider an overdetermined system

$$\sum\_{j=1}^{p} X\_{ij} \beta\_j = y\_i,\ (i=1, 2, \dots, n),$$

of $n$ linear equations in $p$ unknown coefficients,
$\beta_1, \beta_2, \dots, \beta_p$, with $n > p$. (Note: for a linear
model as above, not all elements in $\mathrm{X}$ contains information on
the data points. The first column is populated with ones, $X_{i1} = 1$.
Only the other columns contain actual data. So here $p$ is equal to the
number of regressors plus one.) This can be written in matrix form as

$$\mathrm{X} \boldsymbol{\beta} = \mathbf {y},$$

where

$$\mathrm{X}=\begin{bmatrix}
X\_{11} & X\_{12} & \cdots & X\_{1p} \\\X\_{21} & X\_{22} & \cdots & X\_{2p} \\\\vdots & \vdots & \ddots & \vdots \\\X\_{n1} & X\_{n2} & \cdots & X\_{np}
\end{bmatrix} ,
\qquad \boldsymbol \beta = \begin{bmatrix}
\beta\_1 \\\\ \beta\_2 \\\\ \vdots \\\\ \beta\_p \end{bmatrix} ,
\qquad \mathbf y = \begin{bmatrix}
y\_1 \\\\ y\_2 \\\\ \vdots \\\\ y\_n
\end{bmatrix}.$$

Such a system usually has no exact solution, so the goal is instead to
find the coefficients $\boldsymbol{\beta}$ which fit the equations
""best"", in the sense of solving the quadratic minimization problem

$$\hat{\boldsymbol{\beta}} = \underset{\boldsymbol{\beta}}{\operatorname{arg\,min}}\,S(\boldsymbol{\beta}),$$

where the objective function $S$ is given by

$$S(\boldsymbol{\beta}) = \sum\_{i=1}^n \biggl| y\_i - \sum\_{j=1}^p X\_{ij}\beta\_j\biggr|^2 = \bigl\|\mathbf y - \mathrm{X} \boldsymbol \beta \bigr\|^2.$$

A justification for choosing this criterion is given in Properties
below. This minimization problem has a unique solution, provided that
the $p$ columns of the matrix $\mathrm{X}$ are linearly independent,
given by solving the normal equation

$$(\mathrm{X}^{\mathsf T} \mathrm{X} )\hat{\boldsymbol{\beta}} = \mathrm{X}^{\mathsf T} \mathbf y\ .$$

The matrix $\mathrm{X}^{\mathsf T} \mathrm X$ is known as the Gram
matrix and the matrix $\mathrm{X}^{\mathsf T} \mathbf y$ is known as the
moment matrix of regressand by regressors. Finally,
$\hat{\boldsymbol{\beta}}$ is the coefficient vector of the
least-squares hyperplane, expressed as

$$\hat{\boldsymbol{\beta}}= \left( \mathrm{X}^{\mathsf T} \mathrm{X} \right)^{-1} \mathrm{X}^{\mathsf T} \mathbf y.$$

or

$$\hat{\boldsymbol{\beta}}= \boldsymbol{\beta} + (\mathbf{X}^\top \mathbf{X} )^{-1}\mathbf {X}^\top \boldsymbol{\varepsilon}.$$

\paragraph{Licensing [licensing_12]}

Content obtained and/or adapted from:

\begin{itemize}
    \item \href{https://en.wikipedia.org/wiki/Ordinary_least_squares}{Ordinary least squares,
\end{itemize}
    Wikipedia}
    under a CC BY-SA license

%%===============================================================
\section{Determinants}
\label{sec:determinants}
%%===============================================================

The determinant is a function which associates to a square matrix an
element of the field on which it is defined (commonly the real or
complex numbers). The determinant is required to hold these properties:

\begin{itemize}
    \item It is linear on the rows of the matrix.
\end{itemize}

$$\det \begin{bmatrix}  \ddots & \vdots & \ldots \\\\ \lambda a\_1 + \mu b\_1 & \cdots & \lambda a\_n + \mu b\_n \\\\   \cdots & \vdots & \ddots \end{bmatrix} = \lambda \det \begin{bmatrix}  \ddots & \vdots & \cdots \\\\ a\_1 & \cdots & a\_n \\\\   \cdots & \vdots & \ddots \end{bmatrix} + \mu \det \begin{bmatrix}  \ddots & \vdots & \cdots \\\\  b\_1 & \cdots & b\_n \\\\   \cdots & \vdots & \ddots \end{bmatrix}$$

\begin{itemize}
    \item If the matrix has two equal rows its determinant is zero.
    \item The determinant of the identity matrix is 1.
\end{itemize}

It is possible to prove that $\det A = \det A^T$, making the definition
of the determinant on the rows equal to the one on the columns.

\paragraph{Properties [properties_2]}

\begin{itemize}
    \item The determinant is zero if and only if the rows are linearly
\end{itemize}
    dependent.
\begin{itemize}
    \item Changing two rows changes the sign of the determinant:
\end{itemize}

$$\det \begin{bmatrix} \cdots \\\\ \textrm{row A} \\\\ \cdots \\\\ \textrm{row B} \\\\ \cdots \end{bmatrix} = - \det \begin{bmatrix}\cdots \\\\ \textrm{row B} \\\\ \cdots \\\\ \textrm{row A} \\\\ \cdots \end{bmatrix}$$

\begin{itemize}
    \item The determinant is a \textit{multiplicative map} in the sense that
\end{itemize}

$$\det(AB) = \det(A)\det(B) \,$$ for all \textit{n}-by-\textit{n} matrices $A$ and
$B$. This is generalized by the Cauchy-Binet formula to products of
non-square matrices.

\begin{itemize}
    \item It is easy to see that $\det(rI_n) = r^n \,$ and thus
\end{itemize}

$$\det(rA) = \det(rI\_n \cdot A) = r^n \det(A) \,$$ for all $n$-by-$n$
matrices $A$ and all scalars $r$.

\begin{itemize}
    \item A matrix over a commutative ring \textit{R} is invertible if and only if
\end{itemize}
    its determinant is a unit in \textit{R}. In particular, if \textit{A} is a matrix
    over a field such as the real or complex numbers, then \textit{A} is
    invertible if and only if det(\textit{A}) is not zero. In this case we have

$$\det(A^{-1}) = \det(A)^{-1}. \,$$

Expressed differently: the vectors $v_1, \ldots ,v_n$ in $R^n$
form a basis if and only if det( $v_1, \ldots ,v_n$ ) is non-zero.

A matrix and its transpose have the same determinant:

$$\det(A^\top) = \det(A). \,$$

The determinants of a complex matrix and of its conjugate transpose are
conjugate:

$$\det(A^\*) = \det(A)^\*. \,$$

\paragraph{Licensing [licensing_13]}

Content obtained and/or adapted from:

\begin{itemize}
    \item \href{https://en.wikibooks.org/wiki/Linear_Algebra/Determinant}{Determinant,
\end{itemize}
    Wikibooks}
    under a CC BY-SA license

%%===============================================================
\section{Cramer's Rule}
\label{sec:cramer-s-rule}
%%===============================================================

In linear algebra, \textbf{Cramer''s rule} is an explicit formula for the
solution of a system of linear equations with as many equations as
unknowns, valid whenever the system has a unique solution. It expresses
the solution in terms of the determinants of the (square) coefficient
matrix and of matrices obtained from it by replacing one column by the
column vector of right-hand-sides of the equations. It is named after
Gabriel Cramer (1704–1752), who published the rule for an arbitrary
number of unknowns in 1750, (and possibly knew of it as early as 1729).

Cramer''s rule implemented in a naïve way is computationally inefficient
for systems of more than two or three equations. In the case of
`{{mvar|n}}`{=mediawiki} equations in `{{mvar|n}}`{=mediawiki} unknowns,
it requires computation of `{{math|''''n'''' + 1}}`{=mediawiki}
determinants, while Gaussian elimination produces the result with the
same computational complexity as the computation of a single
determinant. Cramer''s rule can also be numerically unstable even for 2×2
systems. However, it has recently been shown that Cramer''s rule can be
implemented in O(\textit{n}^3^) time, which is comparable to more common
methods of solving systems of linear equations, such as Gaussian
elimination (consistently requiring 2.5 times as many arithmetic
operations for all matrix sizes), while exhibiting comparable numeric
stability in most cases.

\paragraph{General case [general_case]}

Consider a system of `{{mvar|n}}`{=mediawiki} linear equations for
`{{mvar|n}}`{=mediawiki} unknowns, represented in matrix multiplication
form as follows:

$$A\mathbf{x} = \mathbf{b}$$

where the `{{math|''''n'''' × ''''n''''}}`{=mediawiki} matrix
`{{mvar|A}}`{=mediawiki} has a nonzero determinant, and the vector
$\mathbf{x} = (x_1, \ldots, x_n)^\mathsf{T}$ is the column vector of the
variables. Then the theorem states that in this case the system has a
unique solution, whose individual values for the unknowns are given by:

$$x\_i = \frac{\det(A\_i)}{\det(A)} \qquad i = 1, \ldots, n$$

where $A_i$ is the matrix formed by replacing the
`{{mvar|i}}`{=mediawiki}-th column of `{{mvar|A}}`{=mediawiki} by the
column vector `{{math|''''''b''''''}}`{=mediawiki}.

A more general version of Cramer''s rule considers the matrix equation

$$AX = B$$

where the `{{math|''''n'''' × ''''n''''}}`{=mediawiki} matrix
`{{mvar|A}}`{=mediawiki} has a nonzero determinant, and
`{{mvar|X}}`{=mediawiki}, `{{mvar|B}}`{=mediawiki} are
`{{math|''''n'''' × ''''m''''}}`{=mediawiki} matrices. Given sequences
$1 \leq i_1 < i_2 < \cdots < i_k \leq n$ and
$1 \leq j_1 < j_2 < \cdots < j_k \leq m$, let $X_{I,J}$ be the
`{{math|''''k'''' × ''''k''''}}`{=mediawiki} submatrix of
`{{mvar|X}}`{=mediawiki} with rows in $I := (i_1, \ldots, i_k )$ and
columns in $J := (j_1, \ldots, j_k )$. Let $A_{B}(I,J)$ be the
`{{math|''''n'''' × ''''n''''}}`{=mediawiki} matrix formed by replacing the
$i_s$ column of `{{mvar|A}}`{=mediawiki} by the $j_s$ column of
`{{Mvar|B}}`{=mediawiki}, for all $s = 1,\ldots, k$. Then

$$\det X\_{I,J} = \frac{\det(A\_{B}(I,J))}{\det(A)}.$$

In the case $k = 1$, this reduces to the normal Cramer''s rule.

The rule holds for systems of equations with coefficients and unknowns
in any field, not just in the real numbers.

\paragraph{Proof}

The proof for Cramer''s rule uses the following properties of the
determinants: linearity with respect to any given column and the fact
that the determinant is zero whenever two columns are equal, which is
implied by the property that the sign of the determinant flips if you
switch two columns.

Fix the index \textit{j} of a column. Linearity means that if we consider only
column \textit{j} as variable (fixing the others arbitrarily), the resulting
function `{{math|''''''R''''''''''n'''' ? ''''''R''''''}}`{=mediawiki}
(assuming matrix entries are in `{{math|''''''R''''''}}`{=mediawiki}) can be
given by a matrix, with one row and \textit{n} columns, that acts on column
\textit{j}. In fact this is precisely what Laplace expansion does, writing
`{{math|1=det(''''A'''') = ''''C''''1''''a''''1,''''j'''' + ? + ''''Cnan,j''''}}`{=mediawiki}
for certain coefficients \textit{C}\textasciitilde{}1\textasciitilde{}, ..., \textit{C\textasciitilde{}n\textasciitilde{}} that depend on the columns
of `{{mvar|A}}`{=mediawiki} other than column \textit{j} (the precise
expression for these cofactors is not important here). The value
`{{math|det(''''A'''')}}`{=mediawiki} is then the result of applying the
one-line matrix
`{{math|1=''''L''''(''''j'''') = (''''C''''1 ''''C''''2 ? ''''Cn'''')}}`{=mediawiki}
to column \textit{j} of `{{mvar|A}}`{=mediawiki}. If
`{{math|''''L''''(''''j'''')}}`{=mediawiki} is applied to any \textit{other}
column \textit{k} of `{{mvar|A}}`{=mediawiki}, then the result is the
determinant of the matrix obtained from `{{mvar|A}}`{=mediawiki} by
replacing column \textit{j} by a copy of column \textit{k}, so the resulting
determinant is 0 (the case of two equal columns).

Now consider a system of `{{mvar|n}}`{=mediawiki} linear equations in
`{{mvar|n}}`{=mediawiki} unknowns $x_1, \ldots,x_n$, whose coefficient
matrix is `{{mvar|A}}`{=mediawiki}, with det(\textit{A}) assumed to be nonzero:

$$\begin{matrix}
a\_{11}x\_1+a\_{12}x\_2+\cdots+a\_{1n}x\_n&=&b\_1\\\a\_{21}x\_1+a\_{22}x\_2+\cdots+a\_{2n}x\_n&=&b\_2\\\&\vdots&\\\a\_{n1}x\_1+a\_{n2}x\_2+\cdots+a\_{nn}x\_n&=&b\_n.
\end{matrix}$$

If one combines these equations by taking \textit{C}\textasciitilde{}1\textasciitilde{} times the first
equation, plus \textit{C}\textasciitilde{}2\textasciitilde{} times the second, and so forth until \textit{C}\textasciitilde{}\textit{n}\textasciitilde{}
times the last, then the coefficient of
`{{mvar|xj}}`{=mediawiki} will become
`{{math|1=''''C''''1''''a''''1, ''''j'''' + ? + ''''Cnan,j'''' = det(''''A'''')}}`{=mediawiki},
while the coefficients of all other unknowns become 0; the left hand
side becomes simply det(\textit{A})\textit{x\textasciitilde{}j\textasciitilde{}}. The right hand side is
`{{math|''''C''''1''''b''''1 + ? + ''''Cnbn''''}}`{=mediawiki},
which is `{{math|''''L''''(''''j'''')}}`{=mediawiki} applied to the
column vector \textbf{b} of the right hand side
`{{mvar|bi}}`{=mediawiki}. In fact what has been done here is
multiply the matrix equation \textit{A\textbf{}x} = \textbf{b} on the left by
`{{math|''''L''''(''''j'''')}}`{=mediawiki}. Dividing by the nonzero
number det(\textit{A}) one finds the following equation, necessary to satisfy
the system:

$$x\_j=\frac{L\_{(j)}\cdot\mathbf{b}}{\det(A)}.$$

But by construction the numerator is the determinant of the matrix
obtained from `{{mvar|A}}`{=mediawiki} by replacing column \textit{j} by \textbf{b},
so we get the expression of Cramer''s rule as a necessary condition for a
solution. The same procedure can be repeated for other values of \textit{j} to
find values for the other unknowns.

The only point that remains to prove is that these values for the
unknowns, the only possible ones, do indeed together form a solution.
But if the matrix `{{mvar|A}}`{=mediawiki} is invertible with inverse
`{{math|''''A''''-1}}`{=mediawiki}, then \textbf{x} = \textit{A}^-1^\textbf{b}
will be a solution, thus showing its existence. To see that
`{{mvar|A}}`{=mediawiki} is invertible when det(\textit{A}) is nonzero,
consider the `{{math|''''n'''' × ''''n''''}}`{=mediawiki} matrix \textit{M} obtained by
stacking the one-line matrices
`{{math|''''L''''(''''j'''')}}`{=mediawiki} on top of each other for
\textit{j} = 1, ..., \textit{n} (this gives the adjugate matrix for
`{{mvar|A}}`{=mediawiki}). It was shown that
`{{math|1=''''L''''(''''j'''')''''A'''' = (0 ? 0 det(''''A'''') 0 ? 0)}}`{=mediawiki}
where `{{math|det(''''A'''')}}`{=mediawiki} appears at the position \textit{j};
from this it follows that \textit{MA} = det(\textit{A})\textit{I\textasciitilde{}n\textasciitilde{}}. Therefore,

$$\frac1{\det(A)}M=A^{-1},$$

completing the proof.

For other proofs, see below.

\paragraph{Finding inverse matrix [finding_inverse_matrix]}

Let `{{mvar|A}}`{=mediawiki} be an `{{math|''''n'''' × ''''n''''}}`{=mediawiki}
matrix with entries in a field `{{math|''''F''''}}`{=mediawiki}. Then

$$A\,\operatorname{adj}(A) = \operatorname{adj}(A)\,A=\det(A) I$$

where `{{math|adj(''''A'''')}}`{=mediawiki} denotes the adjugate matrix,
`{{math|det(''''A'''')}}`{=mediawiki} is the determinant, and
`{{math|''''I''''}}`{=mediawiki} is the identity matrix. If
`{{math|det(''''A'''')}}`{=mediawiki} is nonzero, then the inverse matrix of
`{{mvar|A}}`{=mediawiki} is

$$A^{-1} = \frac{1}{\det(A)} \operatorname{adj}(A).$$

This gives a formula for the inverse of `{{mvar|A}}`{=mediawiki},
provided `{{math|det(''''A'''') ? 0}}`{=mediawiki}. In fact, this formula
works whenever `{{math|''''F''''}}`{=mediawiki} is a commutative ring,
provided that `{{math|det(''''A'''')}}`{=mediawiki} is a unit. If
`{{math|det(''''A'''')}}`{=mediawiki} is not a unit, then
`{{mvar|A}}`{=mediawiki} is not invertible over the ring (it may be
invertible over a larger ring in which some non-unit elements of
`{{mvar|F}}`{=mediawiki} may be invertible).

\paragraph{Applications [applications_1]}

\paragraph{Explicit formulas for small systems [explicit_formulas_for_small_systems]}

Consider the linear system

$$\left\{\begin{matrix}
a\_1x + b\_1y&= {\color{red}c\_1}\\\a\_2x + b\_2y&= {\color{red}c\_2}
\end{matrix}\right.$$

which in matrix format is

$$\begin{bmatrix} a\_1 & b\_1 \\\\ a\_2 & b\_2 \end{bmatrix}\begin{bmatrix} x \\\\ y \end{bmatrix}=\begin{bmatrix} {\color{red}c\_1} \\\\ {\color{red}c\_2} \end{bmatrix}.$$

Assume
`{{math|''''a''''1''''b''''2 - ''''b''''1''''a''''2}}`{=mediawiki}
nonzero. Then, with help of determinants, `{{mvar|x}}`{=mediawiki} and
`{{mvar|y}}`{=mediawiki} can be found with Cramer''s rule as

\begin{align}
x &= \frac{\begin{vmatrix} {\color{red}{c_1}} & b_1 \\\\ {\color{red}{c_2}} & b_2 \end{vmatrix}}{\begin{vmatrix} a_1 & b_1 \\\\ a_2 & b_2 \end{vmatrix}} = { {\color{red}c_1}b_2 - b_1{\color{red}c_2} \over a_1b_2 - b_1a_2}, \quad
y = \frac{\begin{vmatrix} a_1 & {\color{red}{c_1}} \\\\ a_2 & {\color{red}{c_2}} \end{vmatrix}}{\begin{vmatrix} a_1 & b_1 \\\\ a_2 & b_2 \end{vmatrix}}  = { a_1{\color{red}c_2} - {\color{red}c_1}a_2 \over a_1b_2 - b_1a_2}
\end{aligned}.$$

The rules for `{{math|3 × 3}}`{=mediawiki} matrices are similar. Given

$$\left\{\begin{matrix}
a_1x + b_1y + c_1z&= {\color{red}d_1}\\\a_2x + b_2y + c_2z&= {\color{red}d_2}\\\a_3x + b_3y + c_3z&= {\color{red}d_3}
\end{matrix}\right.$$

which in matrix format is

$$\begin{bmatrix} a_1 & b_1 & c_1 \\\\ a_2 & b_2 & c_2 \\\\ a_3 & b_3 & c_3 \end{bmatrix}\begin{bmatrix} x \\\\ y \\\\ z \end{bmatrix}=\begin{bmatrix} {\color{red}d_1} \\\\ {\color{red}d_2} \\\\ {\color{red}d_3} \end{bmatrix}.$$

Then the values of `{{mvar|x, y}}`{=mediawiki} and
`{{mvar|z}}`{=mediawiki} can be found as follows:

$$x = \frac{\begin{vmatrix} {\color{red}d_1} & b_1 & c_1 \\\\ {\color{red}d_2} & b_2 & c_2 \\\\ {\color{red}d_3} & b_3 & c_3 \end{vmatrix} } { \begin{vmatrix} a_1 & b_1 & c_1 \\\\ a_2 & b_2 & c_2 \\\\ a_3 & b_3 & c_3 \end{vmatrix}}, \quad
y = \frac {\begin{vmatrix} a_1 & {\color{red}d_1} & c_1 \\\\ a_2 & {\color{red}d_2} & c_2 \\\\ a_3 & {\color{red}d_3} & c_3 \end{vmatrix}} {\begin{vmatrix} a_1 & b_1 & c_1 \\\\ a_2 & b_2 & c_2 \\\\ a_3 & b_3 & c_3 \end{vmatrix}}, \text{ and }
z = \frac { \begin{vmatrix} a_1 & b_1 & {\color{red}d_1} \\\\ a_2 & b_2 & {\color{red}d_2} \\\\ a_3 & b_3 & {\color{red}d_3} \end{vmatrix}} {\begin{vmatrix} a_1 & b_1 & c_1 \\\\ a_2 & b_2 & c_2 \\\\ a_3 & b_3 & c_3 \end{vmatrix} }.$$

\paragraph{Differential geometry [differential\_geometry]}

\paragraph{Ricci calculus [ricci\_calculus]}

Cramer''s rule is used in the Ricci calculus in various calculations
involving the Christoffel symbols of the first and second kind.

In particular, Cramer''s rule can be used to prove that the divergence
operator on a Riemannian manifold is invariant with respect to change of
coordinates. We give a direct proof, suppressing the role of the
Christoffel symbols. Let $(M,g)$ be a Riemannian manifold equipped with
local coordinates $(x^1, x^2, \dots, x^n)$. Let
$A=A^i \frac{\partial}{\partial x^i}$ be a vector field. We use the
summation convention throughout.

:   \textbf{Theorem}.
:   \textit{The}divergence'''' of $A$,
    $$\operatorname{div} A = \frac{1}{\sqrt{\det g}} \frac{\partial}{\partial x^i} \left( A^i \sqrt{\det g} \right),$$
:   is invariant under change of coordinates.''''

\textbf{\textit{Proof}}

> Let $(x^1,x^2,\ldots,x^n)\mapsto (\bar x^1,\ldots,\bar x^n)$ be a
> coordinate transformation with non-singular Jacobian. Then the
> classical transformation laws imply that
> $A=\bar A^{k}\frac{\partial}{\partial\bar x^{k}}$ where
> $\bar A^{k}=\frac{\partial \bar x^{k}}{\partial x^{j}}A^{j}$.
> Similarly, if
> $g=g_{mk}\,dx^{m}\otimes dx^{k}=\bar{g}_{ij}\,d\bar x^{i}\otimes d\bar x^{j}$,
> then
> $\bar{g}_{ij}=\,\frac{\partial x^{m}}{\partial\bar x^{i}}\frac{\partial x^{k}}{\partial \bar x^{j}}g_{mk}$.
> Writing this transformation law in terms of matrices yields
> $\bar g=\left(\frac{\partial x}{\partial\bar{x}}\right)^{\text{T}}g\left(\frac{\partial x}{\partial\bar{x}}\right)$,
> which implies
> $\det\bar g=\left(\det\left(\frac{\partial x}{\partial\bar{x}}\right)\right)^{2}\det g$.
>
> Now one computes
>
> \begin{align}
> \operatorname{div} A &=\frac{1}{\sqrt{\det g}}\frac{\partial}{\partial x^{i}}\left( A^{i}\sqrt{\det g}\right)\\\>   &=\det\left(\frac{\partial x}{\partial\bar{x}}\right)\frac{1}{\sqrt{\det\bar g}}\frac{\partial \bar x^k}{\partial x^{i}}\frac{\partial}{\partial\bar x^{k}}\left(\frac{\partial x^{i}}{\partial \bar x^{\ell}}\bar{A}^{\ell}\det\!\left(\frac{\partial x}{\partial\bar{x}}\right)^{\!\!-1}\!\sqrt{\det\bar g}\right).
> \end{align} In order to show that this equals
> $\frac{1}{\sqrt{\det\bar g}}\frac{\partial}{\partial\bar x^{k}}\left(\bar A^{k}\sqrt{\det\bar{g}}\right)$,
> it is necessary and sufficient to show that
>
> $$\frac{\partial\bar x^{k}}{\partial x^{i}}\frac{\partial}{\partial\bar x^{k}}\left(\frac{\partial x^{i}}{\partial \bar x^{\ell}}\det\!\left(\frac{\partial x}{\partial\bar{x}}\right)^{\!\!\!-1}\right)=0\qquad\text{for all } \ell,$$
> which is equivalent to
>
> $$\frac{\partial}{\partial \bar x^{\ell}}\det\left(\frac{\partial x}{\partial\bar{x}}\right)
> =\det\left(\frac{\partial x}{\partial\bar{x}}\right)\frac{\partial\bar x^{k}}{\partial x^{i}}\frac{\partial^{2}x^{i}}{\partial\bar x^{k}\partial\bar x^{\ell}}.$$
> Carrying out the differentiation on the left-hand side, we get:
>
> \begin{align}
>   \frac{\partial}{\partial\bar x^{\ell}}\det\left(\frac{\partial x}{\partial\bar{x}}\right)
>   &=(-1)^{i+j}\frac{\partial^{2}x^{i}}{\partial\bar x^{\ell}\partial\bar x^{j}}\det M(i|j)\\\>   &=\frac{\partial^{2}x^{i}}{\partial\bar x^{\ell}\partial\bar x^{j}}\det\left(\frac{\partial x}{\partial\bar{x}}\right)\frac{(-1)^{i+j}}{\det\left(\frac{\partial x}{\partial\bar{x}}\right)}\det M(i|j)=(\ast),
>   \end{align} where $M(i|j)$ denotes the matrix obtained from
> $\left(\frac{\partial x}{\partial\bar{x}}\right)$ by deleting the
> $i$th row and $j$th column. But Cramer''s Rule says that
>
> $$\frac{(-1)^{i+j}}{\det\left(\frac{\partial x}{\partial\bar{x}}\right)}\det M(i|j)$$
> is the $(j,i)$th entry of the matrix
> $\left(\frac{\partial \bar{x}}{\partial x}\right)$. Thus
>
> $$(\ast)=\det\left(\frac{\partial x}{\partial\bar{x}}\right)\frac{\partial^{2}x^{i}}{\partial\bar x^{\ell}\partial\bar x^{j}}\frac{\partial\bar x^{j}}{\partial x^{i}},$$
> completing the proof.

\paragraph{Computing derivatives implicitly [computing\_derivatives\_implicitly]}

Consider the two equations $F(x, y, u, v) = 0$ and $G(x, y, u, v) = 0$.
When \textit{u} and \textit{v} are independent variables, we can define $x = X(u, v)$
and $y = Y(u, v).$

An equation for $\dfrac{\partial x}{\partial u}$ can be found by
applying Cramer''s rule.

\textbf{\textit{Calculation of $\dfrac{\partial x}{\partial u}$}}

> First, calculate the first derivatives of \textit{F}, \textit{G}, \textit{x}, and \textit{y}:
>
> \begin{align}
> dF &= \frac{\partial F}{\partial x} dx + \frac{\partial F}{\partial y} dy +\frac{\partial F}{\partial u} du +\frac{\partial F}{\partial v} dv = 0 \\\\[6pt]
> dG &= \frac{\partial G}{\partial x} dx + \frac{\partial G}{\partial y} dy +\frac{\partial G}{\partial u} du +\frac{\partial G}{\partial v} dv = 0 \\\\[6pt]
> dx &= \frac{\partial X}{\partial u} du + \frac{\partial X}{\partial v} dv \\\\[6pt]
> dy &= \frac{\partial Y}{\partial u} du + \frac{\partial Y}{\partial v} dv.
> \end{align}
>
> Substituting \textit{dx}, \textit{dy} into \textit{dF} and \textit{dG}, we have:
>
> \begin{align}
> dF &= \left(\frac{\partial F}{\partial x} \frac{\partial x}{\partial u} +\frac{\partial F}{\partial y} \frac{\partial y}{\partial u} + \frac{\partial F}{\partial u} \right) du + \left(\frac{\partial F}{\partial x} \frac{\partial x}{\partial v} +\frac{\partial F}{\partial y} \frac{\partial y}{\partial v} +\frac{\partial F}{\partial v} \right) dv = 0 \\\\ [6pt]
> dG &= \left(\frac{\partial G}{\partial x} \frac{\partial x}{\partial u} +\frac{\partial G}{\partial y} \frac{\partial y}{\partial u} +\frac{\partial G}{\partial u} \right) du + \left(\frac{\partial G}{\partial x} \frac{\partial x}{\partial v} +\frac{\partial G}{\partial y} \frac{\partial y}{\partial v} +\frac{\partial G}{\partial v} \right) dv = 0.
> \end{align}
>
> Since \textit{u}, \textit{v} are both independent, the coefficients of \textit{du}, \textit{dv}
> must be zero. So we can write out equations for the coefficients:
>
> \begin{align}
> \frac{\partial F}{\partial x} \frac{\partial x}{\partial u} +\frac{\partial F}{\partial y} \frac{\partial y}{\partial u} & = -\frac{\partial F}{\partial u} \\\\[6pt]
> \frac{\partial G}{\partial x} \frac{\partial x}{\partial u} +\frac{\partial G}{\partial y} \frac{\partial y}{\partial u} & = -\frac{\partial G}{\partial u} \\\\[6pt]
> \frac{\partial F}{\partial x} \frac{\partial x}{\partial v} +\frac{\partial F}{\partial y} \frac{\partial y}{\partial v} & = -\frac{\partial F}{\partial v} \\\\[6pt]
> \frac{\partial G}{\partial x} \frac{\partial x}{\partial v} +\frac{\partial G}{\partial y} \frac{\partial y}{\partial v} & = -\frac{\partial G}{\partial v}.
> \end{align}
>
> Now, by Cramer''s rule, we see that:
>
> $$\frac{\partial x}{\partial u} = \frac{\begin{vmatrix} -\frac{\partial F}{\partial u} & \frac{\partial F}{\partial y} \\\\ -\frac{\partial G}{\partial u} & \frac{\partial G}{\partial y}\end{vmatrix}}{\begin{vmatrix}\frac{\partial F}{\partial x} & \frac{\partial F}{\partial y} \\\\ \frac{\partial G}{\partial x} & \frac{\partial G}{\partial y}\end{vmatrix}}.$$
>
> This is now a formula in terms of two Jacobians:
>
> $$\frac{\partial x}{\partial u} = -\frac{\left(\frac{\partial (F, G)}{\partial (u, y)}\right)}{\left(\frac{\partial (F, G)}{\partial(x, y)}\right)}.$$
>
> Similar formulas can be derived for
> $\frac{\partial x}{\partial v}, \frac{\partial y}{\partial u}, \frac{\partial y}{\partial v}.$

\paragraph{Integer programming [integer\_programming]}

Cramer''s rule can be used to prove that an integer programming problem
whose constraint matrix is totally unimodular and whose right-hand side
is integer, has integer basic solutions. This makes the integer program
substantially easier to solve.

\paragraph{Ordinary differential equations [ordinary\_differential\_equations]}

Cramer''s rule is used to derive the general solution to an inhomogeneous
linear differential equation by the method of variation of parameters.

\paragraph{Geometric interpretation [geometric\_interpretation]}

Cramer''s rule has a geometric interpretation that can be considered also
a proof or simply giving insight about its geometric nature. These
geometric arguments work in general and not only in the case of two
equations with two unknowns presented here.

Given the system of equations

$$\begin{matrix}a_{11}x_1+a_{12}x_2&=b_1\\\\a_{21}x_1+a_{22}x_2&=b_2\end{matrix}$$

it can be considered as an equation between vectors

$$x_1\binom{a_{11}}{a_{21}}+x_2\binom{a_{12}}{a_{22}}=\binom{b_1}{b_2}.$$

The area of the parallelogram determined by $\binom{a_{11}}{a_{21}}$ and
$\binom{a_{12}}{a_{22}}$ is given by the determinant of the system of
equations:

$$\begin{vmatrix}a_{11}&a_{12}\\\\a_{21}&a_{22}\end{vmatrix}.$$

In general, when there are more variables and equations, the determinant
of `{{mvar|n}}`{=mediawiki} vectors of length `{{mvar|n}}`{=mediawiki}
will give the \textit{volume} of the \textit{parallelepiped} determined by those
vectors in the `{{mvar|n}}`{=mediawiki}-th dimensional Euclidean space.

Therefore, the area of the parallelogram determined by
$x_1\binom{a_{11}}{a_{21}}$ and $\binom{a_{12}}{a_{22}}$ has to be $x_1$
times the area of the first one since one of the sides has been
multiplied by this factor. Now, this last parallelogram, by Cavalieri''s
principle, has the same area as the parallelogram determined by
$\binom{b_1}{b_2}=x_1\binom{a_{11}}{a_{21}}+x_2\binom{a_{12}}{a_{22}}$
and $\binom{a_{12}}{a_{22}}.$

Equating the areas of this last and the second parallelogram gives the
equation

$$\begin{vmatrix}b_1&a_{12}\\\\b_2&a_{22}\end{vmatrix} = \begin{vmatrix}a_{11}x_1&a_{12}\\\\a_{21}x_1&a_{22}\end{vmatrix} =x_1 \begin{vmatrix}a_{11}&a_{12}\\\\a_{21}&a_{22}\end{vmatrix}$$

from which Cramer''s rule follows.

\paragraph{Other proofs [other\_proofs]}

\paragraph{A proof by abstract linear algebra [a\_proof\_by\_abstract\_linear\_algebra]}

This is a restatement of the proof above in abstract language.

Consider the map
$\mathbf{x}=(x_1,\ldots, x_n) \mapsto  \frac{1}{\det A} \left(\det (A_1),\ldots, \det(A_n)\right),$
where $A_i$ is the matrix $A$ with $\mathbf{x}$ substituted in the $i$th
column, as in Cramer''s rule. Because of linearity of determinant in
every column, this map is linear. Observe that it sends the $i$th column
of $A$ to the $i$th basis vector $\mathbf{e}_i=(0,\ldots, 1, \ldots, 0)$
(with 1 in the $i$th place), because determinant of a matrix with a
repeated column is 0. So we have a linear map which agrees with the
inverse of $A$ on the column space; hence it agrees with $A^{-1}$ on the
span of the column space. Since $A$ is invertible, the column vectors
span all of $\mathbb{R}^n$, so our map really is the inverse of $A$.
Cramer''s rule follows.

\paragraph{A short proof [a\_short\_proof]}

A short proof of Cramer''s rule can be given by noticing that $x_1$ is
the determinant of the matrix

$$X_1=\begin{bmatrix}
x_1 & 0 & 0 & \cdots & 0\\\x_2 & 1 & 0 & \cdots & 0\\\x_3 & 0 & 1 & \cdots & 0\\\\vdots & \vdots & \vdots & \ddots &\vdots \\\x_n & 0 & 0 & \cdots & 1
\end{bmatrix}$$

On the other hand, assuming that our original matrix
`{{mvar|A}}`{=mediawiki} is invertible, this matrix $X_1$ has columns
$A^{-1}\mathbf{b}, A^{-1}\mathbf{v}_2, \ldots, A^{-1}\mathbf{v}_n$,
where $\mathbf{v}_n$ is the \textit{n}-th column of the matrix
`{{mvar|A}}`{=mediawiki}. Recall that the matrix $A_1$ has columns
$\mathbf{b}, \mathbf{v}_2, \ldots, \mathbf{v}_n$, and therefore
$X_1=A^{-1}A_1$. Hence, by using that the determinant of the product of
two matrices is the product of the determinants, we have

$$x_1= \det (X_1) = \det (A^{-1}) \det (A_1)= \frac{\det (A_1)}{\det (A)}.$$

The proof for other $x_j$ is similar.

\paragraph{Incompatible and indeterminate cases [incompatible\_and\_indeterminate\_cases]}

A system of equations is said to be incompatible or inconsistent when
there are no solutions and it is called indeterminate when there is more
than one solution. For linear equations, an indeterminate system will
have infinitely many solutions (if it is over an infinite field), since
the solutions can be expressed in terms of one or more parameters that
can take arbitrary values.

Cramer''s rule applies to the case where the coefficient determinant is
nonzero. In the 2×2 case, if the coefficient determinant is zero, then
the system is incompatible if the numerator determinants are nonzero, or
indeterminate if the numerator determinants are zero.

For 3×3 or higher systems, the only thing one can say when the
coefficient determinant equals zero is that if any of the numerator
determinants are nonzero, then the system must be incompatible. However,
having all determinants zero does not imply that the system is
indeterminate. A simple example where all determinants vanish (equal
zero) but the system is still incompatible is the 3×3 system
\textit{x}+\textit{y}+\textit{z}=1, \textit{x}+\textit{y}+\textit{z}=2, \textit{x}+\textit{y}+\textit{z}=3.

\paragraph{Licensing [licensing\_14]}

Content obtained and/or adapted from:

\begin{itemize}
    \item \href{https://en.wikipedia.org/wiki/Cramer\%27s\_rule}{Cramer''s rule,
\end{itemize}
    Wikipedia} under a CC
    BY-SA license

%%===============================================================
\section{Diagonalization}
\label{sec:diagonalization}
%%===============================================================

"The prior subsection defines the relation of similarity and shows that,

although similar matrices are necessarily matrix equivalent, the

converse does not hold. Some matrix-equivalence classes break into two

or more similarity classes (the nonsingular $n \! \times \! n$ matrices,

for instance). This means that the canonical form for matrix

equivalence, a block partial-identity, cannot be used as a canonical

form for matrix similarity because the partial-identities cannot be in

more than one similarity class, so there are similarity classes without

one. This picture illustrates. As earlier in this book, class

representatives are shown with stars.











Linalg\_matrix\_similarity\_equiv\_classes\_2.png









We are developing a canonical form for representatives of the similarity

classes. We naturally try to build on our previous work, meaning first

that the partial identity matrices should represent the similarity

classes into which they fall, and beyond that, that the representatives

should be as simple as possible. The simplest extension of the

partial-identity form is a diagonal form.



> \textbf{Definition 2.1}:

>

> :   A transformation is \textbf{diagonalizable} if it has a diagonal

>     representation with respect to the same basis for the codomain as

>     for the domain. A \textbf{diagonalizable matrix} is one that is similar

>     to a diagonal matrix: $T$ is diagonalizable if there is a

>     nonsingular $P$ such that $PTP^{-1}$ is diagonal.



> \textbf{Example 2.2}: The matrix

>

> $$\begin{pmatrix}

> 4 &-2 \\\> 1 &1

> \end{pmatrix}$$

>

> is diagonalizable.

>

> $$\begin{pmatrix}

> 2  &0   \\\> 0  &3

> \end{pmatrix}

> =

> \begin{pmatrix}

> -1  &2  \\\> 1  &-1

> \end{pmatrix}

> \begin{pmatrix}

> 4  &-2 \\\> 1  &1

> \end{pmatrix}

> \begin{pmatrix}

> -1  &2  \\\> 1  &-1

> \end{pmatrix}^{-1}$$



> \textbf{Example 2.3}: Not every matrix is diagonalizable. The square of

>

> $$N=\begin{pmatrix}

> 0  &0  \\\> 1  &0

> \end{pmatrix}$$

>

> is the zero matrix. Thus, for any map $n$ that $N$ represents (with

> respect to the same basis for the domain as for the codomain), the

> composition $n\circ n$ is the zero map. This implies that no such map

> $n$ can be diagonally represented (with respect to any $B,B$) because

> no power of a nonzero diagonal matrix is zero. That is, there is no

> diagonal matrix in $N$'s similarity class.



That example shows that a diagonal form will not do for a canonical

form— we cannot find a diagonal matrix in each matrix similarity class.

However, the canonical form that we are developing has the property that

if a matrix can be diagonalized then the diagonal matrix is the

canonical representative of the similarity class. The next result

characterizes which maps can be diagonalized.



> \textbf{Corollary 2.4}:

>

> :   A transformation $t$ is diagonalizable if and only if there is a

>     basis $B=\langle \vec{\beta}_1,\ldots,\vec{\beta}_n  \rangle$ and

>     scalars $\lambda_1,\ldots,\lambda_n$ such that

>     $t(\vec{\beta}_i)=\lambda_i\vec{\beta}_i$ for each $i$.



Proof: This follows from the definition by considering a diagonal

representation matrix.



$${\rm Rep}_{B,B}(t)=

\left(\begin{array}{c|c|c}

\vdots                    &       &\vdots                     \\\{\rm Rep}_{B}(t(\vec{\beta}_1)) &\cdots &{\rm Rep}_{B}(t(\vec{\beta}_n))  \\\\vdots                    &       &\vdots

\end{array}\right)

=

\left(\begin{array}{c|c|c}

\lambda_1   &       &0         \\\\vdots      &\ddots &\vdots    \\\0           &       &\lambda_n

\end{array}\right)$$



This representation is equivalent to the existence of a basis satisfying

the stated conditions simply by the definition of matrix representation.



> Example 2.5: To diagonalize

>

> $$T=\begin{pmatrix}

> 3  &2  \\\> 0  &1

> \end{pmatrix}$$

>

> we take it as the representation of a transformation with respect to

> the standard basis $T={\rm Rep}_{\mathcal{E}_2,\mathcal{E}_2}(t)$ and

> we look for a basis $B=\langle \vec{\beta}_1,\vec{\beta}_2 \rangle$

> such that

>

> $${\rm Rep}_{B,B}(t)

> =

> \begin{pmatrix}

> \lambda_1  &0          \\\> 0          &\lambda_2

> \end{pmatrix}$$

>

> that is, such that $t(\vec{\beta}_1)=\lambda_1\vec{\beta}_1$ and

> $t(\vec{\beta}_2)=\lambda_2\vec{\beta}_2$.

>

> $$\begin{pmatrix}

> 3  &2  \\\> 0  &1

> \end{pmatrix}

> \vec{\beta}_1=\lambda_1\cdot\vec{\beta}_1

> \qquad

> \begin{pmatrix}

> 3  &2  \\\> 0  &1

> \end{pmatrix}

> \vec{\beta}_2=\lambda_2\cdot\vec{\beta}_2$$

>

> We are looking for scalars $x$ such that this equation

>

> $$\begin{pmatrix}

> 3  &2  \\\> 0  &1

> \end{pmatrix}

> \begin{pmatrix} b_1 \\\\ b_2 \end{pmatrix}=x\cdot\begin{pmatrix} b_1 \\\\ b_2 \end{pmatrix}$$

>

> has solutions $b_1$ and $b_2$, which are not both zero. Rewrite that

> as a linear system.

>

> $$\begin{array}{*{2}{rc}r}

> (3-x)\cdot b_1  &+  &2\cdot b_2       &=  &0  \\\> &   &(1-x)\cdot b_2   &=  &0

> \end{array}

> \qquad (*)$$

>

> In the bottom equation the two numbers multiply to give zero only if

> at least one of them is zero so there are two possibilities, $b_2=0$

> and $x=1$. In the $b_2=0$ possibility, the first equation gives that

> either $b_1=0$ or $x=3$. Since the case of both $b_1=0$ and $b_2=0$ is

> disallowed, we are left looking at the possibility of $x=3$. With it,

> the first equation in ($*$) is $0\cdot b_1+2\cdot b_2=0$ and so

> associated with $3$ are vectors with a second component of zero and a

> first component that is free.

>

> $$\begin{pmatrix}

> 3  &2  \\\> 0  &1

> \end{pmatrix}

> \begin{pmatrix} b_1 \\\\ 0 \end{pmatrix}=3\cdot\begin{pmatrix} b_1 \\\\ 0 \end{pmatrix}$$

>

> That is, one solution to ($*$) is $\lambda_1=3$, and we have a first

> basis vector.

>

> $$\vec{\beta}_1=\begin{pmatrix} 1 \\\\ 0 \end{pmatrix}$$

>

> In the $x=1$ possibility, the first equation in ($*$) is

> $2\cdot b_1+2\cdot b_2=0$, and so associated with $1$ are vectors

> whose second component is the negative of their first component.

>

> $$\begin{pmatrix}

> 3  &2  \\\> 0  &1

> \end{pmatrix}

> \begin{pmatrix} b_1 \\\\ -b_1 \end{pmatrix}=1\cdot\begin{pmatrix} b_1 \\\\ -b_1 \end{pmatrix}$$

>

> Thus, another solution is $\lambda_2=1$ and a second basis vector is

> this.

>

> $$\vec{\beta}_2=\begin{pmatrix} 1 \\\\ -1 \end{pmatrix}$$

>

> To finish, drawing the similarity diagram

>

> :   

>          title=""Linalg\_matrix\_equivalent\_cd\_3.png"" height=""150"" />

>     Linalg\_matrix\_equivalent\_cd\_3.png

>     

>

> and noting that the matrix ${\rm Rep}_{B,\mathcal{E}_2}(\mbox{id})$ is

> easy leads to this diagonalization.

>

> $$\begin{pmatrix}

> 3  &0  \\\> 0  &1

> \end{pmatrix}

> =

> \begin{pmatrix}

> 1  &1  \\\> 0  &-1

> \end{pmatrix}^{-1}

> \begin{pmatrix}

> 3  &2  \\\> 0  &1

> \end{pmatrix}

> \begin{pmatrix}

> 1  &1  \\\> 0  &-1

> \end{pmatrix}$$



In the next subsection, we will expand on that example by considering

more closely the property of \href{Corollary

2.4} . This includes seeing

another way, the way that we will routinely use, to find the

$\lambda$'s.



\paragraph{Resources [resources\_5]}



\begin{itemize}
    \item \href{https://en.wikibooks.org/wiki/Linear\_Algebra/Diagonalizability}{Diagonalizability},
\end{itemize}
    Wikibooks: Linear Algebra



\paragraph{Licensing [licensing\_16]}



Content obtained and/or adapted from:



\begin{itemize}
    \item \href{https://en.wikibooks.org/wiki/Linear\_Algebra/Diagonalizability}{Diagonalizability, Wikibooks: Linear
\end{itemize}
    Algebra}

    under a CC BY-SA license"

%%===============================================================
\section{Finding Eigenvalues and Eigenvectors}
\label{sec:finding-eigenvalues-and-eigenvectors}
%%===============================================================

Definition



In linear algebra, an eigenvector of a matrix is a nonzero vector that

changes at most by a scalar factor when that linear transformation is

applied to it. The corresponding eigenvalue, often denoted by $\lambda$,

is the factor by which the eigenvector is scaled. That is, given some

eigenvector $v_i$ of a square matrix $\mathbf{A}$,

$\mathbf{A}v_i = \lambda_i v_i$, where $\lambda_i$ is the corresponding

eigenvalue of $v_i$. For example:



Let $\mathbf{A} = \begin{bmatrix}

3 \& 4 \& -2\\\1 \& 4 \& -1\\\2 \& 6 \& -1

\end{bmatrix}$, $v_1 = \begin{bmatrix}

1\\\1\\\2

\end{bmatrix}$



$\mathbf{A}v_1 = \begin{bmatrix}

3 \& 4 \& -2\\\1 \& 4 \& -1\\\2 \& 6 \& -1

\end{bmatrix} \begin{bmatrix}

1\\\1\\\2

\end{bmatrix} = \begin{bmatrix}

3\\\3\\\6

\end{bmatrix} = 3\begin{bmatrix}

1\\\1\\\2

\end{bmatrix} = 3v_1$



Thus, $v_1$ is an eigenvector of matrix $\mathbf{A}$, and its

corresponding eigenvalue $\lambda_1 = 3$.



Geometrically, an eigenvector, corresponding to a real nonzero

eigenvalue, points in a direction in which it is stretched by the

transformation and the eigenvalue is the factor by which it is

stretched. If the eigenvalue is negative, the direction is reversed.

Loosely speaking, in a multidimensional vector space, the eigenvector is

not rotated.



\paragraph{Formal definition [formal\_definition]}



If `{{mvar|T}}`{=mediawiki} is a linear transformation from a vector

space `{{mvar|V}}`{=mediawiki} over a field `{{mvar|F}}`{=mediawiki}

into itself and `{{math|'''v'''}}`{=mediawiki} is a nonzero vector in

`{{mvar|V}}`{=mediawiki}, then `{{math|'''v'''}}`{=mediawiki} is an

eigenvector of `{{mvar|T}}`{=mediawiki} if

`{{math|''T''('''v''')}}`{=mediawiki} is a scalar multiple of

`{{math|'''v'''}}`{=mediawiki}. This can be written as



$$T(\mathbf{v}) = \lambda \mathbf{v},$$



where `{{mvar|?}}`{=mediawiki} is a scalar in `{{mvar|F}}`{=mediawiki},

known as the \textbf{eigenvalue}, \textbf{characteristic value}, or

\textbf{characteristic root} associated with `{{math|'''v'''}}`{=mediawiki}.



There is a direct correspondence between \textit{n}-by-\textit{n} square matrices and

linear transformations from an \textit{n}-dimensional vector space into itself,

given any basis of the vector space. Hence, in a finite-dimensional

vector space, it is equivalent to define eigenvalues and eigenvectors

using either the language of matrices, or the language of linear

transformations.



If `{{mvar|V}}`{=mediawiki} is finite-dimensional, the above equation is

equivalent to



$$A\mathbf{u} = \lambda \mathbf{u}.$$



where `{{mvar|A}}`{=mediawiki} is the matrix representation of

`{{mvar|T}}`{=mediawiki} and `{{math|'''u'''}}`{=mediawiki} is the

coordinate vector of `{{math|'''v'''}}`{=mediawiki}.



\paragraph{Overview [overview\_1]}



Eigenvalues and eigenvectors feature prominently in the analysis of

linear transformations. The prefix \textit{eigen-} is adopted from the German

word \textit{eigen} (cognate with the English word \textit{own}) for ""proper"",

""characteristic"", ""own"". Originally used to study principal axes of the

rotational motion of rigid bodies, eigenvalues and eigenvectors have a

wide range of applications, for example in stability analysis, vibration

analysis, atomic orbitals, facial recognition, and matrix

diagonalization.



In essence, an eigenvector \textbf{v} of a linear transformation \textit{T} is a

nonzero vector that, when \textit{T} is applied to it, does not change

direction. Applying \textit{T} to the eigenvector only scales the eigenvector

by the scalar value \textit{?}, called an eigenvalue. This condition can be

written as the equation



$$T(\mathbf{v}) = \lambda \mathbf{v},$$



referred to as the \textbf{eigenvalue equation} or \textbf{eigenequation}. In

general, \textit{?} may be any scalar. For example, \textit{?} may be negative, in

which case the eigenvector reverses direction as part of the scaling, or

it may be zero or complex.









The Mona Lisa example pictured here provides a simple illustration. Each

point on the painting can be represented as a vector pointing from the

center of the painting to that point. The linear transformation in this

example is called a shear mapping. Points in the top half are moved to

the right, and points in the bottom half are moved to the left,

proportional to how far they are from the horizontal axis that goes

through the middle of the painting. The vectors pointing to each point

in the original image are therefore tilted right or left, and made

longer or shorter by the transformation. Points \textit{along} the horizontal

axis do not move at all when this transformation is applied. Therefore,

any vector that points directly to the right or left with no vertical

component is an eigenvector of this transformation, because the mapping

does not change its direction. Moreover, these eigenvectors all have an

eigenvalue equal to one, because the mapping does not change their

length either.



Linear transformations can take many different forms, mapping vectors in

a variety of vector spaces, so the eigenvectors can also take many

forms. For example, the linear transformation could be a differential

operator like $\tfrac{d}{dx}$, in which case the eigenvectors are

functions called eigenfunctions that are scaled by that differential

operator, such as



$$\frac{d}{dx}e^{\lambda x} = \lambda e^{\lambda x}.$$



Alternatively, the linear transformation could take the form of an \textit{n}

by \textit{n} matrix, in which case the eigenvectors are \textit{n} by 1 matrices. If

the linear transformation is expressed in the form of an \textit{n} by \textit{n}

matrix \textit{A}, then the eigenvalue equation for a linear transformation

above can be rewritten as the matrix multiplication



$$A\mathbf v = \lambda \mathbf v,$$



where the eigenvector \textit{v} is an \textit{n} by 1 matrix. For a matrix,

eigenvalues and eigenvectors can be used to decompose the matrix—for

example by diagonalizing it.



Eigenvalues and eigenvectors give rise to many closely related

mathematical concepts, and the prefix \textit{eigen-} is applied liberally when

naming them:



\begin{itemize}
    \item The set of all eigenvectors of a linear transformation, each paired
\end{itemize}
    with its corresponding eigenvalue, is called the \textbf{eigensystem} of

    that transformation.

\begin{itemize}
    \item The set of all eigenvectors of \textit{T} corresponding to the same
\end{itemize}
    eigenvalue, together with the zero vector, is called an

    \textbf{eigenspace}, or the \textbf{characteristic space} of \textit{T} associated

    with that eigenvalue.

\begin{itemize}
    \item If a set of eigenvectors of \textit{T} forms a basis of the domain of \textit{T},
\end{itemize}
    then this basis is called an \textbf{eigenbasis}.



\paragraph{Eigenvalues and eigenvectors of matrices [eigenvalues\_and\_eigenvectors\_of\_matrices]}



Eigenvalues and eigenvectors are often introduced to students in the

context of linear algebra courses focused on matrices. Furthermore,

linear transformations over a finite-dimensional vector space can be

represented using matrices, which is especially common in numerical and

computational applications.







Matrix A acts by stretching the vector

x, not changing its direction, so x is

an eigenvector of A.





Consider `{{mvar|n}}`{=mediawiki}-dimensional vectors that are formed as

a list of `{{mvar|n}}`{=mediawiki} scalars, such as the

three-dimensional vectors



$$\mathbf x = \begin{bmatrix}1\\\\-3\\\\4\end{bmatrix}\quad\mbox{and}\quad \mathbf y = \begin{bmatrix}-20\\\\60\\\\-80\end{bmatrix}.$$



These vectors are said to be scalar multiples of each other, or parallel

or collinear, if there is a scalar `{{mvar|?}}`{=mediawiki} such that



$$\mathbf x = \lambda \mathbf y.$$



In this case $\lambda = -\frac{1}{20}$.



Now consider the linear transformation of

`{{mvar|n}}`{=mediawiki}-dimensional vectors defined by an

`{{mvar|n}}`{=mediawiki} by `{{mvar|n}}`{=mediawiki} matrix

`{{mvar|A}}`{=mediawiki},



$$A \mathbf v = \mathbf w,$$



or



$$\begin{bmatrix}

    A_{11} & A_{12} & \cdots & A_{1n} \\\    A_{21} & A_{22} & \cdots & A_{2n} \\\    \vdots & \vdots & \ddots & \vdots \\\    A_{n1} & A_{n2} & \cdots & A_{nn} \\\  \end{bmatrix}\begin{bmatrix}

    v_1 \\\\ v_2 \\\\ \vdots \\\\ v_n

  \end{bmatrix} = \begin{bmatrix}

    w_1 \\\\ w_2 \\\\ \vdots \\\\ w_n

  \end{bmatrix}$$



where, for each row,



$$w_i = A_{i1} v_1 + A_{i2} v_2 + \cdots + A_{in} v_n = \sum_{j = 1}^n A_{ij} v_j.$$



If it occurs that `{{mvar|v}}`{=mediawiki} and `{{mvar|w}}`{=mediawiki}

are scalar multiples, that is if



$$A \mathbf v = \mathbf w = \lambda \mathbf v,$$ (\textbf{1}) then

`{{math|'''v'''}}`{=mediawiki} is an \textbf{eigenvector} of the linear

transformation `{{mvar|A}}`{=mediawiki} and the scale factor

`{{mvar|?}}`{=mediawiki} is the \textbf{eigenvalue} corresponding to that

eigenvector. Equation (\textbf{1}) is the \textbf{eigenvalue equation} for the

matrix `{{mvar|A}}`{=mediawiki}.



Equation (\textbf{1}) can be stated equivalently as



$$\left(A - \lambda I \right) \mathbf v = \mathbf 0,$$ (\textbf{2})



where `{{mvar|I}}`{=mediawiki} is the `{{mvar|n}}`{=mediawiki} by

`{{mvar|n}}`{=mediawiki} identity matrix and \textbf{0} is the zero vector.



\paragraph{Eigenvalues and the characteristic polynomial [eigenvalues\_and\_the\_characteristic\_polynomial]}



Equation (\textbf{2}) has a nonzero solution \textit{v} if and only if the

determinant of the matrix (\textit{A} - \textit{?I}) is zero. Therefore, the

eigenvalues of \textit{A} are values of \textit{?} that satisfy the equation



$$|A-\lambda I| = 0$$ (\textbf{3})



Using Leibniz' rule for the determinant, the left-hand side of Equation

(\textbf{3}) is a polynomial function of the variable $\lambda$ and the degree of

this polynomial is $n$, the order of the matrix $A$. Its coefficients

depend on the entries of $A$, except that its term of degree $n$ is

always $(-1)^n \lambda^n$. This polynomial is called the characteristic

polynomial of $A$. Equation (\textbf{3}) is called the characteristic

equation or the secular equation of $A$.



The fundamental theorem of algebra implies that the characteristic

polynomial of an \textit{n}-by-\textit{n} matrix \textit{A}, being a polynomial of degree

\textit{n}, can be factored into the product of \textit{n} linear terms,



$$|A - \lambda I| = (\lambda_1 - \lambda )(\lambda_2 - \lambda) \cdots (\lambda_n - \lambda),$$

(\textbf{4})



where each \textit{?}\textasciitilde{}\textit{i}\textasciitilde{} may be real but in general is a complex number. The

numbers \textit{?}\textasciitilde{}1\textasciitilde{}, \textit{?}\textasciitilde{}2\textasciitilde{}, …, \textit{?}\textasciitilde{}\textit{n}\textasciitilde{}, which may not all have distinct

values, are roots of the polynomial and are the eigenvalues of \textit{A}.



As a brief example, which is described in more detail in the examples

section later, consider the matrix



$$A = \begin{bmatrix}

  2 & 1\\\  1 & 2

\end{bmatrix}.$$



Taking the determinant of (\textit{A} - \textit{?I}), the characteristic polynomial of

\textit{A} is



$$|A - \lambda I| = \begin{vmatrix}

    2 - \lambda & 1 \\\    1           & 2 - \lambda

  \end{vmatrix} =

  3 - 4\lambda + \lambda^2.$$



Setting the characteristic polynomial equal to zero, it has roots at ?=1

and ?=3, which are the two eigenvalues of \textit{A}. The eigenvectors

corresponding to each eigenvalue can be found by solving for the

components of \textbf{v} in the equation

$\left(A - \lambda I\right) \mathbf v = \mathbf 0$. In this example, the

eigenvectors are any nonzero scalar multiples of



$$\mathbf v_{\lambda=1} = \begin{bmatrix} 1 \\\\ -1 \end{bmatrix}, \quad \mathbf v_{\lambda=3} = \begin{bmatrix} 1 \\\\ 1 \end{bmatrix}.$$



If the entries of the matrix \textit{A} are all real numbers, then the

coefficients of the characteristic polynomial will also be real numbers,

but the eigenvalues may still have nonzero imaginary parts. The entries

of the corresponding eigenvectors therefore may also have nonzero

imaginary parts. Similarly, the eigenvalues may be irrational numbers

even if all the entries of \textit{A} are rational numbers or even if they are

all integers. However, if the entries of \textit{A} are all algebraic numbers,

which include the rationals, the eigenvalues are complex algebraic

numbers.



The non-real roots of a real polynomial with real coefficients can be

grouped into pairs of complex conjugates, namely with the two members of

each pair having imaginary parts that differ only in sign and the same

real part. If the degree is odd, then by the intermediate value theorem

at least one of the roots is real. Therefore, any real matrix with odd

order has at least one real eigenvalue, whereas a real matrix with even

order may not have any real eigenvalues. The eigenvectors associated

with these complex eigenvalues are also complex and also appear in

complex conjugate pairs.



\paragraph{Algebraic multiplicity [algebraic\_multiplicity]}



Let $\lambda_i$ be an eigenvalue of an $n$ by $n$ matrix $A$. The

\textbf{algebraic multiplicity} $\mu_A(\lambda_i)$ of the eigenvalue is its

multiplicity as a root of the characteristic polynomial, that is, the

largest integer $k$ such that $(\lambda - \lambda_i)^k$ divides evenly that

polynomial.



Suppose a matrix \textit{A} has dimension \textit{n} and \textit{d} = \textit{n} distinct

eigenvalues. Whereas Equation (\textbf{4}) factors the characteristic

polynomial of \textit{A} into the product of \textit{n} linear terms with some terms

potentially repeating, the characteristic polynomial can instead be

written as the product of \textit{d} terms each corresponding to a distinct

eigenvalue and raised to the power of the algebraic multiplicity,



$$|A - \lambda I| = (\lambda_1 - \lambda)^{\mu_A(\lambda_1)}(\lambda_2 - \lambda)^{\mu_A(\lambda_2)} \cdots (\lambda_d - \lambda)^{\mu_A(\lambda_d)}.$$



If \textit{d} = \textit{n} then the right-hand side is the product of \textit{n} linear terms

and this is the same as Equation (\textbf{4}). The size of each eigenvalue's

algebraic multiplicity is related to the dimension \textit{n} as



\begin{align}
      1 &\leq \mu_A(\lambda_i) \leq n, \\\  \mu_A &= \sum_{i=1}^d \mu_A\left(\lambda_i\right) = n.
\end{align}



If \textit{µ}\textasciitilde{}\textit{A}\textasciitilde{}(\textit{?}\textasciitilde{}\textit{i}\textasciitilde{}) = 1, then \textit{?}\textasciitilde{}\textit{i}\textasciitilde{} is said to be a *simple

eigenvalue\textit{. If }µ\textit{\textasciitilde{}}A\textit{\textasciitilde{}(}?\textit{\textasciitilde{}}i*\textasciitilde{}) equals the geometric multiplicity of

\textit{?}\textasciitilde{}\textit{i}\textasciitilde{}, \textit{?}\textasciitilde{}\textit{A}\textasciitilde{}(\textit{?}\textasciitilde{}\textit{i}\textasciitilde{}), defined in the next section, then \textit{?}\textasciitilde{}\textit{i}\textasciitilde{}

is said to be a \textit{semisimple eigenvalue}.



\paragraph{Eigenspaces, geometric multiplicity, and the eigenbasis for matrices [eigenspaces\_geometric\_multiplicity\_and\_the\_eigenbasis\_for\_matrices]}



Given a particular eigenvalue \textit{?} of the \textit{n} by \textit{n} matrix \textit{A}, define

the set \textit{E} to be all vectors \textbf{v} that satisfy Equation (\textbf{2}),



$$E = \left\{\mathbf{v} : \left(A - \lambda I\right) \mathbf{v} = \mathbf{0}\right\}.$$



On one hand, this set is precisely the kernel or nullspace of the matrix

(\textit{A} - \textit{?I}). On the other hand, by definition, any nonzero vector that

satisfies this condition is an eigenvector of \textit{A} associated with \textit{?}.

So, the set \textit{E} is the union of the zero vector with the set of all

eigenvectors of \textit{A} associated with \textit{?}, and \textit{E} equals the nullspace of

(\textit{A} - \textit{?I}). \textit{E} is called the \textbf{eigenspace} or **characteristic

space** of \textit{A} associated with \textit{?}. In general \textit{?} is a complex number

and the eigenvectors are complex \textit{n} by 1 matrices. A property of the

nullspace is that it is a linear subspace, so \textit{E} is a linear subspace

of $\mathbb{C}^n$.



Because the eigenspace \textit{E} is a linear subspace, it is closed under

addition. That is, if two vectors \textbf{u} and \textbf{v} belong to the set \textit{E},

written \textbf{u}, \textbf{v} ? \textit{E}, then (\textbf{u} + \textbf{v}) ? \textit{E} or equivalently

\textit{A}(\textbf{u} + \textbf{v}) = \textit{?}(\textbf{u} + \textbf{v}). This can be checked using the

distributive property of matrix multiplication. Similarly, because \textit{E}

is a linear subspace, it is closed under scalar multiplication. That is,

if \textbf{v} ? \textit{E} and \textit{a} is a complex number, (\textit{a\textbf{}v}) ? \textit{E} or

equivalently \textit{A}(\textit{a\textbf{}v}) = \textit{?}(\textit{a\textbf{}v}). This can be checked by

noting that multiplication of complex matrices by complex numbers is

commutative. As long as \textbf{u} + \textbf{v} and \textit{a\textbf{}v} are not zero, they

are also eigenvectors of \textit{A} associated with \textit{?}.



The dimension of the eigenspace \textit{E} associated with \textit{?}, or equivalently

the maximum number of linearly independent eigenvectors associated with

\textit{?}, is referred to as the eigenvalue's \textbf{geometric multiplicity}

\textit{?}\textasciitilde{}\textit{A}\textasciitilde{}(\textit{?}). Because \textit{E} is also the nullspace of (\textit{A} - \textit{?I}), the

geometric multiplicity of \textit{?} is the dimension of the nullspace of (\textit{A}

\begin{itemize}
    \item \textit{?I}), also called the \textit{nullity} of (\textit{A} - \textit{?I}), which relates to the
\end{itemize}
dimension and rank of (\textit{A} - \textit{?I}) as



$$\gamma_A(\lambda) = n - \operatorname{rank}(A - \lambda I).$$



Because of the definition of eigenvalues and eigenvectors, an

eigenvalue's geometric multiplicity must be at least one, that is, each

eigenvalue has at least one associated eigenvector. Furthermore, an

eigenvalue's geometric multiplicity cannot exceed its algebraic

multiplicity. Additionally, recall that an eigenvalue's algebraic

multiplicity cannot exceed \textit{n}.



$$1 \le \gamma_A(\lambda) \le \mu_A(\lambda) \le n$$



To prove the inequality $\gamma_A(\lambda)\le\mu_A(\lambda)$, consider

how the definition of geometric multiplicity implies the existence of

$\gamma_A(\lambda)$ orthonormal eigenvectors

$\boldsymbol{v}_1,\, \ldots,\, \boldsymbol{v}_{\gamma_A(\lambda)}$, such

that $A \boldsymbol{v}_k = \lambda \boldsymbol{v}_k$. We can therefore

find a (unitary) matrix $V$ whose first $\gamma_A(\lambda)$ columns are

these eigenvectors, and whose remaining columns can be any orthonormal

set of $n - \gamma_A(\lambda)$ vectors orthogonal to these eigenvectors

of $A$. Then $V$ has full rank and is therefore invertible, and $AV=VD$

with $D$ a matrix whose top left block is the diagonal matrix

$\lambda I_{\gamma_A(\lambda)}$. This implies that

$(A - \xi I)V = V(D - \xi I)$. In other words, $A - \xi I$ is similar to

$D - \xi I$, which implies that $\det(A - \xi I) = \det(D - \xi I)$. But

from the definition of $D$ we know that $\det(D - \xi I)$ contains a

factor $(\xi - \lambda)^{\gamma_A(\lambda)}$, which means that the

algebraic multiplicity of $\lambda$ must satisfy

$\mu_A(\lambda) \ge \gamma_A(\lambda)$.



Suppose $A$ has $d \leq n$ distinct eigenvalues

$\lambda_1, \ldots, \lambda_d$, where the geometric multiplicity of

$\lambda_i$ is $\gamma_A (\lambda_i)$. The total geometric multiplicity

of $A$,



\begin{align}
  \gamma_A &= \sum_{i=1}^d \gamma_A(\lambda_i), \\\         d &\le \gamma_A \le n,
\end{align}



is the dimension of the sum of all the eigenspaces of $A$'s eigenvalues,

or equivalently the maximum number of linearly independent eigenvectors

of $A$. If $\gamma_A=n$, then



\begin{itemize}
    \item The direct sum of the eigenspaces of all of $A$'s eigenvalues is the
\end{itemize}
    entire vector space $\mathbb{C}^n$.

\begin{itemize}
    \item A basis of $\mathbb{C}^n$ can be formed from $n$ linearly
\end{itemize}
    independent eigenvectors of $A$; such a basis is called an

    \textbf{eigenbasis}

\begin{itemize}
    \item Any vector in $\mathbb{C}^n$ can be written as a linear combination
\end{itemize}
    of eigenvectors of $A$.



\paragraph{Additional properties of eigenvalues [additional\_properties\_of\_eigenvalues]}



Let $A$ be an arbitrary $n \times n$ matrix of complex numbers with

eigenvalues $\lambda_1, \ldots, \lambda_n$. Each eigenvalue appears

$\mu_A(\lambda_i)$ times in this list, where $\mu_A(\lambda_i)$ is the

eigenvalue's algebraic multiplicity. The following are properties of

this matrix and its eigenvalues:



\begin{itemize}
    \item The trace of $A$, defined as the sum of its diagonal elements, is
\end{itemize}
    also the sum of all eigenvalues,



    :   $\operatorname{tr}(A) = \sum_{i=1}^n a_{ii} = \sum_{i=1}^n \lambda_i = \lambda_1 + \lambda_2 + \cdots + \lambda_n.$

\begin{itemize}
    \item The determinant of $A$ is the product of all its eigenvalues,
\end{itemize}

    :   $\det(A) = \prod_{i=1}^n \lambda_i = \lambda_1\lambda_2 \cdots \lambda_n.$

\begin{itemize}
    \item The eigenvalues of the $k$^th^ power of $A$; i.e., the eigenvalues
\end{itemize}
    of $A^k$, for any positive integer $k$, are

    $\lambda_1^k, \ldots, \lambda_n^k$.

\begin{itemize}
    \item The matrix $A$ is invertible if and only if every eigenvalue is
\end{itemize}
    nonzero.

\begin{itemize}
    \item If $A$ is invertible, then the eigenvalues of $A^{-1}$ are
\end{itemize}
    $\frac{1}{\lambda_1}, \ldots, \frac{1}{\lambda_n}$ and each

    eigenvalue's geometric multiplicity coincides. Moreover, since the

    characteristic polynomial of the inverse is the reciprocal

    polynomial of the original, the eigenvalues share the same algebraic

    multiplicity.

\begin{itemize}
    \item If $A$ is equal to its conjugate transpose $A^*$, or equivalently if
\end{itemize}
    $A$ is Hermitian, then every eigenvalue is real. The same is true of

    any symmetric real matrix.

\begin{itemize}
    \item If $A$ is not only Hermitian but also positive-definite,
\end{itemize}
    positive-semidefinite, negative-definite, or negative-semidefinite,

    then every eigenvalue is positive, non-negative, negative, or

    non-positive, respectively.

\begin{itemize}
    \item If $A$ is unitary, every eigenvalue has absolute value
\end{itemize}
    $|\lambda_i|=1$.

\begin{itemize}
    \item if $A$ is a $n\times n$ matrix and $\{\lambda_1,\ldots,\lambda_k\}$
\end{itemize}
    are its eigenvalues, then the eigenvalues of matrix $I+A$ (where $I$

    is the identity matrix) are $\{\lambda_1+1,\ldots,\lambda_k+1\}$.

    Moreover, if $\alpha\in\mathbb C$, the eigenvalues of $\alpha I+A$

    are $\{\lambda_1+\alpha,\ldots,\lambda_k+\alpha\}$. More generally,

    for a polynomial $P$ the eigenvalues of matrix $P(A)$ are

    $\{P(\lambda_1), \ldots, P(\lambda_k)\}$.



\paragraph{Left and right eigenvectors [left\_and\_right\_eigenvectors]}



Many disciplines traditionally represent vectors as matrices with a

single column rather than as matrices with a single row. For that

reason, the word ""eigenvector"" in the context of matrices almost always

refers to a \textbf{right eigenvector}, namely a \textit{column} vector that \textit{right}

multiplies the $n \times n$ matrix $A$ in the defining equation,

Equation (\textbf{1}),



$$A \mathbf v = \lambda \mathbf v.$$



The eigenvalue and eigenvector problem can also be defined for \textit{row}

vectors that \textit{left} multiply matrix $A$. In this formulation, the

defining equation is



$$\mathbf u A = \kappa \mathbf u,$$



where $\kappa$ is a scalar and $u$ is a $1 \times n$ matrix. Any row

vector $u$ satisfying this equation is called a \textbf{left eigenvector} of

$A$ and $\kappa$ is its associated eigenvalue. Taking the transpose of

this equation,



$$A^\textsf{T} \mathbf u^\textsf{T} = \kappa \mathbf u^\textsf{T}.$$



Comparing this equation to Equation (\textbf{1}), it follows immediately that

a left eigenvector of $A$ is the same as the transpose of a right

eigenvector of $A^\textsf{T}$, with the same eigenvalue. Furthermore,

since the characteristic polynomial of $A^\textsf{T}$ is the same as the

characteristic polynomial of $A$, the eigenvalues of the left

eigenvectors of $A$ are the same as the eigenvalues of the right

eigenvectors of $A^\textsf{T}$.



\paragraph{Diagonalization and the eigendecomposition [diagonalization\_and\_the\_eigendecomposition]}



Suppose the eigenvectors of \textit{A} form a basis, or equivalently \textit{A} has

\textit{n} linearly independent eigenvectors \textbf{v}\textasciitilde{}1\textasciitilde{}, \textbf{v}\textasciitilde{}2\textasciitilde{}, …, \textbf{v}\textasciitilde{}\textit{n}\textasciitilde{}

with associated eigenvalues \textit{?}\textasciitilde{}1\textasciitilde{}, \textit{?}\textasciitilde{}2\textasciitilde{}, …, \textit{?}\textasciitilde{}\textit{n}\textasciitilde{}. The eigenvalues

need not be distinct. Define a square matrix \textit{Q} whose columns are the

\textit{n} linearly independent eigenvectors of \textit{A},



:   $Q = \begin{bmatrix} \mathbf v_1 & \mathbf v_2 & \cdots & \mathbf v_n \end{bmatrix}.$



Since each column of \textit{Q} is an eigenvector of \textit{A}, right multiplying \textit{A}

by \textit{Q} scales each column of \textit{Q} by its associated eigenvalue,



:   $AQ = \begin{bmatrix} \lambda_1 \mathbf v_1 & \lambda_2 \mathbf v_2 & \cdots & \lambda_n \mathbf v_n \end{bmatrix}.$



With this in mind, define a diagonal matrix ? where each diagonal

element ?\textasciitilde{}\textit{ii}\textasciitilde{} is the eigenvalue associated with the \textit{i}th column of

\textit{Q}. Then



:   $AQ = Q\Lambda.$



Because the columns of \textit{Q} are linearly independent, Q is invertible.

Right multiplying both sides of the equation by \textit{Q}^-1^,



:   $A = Q\Lambda Q^{-1},$



or by instead left multiplying both sides by \textit{Q}^-1^,



:   $Q^{-1}AQ = \Lambda.$



\textit{A} can therefore be decomposed into a matrix composed of its

eigenvectors, a diagonal matrix with its eigenvalues along the diagonal,

and the inverse of the matrix of eigenvectors. This is called the

eigendecomposition and it is a similarity transformation. Such a matrix

\textit{A} is said to be \textit{similar} to the diagonal matrix ? or

\textit{diagonalizable}. The matrix \textit{Q} is the change of basis matrix of the

similarity transformation. Essentially, the matrices \textit{A} and ? represent

the same linear transformation expressed in two different bases. The

eigenvectors are used as the basis when representing the linear

transformation as ?.



Conversely, suppose a matrix \textit{A} is diagonalizable. Let \textit{P} be a

non-singular square matrix such that $P^{-1}AP$ is some diagonal matrix

\textit{D}. Left multiplying both by \textit{P}, \textit{AP} = \textit{PD}. Each column of \textit{P} must

therefore be an eigenvector of \textit{A} whose eigenvalue is the corresponding

diagonal element of \textit{D}. Since the columns of \textit{P} must be linearly

independent for \textit{P} to be invertible, there exist \textit{n} linearly

independent eigenvectors of \textit{A}. It then follows that the eigenvectors

of \textit{A} form a basis if and only if \textit{A} is diagonalizable.



A matrix that is not diagonalizable is said to be defective. For

defective matrices, the notion of eigenvectors generalizes to

generalized eigenvectors and the diagonal matrix of eigenvalues

generalizes to the Jordan normal form. Over an algebraically closed

field, any matrix \textit{A} has a Jordan normal form and therefore admits a

basis of generalized eigenvectors and a decomposition into generalized

eigenspaces.



\paragraph{Variational characterization [variational\_characterization]}



In the Hermitian case, eigenvalues can be given a variational

characterization. The largest eigenvalue of $H$ is the maximum value of

the quadratic form

$\mathbf x^\textsf{T} H \mathbf x / \mathbf x^\textsf{T} \mathbf x$. A

value of $\mathbf x$ that realizes that maximum, is an eigenvector.



\paragraph{Matrix examples [matrix\_examples]}



\paragraph{Two-dimensional matrix example [two\_dimensional\_matrix\_example]}



![The transformation matrix \textit{A} =

$\left[\begin{smallmatrix} 2 & 1\\\\ 1 & 2 \end{smallmatrix}\right]$

preserves the direction of purple vectors parallel to \textbf{v}\textasciitilde{}\textit{?}=1\textasciitilde{} = \[1

-1\]^T^ and blue vectors parallel to \textbf{v}\textasciitilde{}\textit{?}=3\textasciitilde{} = \[1 1\]^T^. The red

vectors are not parallel to either eigenvector, so, their directions are

changed by the transformation. The lengths of the purple vectors are

unchanged after the transformation (due to their eigenvalue of 1), while

blue vectors are three times the length of the original (due to their

eigenvalue of

3).](Eigenvectors.gif ""The transformation matrix A = \left[\begin{smallmatrix} 2 \& 1\\\\ 1 \& 2 \end{smallmatrix}\right] preserves the direction of purple vectors parallel to v?=1 = [1 -1]T and blue vectors parallel to v?=3 = [1 1]T. The red vectors are not parallel to either eigenvector, so, their directions are changed by the transformation. The lengths of the purple vectors are unchanged after the transformation (due to their eigenvalue of 1), while blue vectors are three times the length of the original (due to their eigenvalue of 3)."")



Consider the matrix



$$A = \begin{bmatrix}

  2 & 1\\\  1 & 2

\end{bmatrix}.$$



The figure on the right shows the effect of this transformation on point

coordinates in the plane. The eigenvectors \textit{v} of this transformation

satisfy Equation (\textbf{1}), and the values of \textit{?} for which the

determinant of the matrix (\textit{A} - \textit{?I}) equals zero are the eigenvalues.



Taking the determinant to find characteristic polynomial of \textit{A},



\begin{align}
     |A - \lambda I| 

  &= \left|\begin{bmatrix}

    2 & 1 \\\    1 & 2

  \end{bmatrix} - \lambda\begin{bmatrix}

    1 & 0 \\\    0 & 1

  \end{bmatrix}\right| = \begin{vmatrix}

    2 - \lambda & 1 \\\    1           & 2 - \lambda

  \end{vmatrix} \\\\[6pt]

     &= 3 - 4\lambda + \lambda^2 \\\\[6pt]

     &= (\lambda - 3)(\lambda - 1).
\end{align}



Setting the characteristic polynomial equal to zero, it has roots at

\textit{?}=1 and \textit{?}=3, which are the two eigenvalues of \textit{A}.



For \textit{?}=1, Equation (\textbf{2}) becomes,



$$(A - I)\mathbf{v}_{\lambda=1} = \begin{bmatrix} 1 & 1\\\\ 1 & 1\end{bmatrix}\begin{bmatrix}v_1 \\\\ v_2\end{bmatrix} = \begin{bmatrix}0 \\\\ 0\end{bmatrix}$$



$$1v_1 + 1v_2 = 0$$; $1v_1 + 1v_2 = 0$



Any nonzero vector with \textit{v}\textasciitilde{}1\textasciitilde{} = -\textit{v}\textasciitilde{}2\textasciitilde{} solves this equation.

Therefore,



$$\mathbf{v}_{\lambda=1} = \begin{bmatrix} v_1 \\\\ -v_1 \end{bmatrix} = \begin{bmatrix} 1 \\\\ -1 \end{bmatrix}$$



is an eigenvector of \textit{A} corresponding to \textit{?} = 1, as is any scalar

multiple of this vector.



For \textit{?}=3, Equation (\textbf{2}) becomes



\begin{align}
  (A - 3I)\mathbf{v}_{\lambda=3} &=

  \begin{bmatrix} -1 & 1\\\\ 1 & -1 \end{bmatrix}

    \begin{bmatrix} v_1 \\\\ v_2 \end{bmatrix} =

    \begin{bmatrix} 0 \\\\ 0 \end{bmatrix} \\\  -1v_1 + 1v_2 &= 0;\\\   1v_1 - 1v_2 &= 0
\end{align}



Any nonzero vector with \textit{v}\textasciitilde{}1\textasciitilde{} = \textit{v}\textasciitilde{}2\textasciitilde{} solves this equation. Therefore,



$$\mathbf v_{\lambda=3} = \begin{bmatrix} v_1 \\\\ v_1 \end{bmatrix} = \begin{bmatrix} 1 \\\\ 1 \end{bmatrix}$$



is an eigenvector of \textit{A} corresponding to \textit{?} = 3, as is any scalar

multiple of this vector.



Thus, the vectors \textbf{v}\textasciitilde{}\textit{?}=1\textasciitilde{} and \textbf{v}\textasciitilde{}\textit{?}=3\textasciitilde{} are eigenvectors of \textit{A}

associated with the eigenvalues ?=1 and ?=3, respectively.



\paragraph{General formula for eigenvalues of a two-dimensional matrix [general\_formula\_for\_eigenvalues\_of\_a\_two\_dimensional\_matrix]}



The eigenvalues of the real matrix $A = \begin{bmatrix}

  a \& b\\\  c \& d

\end{bmatrix}$ are



:   $\lambda = \frac{1}{2}(a + d) \pm \sqrt{\left(\frac{1}{2}(a - d)\right)^2 + bc}

    = \frac{1}{2}\left(\operatorname{tr}(A) \pm \sqrt{\left(\operatorname{tr}(A)^2\right) - 4\det(A)}\right) .$



Note that the eigenvalues are always real if \textit{b} and \textit{c} have the same

sign, since the quantity under the radical must be positive.



\paragraph{Three-dimensional matrix example [three\_dimensional\_matrix\_example]}



Consider the matrix



$$A = \begin{bmatrix}

  2 & 0 & 0 \\\  0 & 3 & 4 \\\  0 & 4 & 9

\end{bmatrix}.$$



The characteristic polynomial of \textit{A} is



\begin{align}
  |A-\lambda I| &= \left|\begin{bmatrix}

    2 & 0 & 0 \\\    0 & 3 & 4 \\\    0 & 4 & 9

  \end{bmatrix} - \lambda\begin{bmatrix}

    1 & 0 & 0 \\\    0 & 1 & 0 \\\    0 & 0 & 1

  \end{bmatrix}\right| =

  \begin{vmatrix}

    2 - \lambda & 0 & 0 \\\    0 & 3 - \lambda & 4 \\\    0 & 4 & 9 - \lambda

  \end{vmatrix}, \\\\[6pt]

                &= (2 - \lambda)\bigl[(3 - \lambda)(9 - \lambda) - 16\bigr]

                 = -\lambda^3 + 14\lambda^2 - 35\lambda + 22.
\end{align}



The roots of the characteristic polynomial are 2, 1, and 11, which are

the only three eigenvalues of \textit{A}. These eigenvalues correspond to the

eigenvectors $\begin{bmatrix} 1 & 0 & 0 \end{bmatrix}^\textsf{T}$,

$\begin{bmatrix} 0 & -2 & 1 \end{bmatrix}^\textsf{T}$, and

$\begin{bmatrix} 0 & 1 & 2 \end{bmatrix}^\textsf{T}$, or any nonzero

multiple thereof.



\paragraph{Three-dimensional matrix example with complex eigenvalues [three\_dimensional\_matrix\_example\_with\_complex\_eigenvalues]}



Consider the cyclic permutation matrix



$$A = \begin{bmatrix}

  0 & 1 & 0\\\  0 & 0 & 1\\\  1 & 0 & 0

\end{bmatrix}.$$



This matrix shifts the coordinates of the vector up by one position and

moves the first coordinate to the bottom. Its characteristic polynomial

is 1 - \textit{?}^3^, whose roots are



\begin{align}
  \lambda_1 &= 1 \\\  \lambda_2 &= -\frac{1}{2} + i \frac{\sqrt{3}}{2} \\\  \lambda_3 &= \lambda_2^* = -\frac{1}{2} - i \frac{\sqrt{3}}{2}
\end{align}



where $i$ is an imaginary unit with $i^2 = -1$.



For the real eigenvalue \textit{?}\textasciitilde{}1\textasciitilde{} = 1, any vector with three equal nonzero

entries is an eigenvector. For example,



$$A \begin{bmatrix} 5\\\\ 5\\\\ 5 \end{bmatrix} =

    \begin{bmatrix} 5\\\\ 5\\\\ 5 \end{bmatrix} =

    1 \cdot \begin{bmatrix} 5\\\\ 5\\\\ 5 \end{bmatrix}.$$



For the complex conjugate pair of imaginary eigenvalues,



$$\lambda_2\lambda_3 = 1, \quad \lambda_2^2 = \lambda_3, \quad \lambda_3^2 = \lambda_2.$$"

%%===============================================================
\section{The Geometric Interpretation of the Determinant}
\label{sec:the-geometric-interpretation-of-the-determinant}
%%===============================================================

"This parallelogram picture











Linalg\_parallelogram.png









is familiar from the construction of the sum of the two vectors. One way

to compute the area that it encloses is to draw this rectangle and

subtract the area of each subregion.







         $\begin{array}{l}

\text{area of parallelogram}                    \\\\quad 

=\text{area of rectangle}

-\text{area of }A-\text{area of }B \\\\qquad 

-\cdots-\text{area of }F                   \\\\quad 

=(x_1+x_2)(y_1+y_2)-x_2y_1-x_1y_1/2        \\\\qquad 

-x_2y_2/2-x_2y_2/2-x_1y_1/2-x_2y_1         \\\\quad 

=x_1y_2-x_2y_1        

\end{array}$







The fact that the area equals the value of the determinant



$$\begin{vmatrix}

x_1  &x_2  \\\y_1  &y_2

\end{vmatrix}

=x_1y_2-x_2y_1$$



is no coincidence.



The properties in the definition of determinants make reasonable

postulates for a function that measures the size of the region enclosed

by the vectors in the matrix.



For instance, this shows the effect of multiplying one of the

box-defining vectors by a scalar (the scalar used is $k=1.4$).







         









The region formed by $k\vec{v}$ and $\vec{w}$ is bigger, by a factor of

$k$, than the shaded region enclosed by $\vec{v}$ and $\vec{w}$. That

is,

$\text{size}\, (k\vec{v},\vec{w})=k\cdot\text{size}\, (\vec{v},\vec{w})$

and in general we expect of the size measure that

$\text{size}\, (\dots,k\vec{v},\dots)=k\cdot\text{size}\, (\dots,\vec{v},\dots)$.

Of course, this postulate is already familiar as one of the properties

in the defintion of determinants.



Another property of determinants is that they are unaffected by

pivoting. Here are before-pivoting and after-pivoting boxes (the scalar

used is $k=0.35$).







     









Although the region on the right, the box formed by $v$ and

$k\vec{v}+\vec{w}$, is more slanted than the shaded region, the two have

the same base and the same height and hence the same area. This

illustrates that

$\text{size}\, (\vec{v},k\vec{v}+\vec{w})=\text{size}\, (\vec{v},\vec{w})$.

Generalized, $\text{size}\, (\dots,\vec{v},\dots,\vec{w},\dots)

=\text{size}\, (\dots,\vec{v},\dots,k\vec{v}+\vec{w},\dots)$, which is a

restatement of the determinant postulate.



Of course, this picture











Linalg\_parallelogram\_basis.png









shows that $\text{size}\, (\vec{e}_1,\vec{e}_2)=1$, and we naturally

extend that to any number of dimensions

$\text{size}\,(\vec{e}_1,\dots,\vec{e}_n)=1$, which is a restatement of

the property that the determinant of the identity matrix is one.



With that, because property (2) of determinants is redundant (as

remarked right after the definition), we have that all of the properties

of determinants are reasonable to expect of a function that gives the

size of boxes. We can now cite the work done in the prior section to

show that the determinant exists and is unique to be assured that these

postulates are consistent and sufficient (we do not need any more

postulates). That is, we've got an intuitive justification to interpret

$\det(\vec{v}_1,\dots,\vec{v}_n)$ as the size of the box formed by the

vectors.



> \textbf{Example 1.1:} The volume of this parallelepiped, which can be found

> by the usual formula from high school geometry, is $12$.

>

> 

>

>  height=""300"" alt=""Linalg\_parallelepiped.png"" />         

> $\begin{vmatrix}

> 2 \&0 \&-1\\\> 0 \&3 \&0 \\\> 2 \&1 \&1

> \end{vmatrix}=12$

>

> 



\textbf{Remark 1.2:} Although property (2) of the definition of determinants

is redundant, it raises an important point. Consider these two.

















Linalg\_parallelogram\_orientation\_1.png













Linalg\_parallelogram\_orientation\_2.png











$\begin{vmatrix}

4  \&1   \\\2  \&3

\end{vmatrix}=10$







$\begin{vmatrix}

1  \&4   \\\3  \&2

\end{vmatrix}=-10$









The only difference between them is in the order in which the vectors

are taken. If we take $\vec{u}$ first and then go to $\vec{v}$, follow

the counterclockwise arc shown, then the sign is positive. Following a

clockwise arc gives a negative sign.



The sign returned by the size function reflects the ""orientation"" or

""sense"" of the box. (We see the same thing if we picture the effect of

scalar multiplication by a negative scalar.)



Although it is both interesting and important, the idea of orientation

turns out to be tricky. It is not needed for the development below, and

so we will pass it by.



> \textbf{Definition 1.3:} The \textbf{box} (or \textbf{parallelepiped}) formed by

> $\langle

> \vec{v}_1,\dots,\vec{v}_n \rangle$ (where each vector is from

> $\mathbb{R}^n$) includes all of the set

> $\{t_1\vec{v}_1+\dots+t_n\vec{v}_n \,\big|\, t_1,\ldots,t_n\in [0..1]\}$.

> The \textbf{volume} of a box is the absolute value of the determinant of

> the matrix with those vectors as columns.



> Example 1.4: Volume, because it is an absolute value, does not depend

> on the order in which the vectors are given. The volume of the

> parallelepiped in can also be computed as the absolute value of this

> determinant.

>

> $$\begin{vmatrix}

> 0  &2 &-1 \\\> 3  &0 &0 \\\> 1  &2 &1

> \end{vmatrix}=-12$$



The definition of volume gives a geometric interpretation to something

in the space, boxes made from vectors. The next result relates the

geometry to the functions that operate on spaces.



> \textbf{Theorem 1.5:} A transformation $t:\mathbb{R}^n\to \mathbb{R}^n$

> changes the size of all boxes by the same factor, namely the size of

> the image of a box $\left|t(S)\right|$ is $\left|T\right|$ times the

> size of the box $\left|S\right|$, where $T$ is the matrix representing

> $t$ with respect to the standard basis. That is, for all

> $n \! \times \! n$ matrices, the determinant of a product is the

> product of the determinants

> $\left|TS\right|=\left|T\right|\cdot\left|S\right|$.



The two sentences state the same idea, first in map terms and then in

matrix terms. Although we tend to prefer a map point of view, the second

sentence, the matrix version, is more convienent for the proof and is

also the way that we shall use this result later.



\textbf{Proof:} The two statements are equivalent because

$\left|t(S)\right|=\left|TS\right|$, as both give the size of the box

that is the image of the unit box $\mathcal{E}_n$ under the composition

$t\circ s$ (where $s$ is the map represented by $S$ with respect to the

standard basis).



First consider the case that $\left|T\right|=0$. A matrix has a zero

determinant if and only if it is not invertible. Observe that if $TS$ is

invertible, so that there is an $M$ such that $(TS)M=I$, then the

associative property of matrix multiplication $T(SM)=I$ shows that $T$

is also invertible (with inverse $SM$). Therefore, if $T$ is not

invertible then neither is $TS$ — if $\left|T\right|=0$ then

$\left|TS\right|=0$, and the result holds in this case.



Now consider the case that $\left|T\right|\neq 0$, that $T$ is

nonsingular. Recall that any nonsingular matrix can be factored into a

product of elementary matrices, so that $TS=E_1E_2\cdots E_rS$. In the

rest of this argument, we will verify that if $E$ is an elementary

matrix then $\left|ES\right|=\left|E\right|\cdot\left|S\right|$. The

result will follow because then

$\left|TS\right|=\left|E_1\cdots E_rS\right|=\left|E_1\right|\cdots\left|E_r\right|\cdot\left|S\right|

=\left|E_1\cdots E_r\right|\cdot\left|S\right|=\left|T\right|\cdot\left|S\right|$.



If the elementary matrix $E$ is $M_i(k)$ then $M_i(k)S$ equals $S$

except that row $i$ has been multiplied by $k$. The third property of

determinant functions then gives that

$\left|M_i(k)S\right|=k\cdot\left|S\right|$. But

$\left|M_i(k)\right|=k$, again by the third property because $M_i(k)$ is

derived from the identity by multiplication of row $i$ by $k$, and so

$\left|ES\right|=\left|E\right|\cdot\left|S\right|$ holds for

$E=M_i(k)$. The $E=P_{i,j}=-1$ and $E=C_{i,j}(k)$ checks are similar.



> \textbf{Example 1.6:} Application of the map $t$ represented with respect

> to the standard bases by

>

> $$\begin{pmatrix}

> 1  &1  \\\> -2  &0

> \end{pmatrix}$$

>

> will double sizes of boxes, e.g., from this

>

> 

>

>  title=""Linalg\_parallelogram\_doubled\_1.png"" height=""200""

> alt=""Linalg\_parallelogram\_doubled\_1.png"" />          $\begin{vmatrix}

> 2  \&1  \\\> 1  \&2

> \end{vmatrix}=3$

>

> 

>

> to this

>

> 

>

>  title=""Linalg\_parallelogram\_doubled\_2.png"" height=""300""

> alt=""Linalg\_parallelogram\_doubled\_2.png"" />          $\begin{vmatrix}

> 3  \&3  \\\> -4  \&-2

> \end{vmatrix}=6$

>

> 



> \textbf{Corollary 1.7:} If a matrix is invertible then the determinant of

> its inverse is the inverse of its determinant

> $\left|T^{-1}\right|=1/\left|T\right|$.



> Proof:

> $1=\left|I\right|=\left|TT^{-1}\right|=\left|T\right|\cdot\left|T^{-1}\right|$



Recall that determinants are not additive homomorphisms, $\det(A+B)$

need not equal $\det(A)+\det(B)$. The above theorem says, in contrast,

that determinants are multiplicative homomorphisms: $\det(AB)$ does

equal $\det(A)\cdot \det(B)$.



\paragraph{Licensing [licensing\_15]}



Content obtained and/or adapted from:



\begin{itemize}
    \item \href{https://en.wikibooks.org/wiki/Linear\_Algebra/Determinants\_as\_Size\_Functions}{Determinants as Size Functions, Wikibooks: Linear
\end{itemize}
    Algebra}

    under a CC BY-SA license"

%%===============================================================
\section{Symmetric Matrices}
\label{sec:symmetric-matrices}
%%===============================================================

Symmetric Matrices

%%===============================================================
\section{Quadratic Forms}
\label{sec:quadratic-forms}
%%===============================================================

Quadratic Forms

%%===============================================================
\section{Singular Values}
\label{sec:singular-values}
%%===============================================================

Singular Values

%%===============================================================
\section{test}
\label{sec:test}
%%===============================================================

class="lt-math-14497 editable">Algebraic Procedures



\paragraph{Outcomes}



1.  Use elementary operations to find the solution to a linear system of equations.

2.  Find the row-echelon form and reduced row-echelon form of a matrix.

3.  Determine whether a system of linear equations has no solution, a unique solution or an infinite number of solutions from its .

4.  Solve a system of equations using Gaussian Elimination and Gauss-Jordan Elimination.

5.  Model a physical system with linear equations and then solve.



We have taken an in depth look at graphical representations of systems of equations, as well as how to find possible solutions graphically. Our attention now turns to working with systems algebraically.



\paragraph{Definition \\(\\PageIndex{1}\\): System of Linear Equations}



A \textbf{system of linear equations} is a list of equations, $$\\begin{array}{c} a_{11}x_{1}+a_{12}x_{2}+\\cdots +a_{1n}x_{n}=b_{1} \\\\ a_{21}x_{1}+a_{22}x_{2}+\\cdots +a_{2n}x_{n}=b_{2} \\\\ \\vdots \\\\ a_{m1}x_{1}+a_{m2}x_{2}+\\cdots +a_{mn}x_{n}=b_{m} \\end{array}\\nonumber$$ where $a_{ij}$ and $b_{j}$ are real numbers. The above is a system of $m$ equations in the $n$ variables, $x_{1},x_{2}\\cdots ,x_{n}$. Written more simply in terms of summation notation, the above can be written in the form $$\\sum_{j=1}^{n}a_{ij}x_{j}=b_{i}, \\text{ }i=1,2,3,\\cdots ,m\\nonumber$$



The relative size of $m$ and $n$ is not important here. Notice that we have allowed $a_{ij}$ and $b_{j}$ to be any real number. We can also call these numbers \textbf{scalars} . We will use this term throughout the text, so keep in mind that the term \textbf{scalar} just means that we are working with real numbers.



Now, suppose we have a system where $b_{i} = 0$ for all $i$. In other words every equation equals $0$. This is a special type of system.



\paragraph{Definition \\(\\PageIndex{2}\\): Homogeneous System of Equations}



A system of equations is called \textbf{homogeneous} if each equation in the system is equal to $0$. A homogeneous system has the form $$\\begin{array}{c} a_{11}x_{1}+a_{12}x_{2}+\\cdots +a_{1n}x_{n}= 0 \\\\ a_{21}x_{1}+a_{22}x_{2}+\\cdots +a_{2n}x_{n}= 0 \\\\ \\vdots \\\\ a_{m1}x_{1}+a_{m2}x_{2}+\\cdots +a_{mn}x_{n}= 0 \\end{array}\\nonumber $$ where $a_{ij}$ are scalars and $x_{i}$ are variables.



Recall from the previous section that our goal when working with systems of linear equations was to find the point of intersection of the equations when graphed. In other words, we looked for the solutions to the system. We now wish to find these solutions algebraically. We want to find values for $x_{1},\\cdots ,x_{n}$ which solve all of the equations. If such a set of values exists, we call $\\left( x_{1},\\cdots ,x_{n}\\right)$ the \textbf{solution set}.



Recall the above discussions about the types of solutions possible. We will see that systems of linear equations will have one unique solution, infinitely many solutions, or no solution. Consider the following definition.



\paragraph{Definition \\(\\PageIndex{3}\\): Consistent and Inconsistent Systems}



A system of linear equations is called \textbf{consistent} if there exists at least one solution. It is called \textbf{inconsistent} if there is no solution.



If you think of each equation as a condition which must be satisfied by the variables, consistent would mean there is some choice of variables which can satisfy \textbf{all} the conditions. Inconsistent would mean there is no choice of the variables which can satisfy all of the conditions.



The following sections provide methods for determining if a system is consistent or inconsistent, and finding solutions if they exist.



\paragraph{Elementary Operations}



We begin this section with an example. Recall from \href{https://math.libretexts.org/Bookshelves/Linear_Algebra/A_First_Course_in_Linear_Algebra_(Kuttler}{Example 1.1.1}/01\%3A_Systems_of_Equations/1.01\%3A_Geometry#graphicalsoln "1.1: Geometry") that the solution to the given system was $\\left(x, y \\right) = \\left( -1, 4 \\right)$.



\paragraph{Example \\(\\PageIndex{1}\\): Verifying an Ordered Pair is a Solution}



Algebraically verify that $\\left(x, y \\right) = \\left( -1, 4 \\right)$ is a solution to the following system of equations.



$$\\begin{array}{c} x+y=3 \\\\ y-x=5 \\end{array}\\nonumber $$



\paragraph{Solution}



By graphing these two equations and identifying the point of intersection, we previously found that $\\left(x, y \\right) = \\left( -1, 4 \\right)$ is the unique solution.



We can verify algebraically by substituting these values into the original equations, and ensuring that the equations hold. First, we substitute the values into the first equation and check that it equals $3$. $$x + y = (-1)+(4) = 3\\nonumber $$ This equals $3$ as needed, so we see that $\\left( -1,4 \\right)$ is a solution to the first equation. Substituting the values into the second equation yields $$y -x = (4) - (-1) = 4 + 1 = 5\\nonumber $$ which is true. For $\\left( x,y\\right) =\\left( -1,4\\right)$ each equation is true and therefore, this is a solution to the system.



Now, the interesting question is this: If you were not given these numbers to verify, how could you algebraically determine the solution? Linear algebra gives us the tools needed to answer this question. The following basic operations are important tools that we will utilize.



\paragraph{Definition \\(\\PageIndex{4}\\): Elementary Operations}



\textbf{Elementary operations} are those operations consisting of the following.



1.  Interchange the order in which the equations are listed.

2.  Multiply any equation by a nonzero number.

3.  Replace any equation with itself added to a multiple of another equation.



It is important to note that none of these operations will change the set of solutions of the system of equations. In fact, elementary operations are the _key tool_ we use in linear algebra to find solutions to systems of equations.



Consider the following example.



\paragraph{Example \\(\\PageIndex{2}\\): Effects of an Elementary Operation}



Show that the system $$\\begin{array}{c} x+y=7 \\\\ 2x-y=8 \\end{array}\\nonumber $$ has the same solution as the system $$\\begin{array}{c} x+y=7 \\\\ -3y=-6 \\end{array}\\nonumber $$



\paragraph{Solution}



Notice that the second system has been obtained by taking the second equation of the first system and adding -2 times the first equation, as follows: $$2x-y + (-2)(x+y) = 8 + (-2)(7)\\nonumber $$ By simplifying, we obtain $$-3y=-6\\nonumber $$ which is the second equation in the second system. Now, from here we can solve for $y$ and see that $y=2$. Next, we substitute this value into the first equation as follows $$x+y=x+2=7\\nonumber $$ Hence $x=5$ and so $\\left( x,y\\right) = \\left(5,2 \\right)$ is a solution to the second system. We want to check if $\\left(5,2 \\right)$ is also a solution to the first system. We check this by substituting $\\left(x, y \\right) = \\left(5,2 \\right)$ into the system and ensuring the equations are true. $$\\begin{array}{c} x+y = \\left(5 \\right)+ \\left( 2 \\right) = 7 \\\\ 2x-y= 2 \\left(5 \\right) - \\left( 2 \\right) = 8 \\end{array}\\nonumber $$ Hence, $\\left(5,2 \\right)$ is also a solution to the first system.



This example illustrates how an elementary operation applied to a system of two equations in two variables does not affect the solution set. However, a linear system may involve many equations and many variables and there is no reason to limit our study to small systems. For any size of system in any number of variables, the solution set is still the collection of solutions to the equations. In every case, the above operations of \href{Definition \\(\\PageIndex{4}\\)} do not change the set of solutions to the system of linear equations.



In the following theorem, we use the notation $E_i$ to represent an equation, while $b_i$ denotes a constant.



\paragraph{Theorem \\(\\PageIndex{1}\\): Elementary Operations and Solutions}



Suppose you have a system of two linear equations $$\\begin{array}{c} E_{1}=b_{1}\\\\ E_{2}=b_{2} \\end{array} \\label{system}$$ Then the following systems have the same solution set as \\(\\eqref{system}\\):



1.  $$\\begin{array}{c} E_{2}=b_{2}\\\\ E_{1}=b_{1} \\end{array} \\label{thm1.9.1}$$

2.  $$\\begin{array}{c} E_{1}=b_{1} \\\\ kE_{2}=kb_{2}\\\\ \\end{array} \\label{thm1.9.2}$$ for any scalar $k$, provided $k\\neq0$.

3.  $$\\begin{array}{c} E_{1}=b_{1} \\\\ E_{2}+kE_{1}=b_{2}+kb_{1} \\end{array} \\label{thm1.9.3}$$ for any scalar $k$ (including $k=0$).



Before we proceed with the proof of \href{Theorem \\(\\PageIndex{1}\\)}, let us consider this theorem in context of \href{Example \\(\\PageIndex{2}\\)}. Then, $$\\begin{array}{cc} E_{1} = x+y, & b_{1} = 7 \\\\ E_{2} = 2x-y, & b_{2} = 8 \\end{array}\\nonumber $$ Recall the elementary operations that we used to modify the system in the solution to the example. First, we added $\\left( -2 \\right)$ times the first equation to the second equation. In terms of \href{Theorem \\(\\PageIndex{1}\\)}, this action is given by $$E_{2} + \\left( -2 \\right) E_{1} = b_{2} + \\left( -2 \\right)b_{1}\\nonumber $$ or $$2x-y + \\left( -2 \\right) \\left(x+y \\right) = 8 + \\left( -2 \\right) 7\\nonumber $$ This gave us the second system in \href{Example \\(\\PageIndex{2}\\)}, given by $$\\begin{array}{c} E_{1} = b_{1} \\\\ E_{2} + \\left( -2 \\right) E_{1} = b_{2} + \\left( -2 \\right) b_{1} \\end{array}\\nonumber $$



From this point, we were able to find the solution to the system. \href{Theorem \\(\\PageIndex{1}\\)} tells us that the solution we found is in fact a solution to the original system.



We will now prove \href{Theorem \\(\\PageIndex{1}\\)}.



\textbf{Proof}



1.  The proof that the systems \\(\\eqref{system}\\) and \\(\\eqref{thm1.9.1}\\) have the same solution set is as follows. Suppose that $\\left( x_{1},\\cdots ,x_{n}\\right)$ is a solution to $E_{1}=b_{1},E_{2}=b_{2}$. We want to show that this is a solution to the system in \\(\\eqref{thm1.9.1}\\) above. This is clear, because the system in \\(\\eqref{thm1.9.1}\\) is the original system, but listed in a different order. Changing the order does not effect the solution set, so $\\left( x_{1},\\cdots ,x_{n}\\right)$ is a solution to \\(\\eqref{thm1.9.1}\\).

2.  Next we want to prove that the systems \\(\\eqref{system}\\) and \\(\\eqref{thm1.9.2}\\) have the same solution set. That is $E_{1}=b_{1},E_{2}=b_{2}$ has the same solution set as the system $E_{1}=b_{1},kE_{2}=kb_{2}$ provided $k\\neq 0$. Let $\\left( x_{1},\\cdots ,x_{n}\\right)$ be a solution of $E_{1}=b_{1},E_{2}=b_{2},$. We want to show that it is a solution to $E_{1}=b_{1},kE_{2}=kb_{2}$. Notice that the only difference between these two systems is that the second involves multiplying the equation, $E_{2}=b_{2}$ by the scalar $k$. Recall that when you multiply both sides of an equation by the same number, the sides are still equal to each other. Hence if $\\left( x_{1},\\cdots ,x_{n}\\right)$ is a solution to $E_{2}=b_{2}$, then it will also be a solution to $kE_{2}=kb_{2}$. Hence, $\\left( x_{1},\\cdots ,x_{n}\\right)$ is also a solution to \\(\\eqref{thm1.9.2}\\).  

    Similarly, let $\\left( x_{1},\\cdots ,x_{n}\\right)$ be a solution of $E_{1}=b_{1},kE_{2}=kb_{2}$. Then we can multiply the equation $kE_{2}=kb_{2}$ by the scalar $1/k$, which is possible only because we have required that $k\\neq 0$. Just as above, this action preserves equality and we obtain the equation $E_{2} = b_{2}$. Hence $\\left( x_{1},\\cdots ,x_{n}\\right)$ is also a solution to $E_{1}=b_{1},E_{2}=b_{2}.$

3.  Finally, we will prove that the systems \\(\\eqref{system}\\) and \\(\\eqref{thm1.9.3}\\) have the same solution set. We will show that any solution of $E_{1}=b_{1},E_{2}=b_{2}$ is also a solution of \\(\\eqref{thm1.9.3}\\). Then, we will show that any solution of \\(\\eqref{thm1.9.3}\\) is also a solution of $E_{1}=b_{1},E_{2}=b_{2}$. Let $\\left( x_{1},\\cdots ,x_{n}\\right)$ be a solution to $E_{1}=b_{1},E_{2}=b_{2}$. Then in particular it solves $E_{1} = b_{1}$. Hence, it solves the first equation in \\(\\eqref{thm1.9.3}\\). Similarly, it also solves $E_{2} = b_{2}$. By our proof of \\(\\eqref{thm1.9.2}\\), it also solves $kE_{1}=kb_{1}$. Notice that if we add $E_{2}$ and $kE_{1}$, this is equal to $b_{2} + kb_{1}$. Therefore, if $\\left( x_{1},\\cdots ,x_{n}\\right)$ solves $E_{1}=b_{1},E_{2}=b_{2}$ it must also solve $E_{2}+kE_{1}=b_{2}+kb_{1}$.  

    Now suppose $\\left( x_{1},\\cdots ,x_{n}\\right)$ solves the system $E_{1}=b_{1}, E_{2}+kE_{1}=b_{2}+kb_{1}$. Then in particular it is a solution of $E_{1} = b_{1}$. Again by our proof of \\(\\eqref{thm1.9.2}\\), it is also a solution to $kE_{1}=kb_{1}$. Now if we subtract these equal quantities from both sides of $E_{2}+kE_{1}=b_{2}+kb_{1}$ we obtain $E_{2}=b_{2}$, which shows that the solution also satisfies $E_{1}=b_{1},E_{2}=b_{2}.$



Stated simply, the above theorem shows that the elementary operations do not change the solution set of a system of equations.



We will now look at an example of a system of three equations and three variables. Similarly to the previous examples, the goal is to find values for $x,y,z$ such that each of the given equations are satisfied when these values are substituted in.



\paragraph{Example \\(\\PageIndex{3}\\): Solving a System of Equations with Elementary Operations}



Find the solutions to the system,



$$\\begin{array}{c} x+3y+6z=25 \\\\ 2x+7y+14z=58 \\\\ 2y+5z=19 \\end{array} \\label{solvingasystem1}$$



\paragraph{Solution}



We can relate this system to \href{Theorem \\(\\PageIndex{1}\\)} above. In this case, we have $$\\begin{array}{c c} E_{1} = x + 3y + 6z, & b_{1} = 25\\\\ E_{2} = 2x+7y+14z, & b_{2} = 58 \\\\ E_{3} = 2y+5z, & b_{3} = 19 \\end{array}\\nonumber $$ \href{Theorem \\(\\PageIndex{1}\\)} claims that if we do elementary operations on this system, we will not change the solution set. Therefore, we can solve this system using the elementary operations given in \href{Definition \\(\\PageIndex{4}\\)}. First, replace the second equation by $\\left( -2\\right)$ times the first equation added to the second. This yields the system $$\\begin{array}{c} x+3y+6z=25 \\\\ y+2z=8 \\\\ 2y+5z=19 \\end{array} \\label{solvingasystem2}$$ Now, replace the third equation with $\\left( -2\\right)$ times the second added to the third. This yields the system $$\\begin{array}{c} x+3y+6z=25 \\\\ y+2z=8 \\\\ z=3 \\end{array} \\label{solvingasystem3}$$ At this point, we can easily find the solution. Simply take $z=3$ and substitute this back into the previous equation to solve for $y$, and similarly to solve for $x$. $$\\begin{array}{c} x + 3y + 6 \\left(3 \\right) = x + 3y + 18 = 25\\\\ y + 2 \\left(3 \\right) = y + 6 = 8 \\\\ z = 3 \\end{array}\\nonumber $$ The second equation is now $$y+6=8\\nonumber $$ You can see from this equation that $y = 2$. Therefore, we can substitute this value into the first equation as follows: $$x + 3 \\left(2 \\right) + 18 = 25\\nonumber $$ By simplifying this equation, we find that $x=1$. Hence, the solution to this system is $\\left( x,y,z \\right) = \\left( 1,2,3 \\right)$. This process is called \textbf{back substitution}.



Alternatively, in \\(\\eqref{solvingasystem3}\\) you could have continued as follows. Add $\\left( -2\\right)$ times the third equation to the second and then add $\\left( -6\\right)$ times the second to the first. This yields $$ \\begin{array}{c} x+3y=7 \\\\ y=2 \\\\ z=3 \\end{array}\\nonumber $$ Now add $\\left( -3\\right)$ times the second to the first. This yields $$ \\begin{array}{c} x=1 \\\\ y=2 \\\\ z=3 \\end{array}\\nonumber $$ a system which has the same solution set as the original system. This avoided back substitution and led to the same solution set. It is your decision which you prefer to use, as both methods lead to the correct solution, $\\left( x,y,z \\right) = \\left(1,2,3\\right)$.
